\documentclass[11pt,a4paper]{report}
% Compilation : tectonic Cours_Calcul_Quantique.tex (ou xelatex / pdflatex, deux passes)
% Figures à données : python3 tools/make_figures.py (régénère figures/*.tex depuis les expériences du projet)
\usepackage{iftex}
\ifPDFTeX
  \usepackage[utf8]{inputenc}
  \usepackage[T1]{fontenc}
  \usepackage{lmodern}
\else
  \usepackage{fontspec}
\fi
\usepackage[french]{babel}
\usepackage[a4paper,margin=2.3cm]{geometry}
\usepackage{amsmath,amssymb,mathtools}
\usepackage{xcolor}
\usepackage{tikz}
\usetikzlibrary{arrows.meta,calc,positioning,shapes.geometric,decorations.pathreplacing,babel}
\usepackage{pgfplots}
\pgfplotsset{compat=1.18}
\usepgfplotslibrary{groupplots}
\usepackage{tabularx,booktabs,array}
\usepackage{enumitem}
\setlist{nosep,leftmargin=*}
\usepackage{listings}
\usepackage[font=small,labelfont=bf]{caption}
\usepackage[raggedright]{titlesec}
\usepackage[most]{tcolorbox}
\usepackage[colorlinks=true,linkcolor=blue!50!black,urlcolor=blue!50!black,citecolor=blue!50!black]{hyperref}

\definecolor{cblue}{HTML}{2A78D6}
\definecolor{corange}{HTML}{EB6834}
\definecolor{cgreen}{HTML}{1BAF7A}
\definecolor{cink}{HTML}{1F2328}
\definecolor{cmuted}{HTML}{6B6F76}

% ---------------------------------------------------------------- notations
\newcommand{\ket}[1]{\lvert #1\rangle}
\newcommand{\bra}[1]{\langle #1\rvert}
\newcommand{\braket}[2]{\langle #1\vert #2\rangle}
\newcommand{\expval}[1]{\langle #1\rangle}
\newcommand{\C}{\mathbb{C}}
\newcommand{\R}{\mathbb{R}}
\usepackage{dsfont}
\newcommand{\Id}{\mathds{1}}
\newcommand{\tr}{\operatorname{tr}}

% ---------------------------------------------------------------- encadrés
\tcbset{coursbox/.style={enhanced,breakable,boxrule=0pt,leftrule=2.5pt,arc=0pt,
  left=6pt,right=6pt,top=4pt,bottom=4pt,fonttitle=\bfseries\small,coltitle=cink,
  attach title to upper={\ }}}
\newtcolorbox{definition}[1][]{coursbox,colback=cblue!5,colframe=cblue,title={Définition#1.}}
\newtcolorbox{resultat}[1][]{coursbox,colback=cgreen!6,colframe=cgreen,title={À retenir#1.}}
\newtcolorbox{projet}[1][]{coursbox,colback=corange!6,colframe=corange,title={Dans notre projet#1.}}
\newtcolorbox{piege}[1][]{coursbox,colback=cmuted!7,colframe=cmuted,title={Piège#1.}}

% ---------------------------------------------------------------- exercices
\newcounter{exo}[chapter]
\renewcommand{\theexo}{\thechapter.\arabic{exo}}
\newenvironment{exercice}[1][]{\refstepcounter{exo}\par\medskip\noindent
  \textbf{\color{cblue}Exercice~\theexo}\if\relax\detokenize{#1}\relax\else\ (#1)\fi.\ }{\par}
\newcommand{\exercices}{\section*{Exercices}\addcontentsline{toc}{section}{Exercices}
  \noindent\textit{Corrigés en annexe~\ref{annexe:corriges}. Chercher d'abord sans les regarder.}}

% ---------------------------------------------------------------- circuits
\tikzset{
  gate/.style={draw=cink,fill=white,minimum width=7mm,minimum height=6mm,inner sep=1.5pt,font=\small},
  wgate/.style={gate,minimum width=13mm},
  ctrl/.style={circle,fill=cink,inner sep=1.5pt},
  flow/.style={draw=cink,rounded corners=2pt,fill=white,align=center,font=\small,inner sep=4pt},
  arr/.style={-{Stealth[length=2.2mm]},cink,thick},
}
\newcommand{\targ}[2]{\draw[cink] (#1,#2) circle (2.1mm);
  \draw[cink] ([xshift=-2.1mm]#1,#2) -- ([xshift=2.1mm]#1,#2);
  \draw[cink] ([yshift=-2.1mm]#1,#2) -- ([yshift=2.1mm]#1,#2);}
\newcommand{\meter}[2]{\node[gate,minimum width=8mm] at (#1,#2) {};
  \draw[cink] ([xshift=-2.6mm,yshift=-1.2mm]#1,#2) arc (160:20:2.75mm);
  \draw[cink,-{Stealth[length=1.1mm]}] ([yshift=-1.4mm]#1,#2) -- ++(60:3mm);}

\lstset{basicstyle=\ttfamily\footnotesize,breaklines=true,columns=fullflexible,frame=single,rulecolor=\color{cmuted!50},
  backgroundcolor=\color{cmuted!6},keywordstyle=\color{cblue}\bfseries,commentstyle=\color{cmuted}\itshape,
  stringstyle=\color{corange},showstringspaces=false,language=Python,
  literate={é}{\'e}1{è}{\`e}1{à}{\`a}1{ê}{\^e}1{ç}{\c c}1{ô}{\^o}1{î}{\^i}1{ù}{\`u}1{É}{\'E}1{→}{$\to$}1{γ}{$\gamma$}1{β}{$\beta$}1}

\setcounter{secnumdepth}{2}
\setcounter{tocdepth}{1}
\setlength{\parindent}{0pt}
\setlength{\parskip}{0.45em}

% ============================================================================
% Itérations du projet HackathonQuantique : couleurs et marqueurs communs.
% Inclus par tous les cours et par la fiche projet, pour qu'une itération ait
% la même couleur partout. Pour ajouter une itération : une couleur, un nom,
% une date, une ligne dans \legendeiterations.
% ============================================================================
\definecolor{itV0}{HTML}{8D6E63}   % brun   : V0, modèle de la note d'intention
\definecolor{itV1}{HTML}{C2185B}   % framboise : V1, chaîne complète conforme à la note
\definecolor{itV2}{HTML}{00838F}   % cyan foncé : V2, problème reformulé
\definecolor{itV3}{HTML}{E09A00}   % ambre  : V3, concurrent équitable et problème plus grand
\definecolor{itV4}{HTML}{5E35B1}   % indigo : V4, réservé

\makeatletter
\newcommand{\itername}[1]{\@nameuse{itname@#1}}
\newcommand{\iterdate}[1]{\@nameuse{itdate@#1}}
\newcommand{\defiter}[3]{\@namedef{itname@#1}{#2}\@namedef{itdate@#1}{#3}}
\makeatother
\defiter{V0}{Modèle de la note d'intention}{24--30 oct. 2026}
\defiter{V1}{Chaîne complète conforme à la note}{4--5 nov. 2026}
\defiter{V2}{Problème reformulé}{5--6 nov. 2026}
\defiter{V3}{Concurrent équitable, problème plus grand}{à partir du 6 nov. 2026}

% Pastille : \iterbadge{V1}
\newcommand{\iterbadge}[1]{\tikz[baseline=(b.base)]\node[fill=it#1,text=white,rounded corners=2pt,
  inner xsep=3pt,inner ysep=1.5pt,font=\sffamily\bfseries\scriptsize] (b) {#1};}

% Bandeau de début de chapitre : \iterbanner{V0}
\newcommand{\iterbanner}[1]{\par\noindent\begin{tikzpicture}
  \node[fill=it#1!12,draw=it#1,line width=0.6pt,rounded corners=2pt,text width=\dimexpr\linewidth-12pt\relax,
    inner sep=4pt,anchor=west,font=\small] at (0,0)
    {\iterbadge{#1}\ \textbf{\color{it#1}Itération #1} --- \itername{#1} \hfill{\color{it#1}\iterdate{#1}}};
  \end{tikzpicture}\par\medskip}

% Encadré d'évolution : \begin{evolution}{V1}{titre} ... \end{evolution}
\newtcolorbox{evolution}[2]{enhanced,breakable,boxrule=0pt,leftrule=3pt,arc=0pt,
  left=6pt,right=6pt,top=4pt,bottom=4pt,colback=it#1!7,colframe=it#1,
  fonttitle=\bfseries\small,coltitle=it#1,attach title to upper={\ },
  title={\iterbadge{#1}\ Évolution #1 : #2.}}

% Légende des itérations, pour les avant-propos
\newcommand{\legendeiterations}{%
\begin{center}\small
\begin{tabular}{lll}
\iterbadge{V0} & \itername{V0} & \iterdate{V0}\\
\iterbadge{V1} & \itername{V1} & \iterdate{V1}\\
\iterbadge{V2} & \itername{V2} & \iterdate{V2}\\
\iterbadge{V3} & \itername{V3} & \iterdate{V3}\\
\end{tabular}
\end{center}}

% Encadré « pistes d'amélioration » : défini ici s'il manque (cours V0 sans cet encadré)
\providecolor{cviolet}{HTML}{7A4FB3}
\ifcsname pistes\endcsname\else
\newtcolorbox{pistes}[1][]{enhanced,breakable,boxrule=0pt,leftrule=2.5pt,arc=0pt,
  left=6pt,right=6pt,top=4pt,bottom=4pt,fonttitle=\bfseries\small,coltitle=cink,attach title to upper={\ },
  colback=cviolet!6,colframe=cviolet,title={Pistes d'amélioration#1.}}
\fi


\title{\textbf{Calcul quantique pour l'optimisation}\\[0.4em]\large Du qubit au QAOA : tout ce qu'il faut pour la formulation quantique\\ du compromis traînée / SER (hackathon quantique)}
\author{Notes de travail, projet \emph{HackathonQuantique}}
\date{Version du 30 octobre 2026}

\begin{document}
\pagenumbering{Alph}
\maketitle
\pagenumbering{roman}

\chapter*{Avant-propos}
\addcontentsline{toc}{chapter}{Avant-propos}
\paragraph{Pour qui, pour quoi.} Ce cours s'adresse à un ingénieur à l'aise avec l'algèbre linéaire, l'analyse et la mécanique des fluides, mais qui n'a \emph{jamais} fait de physique quantique. Il part des postulats et s'arrête exactement là où commence notre implémentation : la formulation \textbf{QUBO} d'un problème d'optimisation de forme, sa traduction en \textbf{hamiltonien d'Ising}, et sa résolution par \textbf{recuit} (simulé ou quantique) et par \textbf{QAOA}. Il ne couvre ni l'information quantique générale (téléportation, codes correcteurs), ni les algorithmes de Shor ou de Grover, sauf pour situer ce que le quantique promet et ne promet pas.

\paragraph{Fil conducteur.} Pourquoi le quantique, et la notion d'interférence (ch.~\ref{ch:intro}) $\to$ les bases sans prérequis : nombres complexes, expériences fondatrices, postulats (ch.~\ref{ch:bases}) $\to$ un qubit (ch.~\ref{ch:qubit}) $\to$ observables et évolution, c'est-à-dire de l'algèbre linéaire sur $\C^n$ (ch.~\ref{ch:algebre}) $\to$ $n$ qubits et produit tensoriel (ch.~\ref{ch:multi}) $\to$ portes et circuits (ch.~\ref{ch:circuits}) $\to$ hamiltoniens : un problème d'optimisation devient « trouver l'état d'énergie minimale » (ch.~\ref{ch:hamiltoniens}) $\to$ QUBO et Ising (ch.~\ref{ch:qubo}) $\to$ recuit simulé et quantique (ch.~\ref{ch:recuit}) $\to$ QAOA (ch.~\ref{ch:qaoa}) $\to$ simulateurs et outils (ch.~\ref{ch:outils}) $\to$ la chaîne complète du hackathon (ch.~\ref{ch:projet}).

\paragraph{Si vous n'avez jamais fait de quantique.} Lire les chapitres~\ref{ch:intro} à \ref{ch:multi} dans l'ordre, sans sauter : chaque notion y est d'abord motivée par une expérience ou un exemple chiffré, puis formalisée. Les chapitres~\ref{ch:circuits} à \ref{ch:qaoa} n'utilisent rien d'autre.

\paragraph{Lien avec le projet.} Les encadrés orange \emph{Dans notre projet} relient chaque notion à la chaîne du hackathon : profil NACA $(m,p,t)$ codé sur 9 bits, modèle quadratique de $C_D$ ajusté sur le modèle aérodynamique de la phase~1, QUBO à 9 qubits, QAOA et recuit. \textbf{Tous les nombres et toutes les courbes de résultats de ce cours sont calculés}, pas illustratifs : ils sortent de \texttt{Code/quantique/} (simulation exacte par vecteur d'état en NumPy), script \texttt{experiences.py}. Le chapitre~10 explique ce code.

\paragraph{Comment utiliser ce document.} Chaque chapitre se termine par un ou deux exercices courts, corrigés en annexe~\ref{annexe:corriges}. L'annexe~\ref{annexe:memo} est un aide-mémoire d'une page à garder sous la main pendant le hackathon.

\paragraph{Conventions.} Unités où $\hbar=1$. Les bits sont notés $b_k\in\{0,1\}$ et les spins $s_k\in\{-1,+1\}$, avec $b_k=(1-s_k)/2$. Pour un registre de $n$ qubits, on écrit $\ket{b_0b_1\cdots b_{n-1}}$, le qubit~0 \emph{à gauche}, et on indexe le vecteur d'état par l'entier $x=\sum_k b_k2^k$. Attention : Qiskit affiche les chaînes de bits dans l'ordre inverse (§\ref{sec:pieges-outils}).

\paragraph{Itérations du projet.} Le projet avance par itérations, chacune avec sa couleur dans tous les cours et dans la fiche projet (dossier \texttt{Fiche\_Projet/} du projet). Chaque chapitre commence par un bandeau qui indique l'itération de son contenu ; un ajout ou une révision faite lors d'une itération ultérieure apparaît dans un encadré « Évolution » de la couleur correspondante.
\legendeiterations

\tableofcontents
\cleardoublepage
\pagenumbering{arabic}

\chapter{Pourquoi le calcul quantique, et ce qu'il fait vraiment}\label{ch:intro}
\iterbanner{V0}

Avant toute formule, situons la place du quantique dans le hackathon. Nous voulons des profils qui réalisent un bon compromis entre traînée $C_D$ et surface équivalente radar (SER). La partie quantique ne calcule \emph{ni} l'aérodynamique \emph{ni} l'électromagnétisme : elle sert à \textbf{chercher les meilleures combinaisons de bits} dans un problème d'optimisation dont la fonction de coût est un \emph{modèle} quadratique ajusté sur quelques simulations classiques.

\begin{figure}[htbp]
\centering
\begin{tikzpicture}[node distance=6mm and 7mm]
\node[flow] (a) {Profils simulés\\ (phase 1 : $C_D$ ;\\ phase 2 : SER)};
\node[flow,right=of a] (b) {Modèle\\ quadratique\\ (moindres carrés)};
\node[flow,right=of b,fill=corange!10] (c) {QUBO\\ $\min_{\mathbf b}\ \mathbf b^\top Q\mathbf b$\\ (ch.~7)};
\node[flow,right=of c,fill=cblue!10] (d) {QAOA ou recuit\\ (ch.~8, 9)\\ \emph{partie quantique}};
\node[flow,below=9mm of c] (e) {Meilleurs candidats\\ $\to$ simulés à leur tour};
\draw[arr] (a) -- (b); \draw[arr] (b) -- (c); \draw[arr] (c) -- (d);
\draw[arr] (d) |- (e); \draw[arr] (e) -| (a);
\end{tikzpicture}
\caption{La boucle de la formulation quantique. Seul le bloc bleu fait appel au calcul quantique ; tout le reste est classique.}\label{fig:boucle-intro}
\end{figure}

\section{Bits, bits aléatoires, qubits}
\paragraph{Bit classique.} Un registre de $n$ bits est dans \emph{un} état parmi $2^n$. Le profil NACA $(m,p,t)$ codé sur 9 bits est l'un des $2^9=512$ profils de notre grille.

\paragraph{Bit aléatoire.} Si l'on ne connaît l'état qu'en probabilité, on le décrit par un vecteur $\mathbf p=(p_0,p_1)$ avec $p_i\ge0$ et $p_0+p_1=1$. Une opération aléatoire, par exemple « lancer une pièce » qui remplace le bit par un bit uniforme, est une matrice \emph{stochastique} (coefficients positifs, colonnes de somme 1) :
\[
\text{Pièce}=\begin{pmatrix}1/2&1/2\\1/2&1/2\end{pmatrix},\qquad \text{Pièce}\begin{pmatrix}1\\0\end{pmatrix}=\begin{pmatrix}1/2\\1/2\end{pmatrix}.
\]
Pour $n$ bits aléatoires, il faut $2^n$ probabilités. Le nombre de paramètres croît déjà exponentiellement. Ce n'est donc \emph{pas} cette croissance qui fait la spécificité du quantique.

\paragraph{Qubit.} Un qubit est décrit par deux \textbf{amplitudes} complexes $(\alpha,\beta)$, avec $|\alpha|^2+|\beta|^2=1$ ; les \emph{probabilités} de lire 0 ou 1 sont $|\alpha|^2$ et $|\beta|^2$. Les opérations sont des matrices \textbf{unitaires} (elles conservent la norme euclidienne), par exemple la porte de Hadamard
\[
H=\frac{1}{\sqrt2}\begin{pmatrix}1&1\\1&-1\end{pmatrix},\qquad H\begin{pmatrix}1\\0\end{pmatrix}=\begin{pmatrix}1/\sqrt2\\1/\sqrt2\end{pmatrix}.
\]
Après une seule porte, la mesure donne 0 ou 1 avec probabilité $1/2$, exactement comme la pièce. La différence n'apparaît qu'au second pas.

\section{L'interférence : la vraie différence}
Appliquons deux fois chaque opération à partir de 0.
\begin{itemize}
\item Deux pièces : $\text{Pièce}^2\,(1,0)^\top=(1/2,1/2)^\top$. Le résultat reste aléatoire.
\item Deux Hadamard : $H^2=\Id$, donc $H^2\,(1,0)^\top=(1,0)^\top$. On lit 0 \emph{avec certitude}.
\end{itemize}
Le mécanisme est visible sur l'arbre des amplitudes (figure~\ref{fig:interference}). Pour aboutir en 1, il y a deux chemins, d'amplitudes $+\tfrac12$ et $-\tfrac12$ : ils \textbf{s'annulent}. Pour aboutir en 0, les deux chemins valent $+\tfrac12$ : ils \textbf{s'ajoutent}. Des probabilités, positives, ne peuvent jamais se compenser ; des amplitudes, de signe (ou de phase) quelconque, le peuvent.

\begin{figure}[htbp]
\centering
\begin{tikzpicture}[x=2.6cm,y=1.1cm,lab/.style={font=\small,fill=white,inner sep=1pt}]
\node[flow] (s) at (0,0) {$\ket0$};
\node[flow] (a0) at (1,1) {$\ket0$}; \node[flow] (a1) at (1,-1) {$\ket1$};
\node[flow] (b00) at (2,1.6) {$\ket0$}; \node[flow] (b01) at (2,0.55) {$\ket1$};
\node[flow] (b10) at (2,-0.55) {$\ket0$}; \node[flow] (b11) at (2,-1.6) {$\ket1$};
\draw[arr] (s) -- node[lab,above left]{$\frac{1}{\sqrt2}$} (a0);
\draw[arr] (s) -- node[lab,below left]{$\frac{1}{\sqrt2}$} (a1);
\draw[arr] (a0) -- node[lab,above]{$\frac{1}{\sqrt2}$} (b00);
\draw[arr] (a0) -- node[lab,below]{$\frac{1}{\sqrt2}$} (b01);
\draw[arr] (a1) -- node[lab,above]{$\frac{1}{\sqrt2}$} (b10);
\draw[arr,corange] (a1) -- node[lab,below]{$-\frac{1}{\sqrt2}$} (b11);
\node[anchor=west,font=\small] at (2.25,1.6) {$+\tfrac12$};
\node[anchor=west,font=\small] at (2.25,0.55) {$+\tfrac12$};
\node[anchor=west,font=\small] at (2.25,-0.55) {$+\tfrac12$};
\node[anchor=west,font=\small,corange] at (2.25,-1.6) {$-\tfrac12$};
\draw[decorate,decoration={brace,amplitude=4pt}] (2.75,1.75) -- (2.75,-0.7) node[midway,right=5pt,font=\small,align=left] {$\ket0$ : $\tfrac12+\tfrac12=1$};
\node[font=\small,align=left,anchor=west] at (2.85,-1.1) {$\ket1$ : $\tfrac12-\tfrac12=0$};
\node[font=\small] at (0.5,-2.3) {porte $H$}; \node[font=\small] at (1.5,-2.3) {porte $H$};
\end{tikzpicture}
\caption{Arbre des amplitudes pour $H$ puis $H$. On \emph{additionne} les amplitudes des chemins qui mènent au même état, puis on prend le module au carré. Les deux chemins vers $\ket1$ interfèrent destructivement.}\label{fig:interference}
\end{figure}

\begin{resultat}
Un algorithme quantique est une façon d'orchestrer les interférences : faire s'annuler les amplitudes des mauvaises réponses et s'ajouter celles des bonnes. QAOA (ch.~\ref{ch:qaoa}) le fait en codant la fonction de coût dans des \emph{phases} $e^{-i\gamma E(x)}$, puis en mélangeant les états pour que ces phases interfèrent.
\end{resultat}

\section{Ce que le quantique promet, et ce qu'il ne promet pas}
\begin{itemize}
\item \textbf{Parallélisme ?} Un registre de $n$ qubits porte $2^n$ amplitudes, mais une mesure ne renvoie que \emph{$n$ bits}, tirés au hasard selon les $|{\rm amplitude}|^2$. On ne « lit » pas les 512 profils en même temps. Tout l'enjeu est de rendre la bonne réponse \emph{probable} avant de mesurer.
\item \textbf{Accélérations prouvées} : elles concernent des problèmes à structure particulière. Shor (factorisation) est exponentiel. Grover (recherche non structurée) est quadratique : environ $\sqrt N$ requêtes au lieu de $N$.
\item \textbf{Optimisation combinatoire} (notre cas) : QAOA et le recuit quantique sont des \emph{heuristiques}. Il n'existe pas de preuve d'avantage exponentiel pour les problèmes NP-difficiles, et l'on ne s'attend pas à en trouver.
\item \textbf{Matériel actuel} : il est bruité et de taille intermédiaire (ère \emph{NISQ}). La profondeur des circuits exécutables est limitée par les erreurs (ch.~\ref{ch:outils}).
\end{itemize}

\begin{projet}
La note d'intention le dit explicitement : \emph{aucun gain de temps de calcul n'est attendu}. Le cas (9 qubits) est volontairement assez petit pour être résolu exactement, et le quantique y tourne \emph{sur simulateur}. Le livrable, c'est une \textbf{formulation} correcte, une \textbf{qualité de solution} mesurée (nombre de simulations aérodynamiques pour retrouver le front de Pareto exact) et une analyse des \textbf{ressources} (qubits, portes, profondeur) en fonction de la taille du problème.
\end{projet}

\exercices
\begin{exercice}[interférence]
On part de $\ket1=(0,1)^\top$. Calculer l'état après une puis deux portes $H$, et les probabilités de mesure correspondantes. Comparer avec deux « pièces » appliquées au bit 1.
\end{exercice}
\begin{exercice}[parallélisme]
Un collègue affirme : « avec 9 qubits on évalue les 512 profils en parallèle, donc QAOA donne l'optimum en un seul passage ». Expliquer en trois phrases pourquoi c'est faux, et ce que l'on obtient réellement à la sortie d'un passage.
\end{exercice}

\chapter{Les bases : outils mathématiques et physique quantique sans prérequis}\label{ch:bases}
\iterbanner{V0}

Ce chapitre ne suppose \emph{aucune} connaissance de physique quantique. Il rassemble les outils mathématiques (nombres complexes, vecteurs et matrices complexes), puis les expériences qui obligent à abandonner la description classique. Il se termine par la liste des quatre postulats que les chapitres suivants détaillent un par un.

\section{Boîte à outils : les nombres complexes}
La mécanique quantique utilise des nombres complexes partout. Voici tout ce qui est nécessaire.
\begin{itemize}
\item $z=a+ib$ avec $i^2=-1$ ; \textbf{conjugué} $\bar z=a-ib$ ; \textbf{module} $|z|=\sqrt{a^2+b^2}$, et $|z|^2=z\bar z$.
\item \textbf{Forme polaire} : $z=r\,e^{i\varphi}$ avec $r=|z|$ et $\varphi$ l'\emph{argument} ou \emph{phase}, où $e^{i\varphi}=\cos\varphi+i\sin\varphi$ (formule d'Euler).
\item Multiplier par $e^{i\varphi}$, c'est \textbf{tourner} d'un angle $\varphi$ dans le plan complexe sans changer le module : $|e^{i\varphi}z|=|z|$.
\item $e^{i\alpha}e^{i\beta}=e^{i(\alpha+\beta)}$ : les phases s'\emph{additionnent} quand on multiplie.
\item Deux nombres de même module mais de phases opposées peuvent s'annuler : $e^{i0}+e^{i\pi}=1-1=0$. C'est l'interférence destructive du chapitre~\ref{ch:intro}.
\end{itemize}
\begin{figure}[htbp]
\centering
\begin{tikzpicture}[scale=1.6,>={Stealth[length=2mm]}]
\draw[->,cmuted] (-1.3,0) -- (1.4,0) node[right,cink,font=\small] {Re};
\draw[->,cmuted] (0,-1.2) -- (0,1.3) node[above,cink,font=\small] {Im};
\draw[cmuted,dashed] (0,0) circle (1);
\draw[->,cblue,very thick] (0,0) -- (40:1) node[anchor=south west,font=\small] {$e^{i\varphi}$};
\draw[corange,thick] (0.3,0) arc (0:40:0.3); \node[corange,font=\small] at (20:0.45) {$\varphi$};
\draw[->,corange,very thick] (0,0) -- (220:1) node[anchor=north east,font=\small] {$e^{i(\varphi+\pi)}=-e^{i\varphi}$};
\fill (1,0) circle (1pt) node[below right,font=\small] {$1$};
\fill (0,1) circle (1pt) node[above left,font=\small] {$i$};
\end{tikzpicture}
\caption{Le cercle unité. Une \emph{phase} est un point de ce cercle ; deux phases opposées se compensent quand on les additionne.}
\end{figure}

\section{Boîte à outils : vecteurs et matrices complexes}
\begin{itemize}
\item Un \textbf{vecteur} de $\C^N$ est une colonne de $N$ nombres complexes. Le \textbf{produit hermitien} de $u$ et $v$ vaut $\sum_k\bar u_kv_k$ (on conjugue le premier) ; la norme est $\|v\|^2=\sum_k|v_k|^2$.
\item La \textbf{transposée conjuguée} d'une matrice, notée $A^\dagger$ (« A dague »), est la transposée dont on conjugue chaque coefficient.
\item $A$ est \textbf{hermitienne} si $A^\dagger=A$. C'est la version complexe d'une matrice symétrique : ses valeurs propres sont \emph{réelles} et ses vecteurs propres forment une base orthonormée (théorème spectral).
\item $U$ est \textbf{unitaire} si $U^\dagger U=\Id$. C'est la version complexe d'une matrice de rotation : elle conserve les normes et les produits hermitiens.
\item \textbf{Valeur propre} : $Av=\lambda v$ avec $v\neq0$. Pour une matrice diagonale, les valeurs propres sont les coefficients diagonaux et les vecteurs propres les vecteurs de la base canonique.
\end{itemize}
Retenir l'analogie : \emph{hermitienne} $\leftrightarrow$ quelque chose qu'on mesure (valeurs réelles) ; \emph{unitaire} $\leftrightarrow$ quelque chose qu'on fait subir à l'état (une rotation).

\section{Une expérience qui casse l'intuition classique : la polarisation}
La lumière a une \textbf{polarisation} : la direction dans laquelle oscille le champ électrique. Un \textbf{polariseur} orienté à l'angle $\theta$ ne laisse passer que la composante du champ selon $\theta$.

\paragraph{En lumière intense (classique).} Une lumière polarisée à $0^\circ$ traversant un polariseur à $\theta$ voit son intensité multipliée par $\cos^2\theta$ (loi de Malus) : rien n'est surprenant, on projette un vecteur.

\paragraph{Photon par photon (quantique).} En baissant l'intensité jusqu'à n'envoyer qu'\emph{un photon à la fois}, on observe que chaque photon passe \emph{entièrement} ou est \emph{entièrement} absorbé : il n'y a jamais de « fraction de photon ». La loi de Malus devient une loi de \textbf{probabilité} : le photon passe avec probabilité $\cos^2\theta$. Et s'il passe, il ressort polarisé selon $\theta$, et non plus selon $0^\circ$ : la mesure a \emph{changé} l'état.

\paragraph{Le paradoxe des trois polariseurs.} Deux polariseurs croisés à $0^\circ$ puis $90^\circ$ ne laissent rien passer ($\cos^290^\circ=0$). Insérons un troisième polariseur à $45^\circ$ \emph{entre les deux} : de la lumière ressort ! La fraction transmise vaut $\cos^245^\circ\times\cos^245^\circ=\tfrac12\cdot\tfrac12=\tfrac14$ de ce qui sort du premier. Ajouter un filtre \emph{augmente} la transmission. Classiquement, c'est absurde ; quantiquement, c'est naturel : le polariseur à $45^\circ$ \emph{mesure} le photon et le \emph{prépare} dans un nouvel état, qui a une composante non nulle selon $90^\circ$.

\begin{figure}[htbp]
\centering
\begin{tikzpicture}[x=1cm,y=1cm]
\tikzset{pol/.style={draw=cink,fill=cmuted!10,minimum width=4mm,minimum height=16mm}}
\foreach \y/\lab/\mid in {1.4/{sans filtre central}/0,-1.0/{avec filtre à $45^\circ$}/1}{
  \node[font=\small,anchor=east] at (-0.3,\y) {\lab};
  \draw[corange,very thick,->] (0,\y) -- (1.1,\y);
  \node[pol] at (1.4,\y) {}; \draw[cink,thick] (1.4,\y-0.6) -- (1.4,\y+0.6);
  \node[font=\scriptsize,below] at (1.4,\y-0.85) {$0^\circ$};
}
% sans filtre
\draw[corange,very thick,->] (1.7,1.4) -- (4.6,1.4);
\node[pol] at (4.9,1.4) {}; \draw[cink,thick] (4.6,1.4) -- (5.2,1.4);
\node[font=\scriptsize,below] at (4.9,0.55) {$90^\circ$};
\node[font=\small,anchor=west] at (5.5,1.4) {rien ne passe};
% avec filtre
\draw[corange,very thick,->] (1.7,-1.0) -- (2.85,-1.0);
\node[pol] at (3.15,-1.0) {}; \draw[cink,thick] (2.8,-1.35) -- (3.5,-0.65);
\node[font=\scriptsize,below] at (3.15,-1.85) {$45^\circ$};
\draw[corange,thick,->] (3.45,-1.0) -- (4.6,-1.0);
\node[pol] at (4.9,-1.0) {}; \draw[cink,thick] (4.6,-1.0) -- (5.2,-1.0);
\node[font=\scriptsize,below] at (4.9,-1.85) {$90^\circ$};
\draw[corange,thin,->] (5.2,-1.0) -- (5.9,-1.0);
\node[font=\small,anchor=west] at (6.0,-1.0) {$\tfrac14$ passe};
\end{tikzpicture}
\caption{Le paradoxe des trois polariseurs. Le filtre central ne « bloque » pas : il mesure et re-prépare l'état du photon.}\label{fig:polariseurs}
\end{figure}

\paragraph{La description qui marche.} On représente la polarisation d'un photon par un \textbf{vecteur unitaire} $\ket\psi=\cos\phi\,\ket{0^\circ}+\sin\phi\,\ket{90^\circ}$. Un polariseur à $\theta$ :
\begin{enumerate}
\item laisse passer avec probabilité $|\langle\theta|\psi\rangle|^2=\cos^2(\theta-\phi)$, le carré de la projection ;
\item laisse le photon transmis dans l'état $\ket\theta$.
\end{enumerate}
\emph{Ce sont exactement les règles du qubit} : la polarisation d'un photon \emph{est} un qubit, avec $\ket0=\ket{0^\circ}$ et $\ket1=\ket{90^\circ}$. L'état $\ket{45^\circ}=\tfrac1{\sqrt2}(\ket0+\ket1)$ sera noté $\ket+$ au chapitre suivant.

\section{Ce que veut dire « superposition »}
Un photon polarisé à $45^\circ$ est dans la \textbf{superposition} $\tfrac1{\sqrt2}(\ket{0^\circ}+\ket{90^\circ})$. Trois contresens fréquents :
\begin{itemize}
\item \emph{« Il est à $0^\circ$ et à $90^\circ$ en même temps. »} Non : il est dans un état précis, $45^\circ$, qui passe à coup sûr un polariseur à $45^\circ$. « Superposé » est relatif à une base : $\ket{45^\circ}$ est superposé dans la base $\{0^\circ,90^\circ\}$, mais c'est un état de base dans la base $\{45^\circ,135^\circ\}$.
\item \emph{« Il est en réalité à $0^\circ$ ou à $90^\circ$, on ne sait juste pas lequel. »} Non : un tel mélange statistique, 50~\% de photons à $0^\circ$ et 50~\% à $90^\circ$, ne passerait le polariseur à $45^\circ$ qu'une fois sur deux. L'état superposé passe toujours. Les \emph{amplitudes} ne sont pas de l'ignorance.
\item \emph{« La mesure révèle une valeur préexistante. »} Non : elle \emph{produit} un résultat aléatoire et \emph{modifie} l'état (troisième polariseur).
\end{itemize}

\section{Le spin et l'expérience de Stern et Gerlach}
Un autre système à deux états, historiquement fondateur, est le \textbf{spin} d'un électron ou d'un atome d'argent : un petit moment magnétique. Envoyé dans un champ magnétique inhomogène orienté selon $z$ (appareil de Stern-Gerlach), un faisceau d'atomes se sépare en \emph{deux} taches seulement : « spin haut » ($+1$) et « spin bas » ($-1$). Classiquement, on attendrait un continuum. Les mêmes phénomènes réapparaissent :
\begin{itemize}
\item on trie les atomes « haut selon $z$ », puis on mesure selon $x$ : on obtient 50/50 ;
\item on remesure selon $z$ : on obtient de nouveau 50/50. La mesure selon $x$ a effacé l'information sur $z$.
\end{itemize}
C'est l'origine du vocabulaire « spin » ($s=\pm1$) utilisé pour le modèle d'Ising au chapitre~\ref{ch:qubo}. Les valeurs $\pm1$ y sont celles de l'observable $Z$.

\section{Énergie et niveaux discrets}
Le mot « quantique » vient de ce que l'énergie d'un système lié (atome, circuit supraconducteur) ne prend que des valeurs \emph{discrètes} : les \textbf{niveaux d'énergie}. Le plus bas est l'\textbf{état fondamental}. Un système en contact avec un environnement froid tend à y descendre.
\begin{resultat}
Toute l'idée de l'optimisation quantique tient en une phrase : \emph{construire un système dont les niveaux d'énergie sont les valeurs de notre fonction de coût}, puis l'amener dans son état fondamental. Le lire, c'est lire la solution. Les chapitres~\ref{ch:hamiltoniens} à \ref{ch:qaoa} réalisent ce programme.
\end{resultat}

\section{Qu'est-ce qu'un qubit physique ?}
N'importe quel système à deux niveaux bien isolé et contrôlable :
\begin{center}
\begin{tabular}{p{4.2cm}p{5.2cm}p{4.6cm}}
\toprule
Technologie & $\ket0$ / $\ket1$ & Exemples\\
\midrule
Circuits supraconducteurs & deux niveaux d'énergie d'un circuit LC non linéaire (transmon), à $\sim10$~mK & IBM, Google ; \textbf{D-Wave} (qubits de flux, recuit)\\
Ions piégés & deux niveaux d'un ion & Quantinuum, IonQ\\
Atomes neutres & niveaux d'atomes tenus par lasers & Pasqal (France)\\
Photons & polarisation ou chemin & Quandela (France)\\
Spins dans le silicium & spin d'un électron & Quobly (France)\\
\bottomrule
\end{tabular}
\end{center}
Pour ce cours, la technologie n'importe pas : on manipule des qubits abstraits, simulés sur ordinateur classique.

\section{Les quatre postulats, en une page}
Tout le calcul quantique repose sur quatre règles. Les chapitres suivants les déroulent.
\begin{center}
\begin{tabular}{p{2.6cm}p{6.6cm}p{4.3cm}}
\toprule
Postulat & Énoncé & Chapitre\\
\midrule
1. \textbf{État} & un vecteur unitaire d'un espace $\C^N$ (à une phase globale près) & ch.~\ref{ch:qubit}\\
2. \textbf{Mesure} & résultat $k$ avec probabilité $|\langle u_k|\psi\rangle|^2$, puis l'état devient $\ket{u_k}$ & ch.~\ref{ch:qubit}, \ref{ch:algebre}\\
3. \textbf{Évolution} & $\ket\psi\mapsto U\ket\psi$ avec $U$ unitaire, $U=e^{-iHt}$ pour un hamiltonien $H$ & ch.~\ref{ch:algebre}, \ref{ch:hamiltoniens}\\
4. \textbf{Composition} & $n$ qubits : produit tensoriel, espace de dimension $2^n$ & ch.~\ref{ch:multi}\\
\bottomrule
\end{tabular}
\end{center}

\exercices
\begin{exercice}[nombres complexes]
Soit $z_1=\tfrac1{\sqrt2}e^{i\pi/4}$ et $z_2=\tfrac1{\sqrt2}e^{i5\pi/4}$. Calculer $|z_1|^2$, $|z_2|^2$, $|z_1+z_2|^2$ et $|z_1|^2+|z_2|^2$. Lequel des deux derniers est une « somme d'amplitudes », lequel une « somme de probabilités » ?
\end{exercice}
\begin{exercice}[trois polariseurs]
Un photon sort d'un polariseur à $0^\circ$. Il traverse ensuite un polariseur à $30^\circ$, puis un polariseur à $90^\circ$. Quelle est la probabilité qu'il sorte du dernier ? Et sans le polariseur à $30^\circ$ ?
\end{exercice}

\chapter{Le qubit}\label{ch:qubit}
\iterbanner{V0}

On formalise ici ce que la polarisation du photon (ch.~\ref{ch:bases}) a montré : un système à deux résultats possibles, dont l'état est un vecteur, que la mesure perturbe. On parlera désormais de \emph{qubit} sans se soucier du support physique.

\section{État d'un qubit et notation de Dirac}
\begin{definition}[ (postulat 1)]
L'état d'un qubit est un vecteur unitaire de $\C^2$ :
\[\ket\psi=\alpha\ket0+\beta\ket1,\qquad \ket0=\begin{pmatrix}1\\0\end{pmatrix},\ \ket1=\begin{pmatrix}0\\1\end{pmatrix},\qquad |\alpha|^2+|\beta|^2=1.\]
$(\ket0,\ket1)$ est la \emph{base de calcul}. Deux vecteurs qui ne diffèrent que d'un facteur $e^{i\chi}$ représentent le même état physique.
\end{definition}
\paragraph{Notation de Dirac.} C'est simplement de l'algèbre linéaire avec une typographie commode :
\begin{itemize}
\item \emph{ket} $\ket\psi$ : vecteur colonne ;
\item \emph{bra} $\bra\psi=\ket\psi^\dagger$ : vecteur ligne conjugué ($\dagger$ = transposée conjuguée) ;
\item produit hermitien $\braket\phi\psi=\sum_i\overline{\phi_i}\psi_i$, linéaire à droite, antilinéaire à gauche ;
\item $\ket\psi\bra\phi$ : une \emph{matrice} (produit extérieur). Par exemple $\ket0\bra0=\begin{pmatrix}1&0\\0&0\end{pmatrix}$ est le projecteur sur $\ket0$.
\end{itemize}
Norme : $\braket\psi\psi=|\alpha|^2+|\beta|^2=1$.

\paragraph{Exemple chiffré.} $\ket\psi=\tfrac35\ket0+\tfrac{4i}{5}\ket1$ est bien normé : $\tfrac9{25}+\tfrac{16}{25}=1$. Son bra est $\bra\psi=\tfrac35\bra0-\tfrac{4i}{5}\bra1$ (on conjugue). Avec $\ket\phi=\ket1$, on obtient $\braket\phi\psi=\tfrac{4i}5$ et $|\braket\phi\psi|^2=\tfrac{16}{25}$ : c'est la probabilité de mesurer 1 (section suivante).

\paragraph{Pourquoi des nombres complexes ?} Avec des amplitudes réelles, on ne disposerait que des signes $\pm$ pour faire interférer les chemins. La phase complexe $e^{i\varphi}$ donne un continuum de « directions » possibles, et les évolutions physiques (équation de Schrödinger, ch.~\ref{ch:hamiltoniens}) produisent naturellement de telles phases.

\section{La mesure : règle de Born}
\begin{definition}[ (postulat de la mesure, base de calcul)]
Mesurer $\ket\psi=\alpha\ket0+\beta\ket1$ donne le résultat 0 avec probabilité $|\alpha|^2=|\braket0\psi|^2$ et 1 avec probabilité $|\beta|^2=|\braket1\psi|^2$. Après la mesure, l'état \emph{est} $\ket0$ ou $\ket1$ selon le résultat (réduction du paquet d'ondes) : l'information sur $\alpha,\beta$ est perdue.
\end{definition}
Conséquences pratiques :
\begin{itemize}
\item une seule mesure ne révèle presque rien de $\ket\psi$ (un bit) ;
\item pour estimer $|\alpha|^2$, il faut \emph{préparer et mesurer} l'état $N$ fois. Ce sont les \textbf{shots}. La fréquence observée a un écart-type $\sqrt{p(1-p)/N}$ ;
\item la mesure est destructive : on ne peut pas « regarder » l'état en cours de calcul sans le modifier.
\end{itemize}

\section{Phase globale et phase relative}\label{sec:phases}
Écrivons les amplitudes sous forme polaire, $\alpha=|\alpha|e^{i\varphi_0}$ et $\beta=|\beta|e^{i\varphi_1}$. L'état contient quatre nombres réels, dont seuls deux comptent :
\[
\ket\psi=e^{i\varphi_0}\Big(\,|\alpha|\,\ket0+|\beta|\,e^{i(\varphi_1-\varphi_0)}\ket1\Big).
\]
\begin{itemize}
\item La norme impose $|\alpha|^2+|\beta|^2=1$ : un seul réel pour les modules, qui fixe les probabilités en base de calcul.
\item Le facteur commun $e^{i\varphi_0}$ est la \textbf{phase globale}.
\item La différence $\varphi=\varphi_1-\varphi_0$ est la \textbf{phase relative}.
\end{itemize}

\paragraph{La phase globale ne se voit jamais.} Pour n'importe quel état de référence $\ket u$, on a
\[|\braket{u}{e^{i\chi}\psi}|^2=|e^{i\chi}|^2\,|\braket u\psi|^2=|\braket u\psi|^2,\qquad\text{car } |e^{i\chi}|=1.\] Aucune mesure, dans aucune base, ne la détecte : $\ket\psi$ et $e^{i\chi}\ket\psi$ sont \emph{le même état physique}, écrit de deux façons.

\paragraph{La phase relative se voit, mais pas en base de calcul.} Les probabilités en base de calcul, $|\alpha|^2$ et $|\beta|^2$, ne dépendent pas de $\varphi$. Pourtant, $\varphi$ change l'état physique. Les quatre états suivants ont tous les mêmes modules, $\tfrac1{\sqrt2}$ et $\tfrac1{\sqrt2}$, et ne diffèrent que par $\varphi$ :
\begin{align*}
\ket{+}&=\tfrac{1}{\sqrt2}(\ket0+\ket1)\quad (\varphi=0), &
\ket{+i}&=\tfrac{1}{\sqrt2}(\ket0+i\ket1)\quad (\varphi=\tfrac\pi2),\\
\ket{-}&=\tfrac{1}{\sqrt2}(\ket0-\ket1)\quad (\varphi=\pi), &
\ket{-i}&=\tfrac{1}{\sqrt2}(\ket0-i\ket1)\quad (\varphi=\tfrac{3\pi}2).
\end{align*}

\begin{piege}[ : le signe dans $\ket-$ et dans $-\ket-$]
Ce qui est écrit \emph{à l'intérieur} d'un ket est une simple \textbf{étiquette}, un nom. $\ket-$ est le nom du vecteur $\tfrac1{\sqrt2}\begin{pmatrix}1\\-1\end{pmatrix}$, et $\ket+$ celui de $\tfrac1{\sqrt2}\begin{pmatrix}1\\1\end{pmatrix}$. On aurait pu les appeler $\ket{\nearrow}$ et $\ket{\searrow}$.

Ce qui est écrit \emph{devant} un ket est un \textbf{nombre} qui multiplie le vecteur. Ainsi
\[
-\ket-=(-1)\times\tfrac1{\sqrt2}\begin{pmatrix}1\\-1\end{pmatrix}=\tfrac1{\sqrt2}\begin{pmatrix}-1\\1\end{pmatrix}=\tfrac1{\sqrt2}(-\ket0+\ket1).
\]
Comme $-1=e^{i\pi}$, $-\ket-$ ne diffère de $\ket-$ que par une phase \emph{globale}. Les deux vecteurs sont différents, mais représentent le même état physique. À l'inverse, $\ket-$ et $\ket+$ ne diffèrent que par le signe \emph{d'une seule} composante : c'est une phase \emph{relative}, et ce sont deux états physiques différents.
\end{piege}

\paragraph{Ce que la phase relative est physiquement : la polarisation.} Revenons au photon du chapitre~\ref{ch:bases}, avec $\ket0=$ polarisation horizontale ($0^\circ$) et $\ket1=$ verticale ($90^\circ$). Le vecteur $(\alpha,\beta)$ est exactement le \emph{vecteur de Jones} de l'optique classique : les deux composantes du champ électrique d'une onde lumineuse sont
\[
E_x(t)=|\alpha|\cos(\omega t),\qquad E_y(t)=|\beta|\cos(\omega t-\varphi).
\]
\begin{itemize}
\item La \textbf{phase relative} $\varphi$ est le déphasage entre les deux oscillations. Il fixe la \emph{forme} décrite par la pointe du champ : un segment à $45^\circ$ ($\varphi=0$), un cercle ($\varphi=\pi/2$), un segment à $135^\circ$ ($\varphi=\pi$), une ellipse sinon (figure~\ref{fig:jones}). Ce sont des lumières physiquement différentes.
\item La \textbf{phase globale} revient à décaler l'origine des temps, $t\to t+t_0$. La pointe du champ parcourt exactement la même courbe : rien de physique ne change. C'est l'intuition à garder : une phase globale est un choix d'horloge, pas une propriété du système.
\end{itemize}

\begin{figure}[htbp]
\centering
\begin{tikzpicture}[scale=0.85]
\foreach \phi/\lab/\nom [count=\i] in {0/{$\varphi=0$}/{$\ket+$ : linéaire $45^\circ$},
                                     45/{$\varphi=\pi/4$}/{elliptique},
                                     90/{$\varphi=\pi/2$}/{$\ket{+i}$ : circulaire},
                                     180/{$\varphi=\pi$}/{$\ket-$ : linéaire $135^\circ$}}{
  \begin{scope}[xshift={(\i-1)*3.9cm}]
    \draw[->,cmuted] (-1.35,0) -- (1.45,0) node[right,font=\scriptsize,cink] {$E_x$};
    \draw[->,cmuted] (0,-1.35) -- (0,1.45) node[above,font=\scriptsize,cink] {$E_y$};
    \draw[cblue,very thick] plot[domain=0:360,samples=73,smooth] ({cos(\x)},{cos(\x-\phi)});
    \fill[corange] ({cos(0)},{cos(-\phi)}) circle (2pt);
    \draw[corange,thick,-{Stealth[length=2mm]}] plot[domain=0:40,samples=10,smooth] ({cos(\x)},{cos(\x-\phi)});
    \node[font=\small] at (0,-1.75) {\lab};
    \node[font=\footnotesize,cmuted] at (0,-2.2) {\nom};
  \end{scope}}
\end{tikzpicture}
\caption{Trajectoire de la pointe du champ électrique pour $|\alpha|=|\beta|=\frac1{\sqrt2}$ et quatre phases relatives. Les quatre lumières passent toutes à 50~\% un polariseur horizontal, mais ce sont des lumières différentes. Une phase globale déplacerait seulement le point orange (l'instant $t=0$) le long de la même courbe.}\label{fig:jones}
\end{figure}

Un photon unique dans l'état $\ket+$ est donc un photon polarisé à $45^\circ$, et un photon dans l'état $\ket{+i}$ est polarisé circulairement. Le tableau suivant donne la probabilité de passer trois filtres différents, ce qui revient à mesurer dans trois bases différentes (§\ref{sec:autre-base}) :
\begin{center}
\begin{tabular}{lccc}
\toprule
État & polariseur $0^\circ$ (base $\ket0,\ket1$) & polariseur $45^\circ$ (base $\ket\pm$) & filtre circulaire (base $\ket{\pm i}$) \\
\midrule
$\ket+$ & $1/2$ & $1$ & $1/2$ \\
$\ket-$ & $1/2$ & $0$ & $1/2$ \\
$\ket{+i}$ & $1/2$ & $1/2$ & $1$ \\
$\ket{-i}$ & $1/2$ & $1/2$ & $0$ \\
\bottomrule
\end{tabular}
\end{center}
La première colonne ne distingue rien : c'est la seule « question » qui ignore $\varphi$. Les deux autres révèlent la phase relative.

\paragraph{Pourquoi c'est utile : l'interférence.} Appliquons la porte $H$ (chapitre~1) à un état quelconque :
\[
H\ket\psi=\alpha\,\frac{\ket0+\ket1}{\sqrt2}+\beta\,\frac{\ket0-\ket1}{\sqrt2}
=\frac{\alpha+\beta}{\sqrt2}\ket0+\frac{\alpha-\beta}{\sqrt2}\ket1 .
\]
Le résultat 0 est désormais atteint par \emph{deux chemins}, l'un venant de $\ket0$ (amplitude $\alpha$), l'autre de $\ket1$ (amplitude $\beta$), et \emph{les amplitudes s'additionnent avant d'être mises au carré}. Avec $|\alpha|=|\beta|=\tfrac1{\sqrt2}$ :
\[
P(0)=\frac{|\alpha+\beta|^2}{2}=\frac{1+\cos\varphi}{2}=\cos^2\frac\varphi2,\qquad P(1)=\sin^2\frac\varphi2 .
\]
C'est exactement la figure de franges d'une expérience de Young, ou l'addition de deux houles de même amplitude : en phase, elles se renforcent ($\varphi=0$, $P(0)=1$) ; en opposition, elles s'annulent ($\varphi=\pi$, $P(0)=0$). La phase relative, invisible tant que les deux chemins restent séparés, devient une \emph{probabilité} dès qu'on les fait se recombiner (figure~\ref{fig:interference}).

\begin{figure}[htbp]
\centering
\begin{tikzpicture}
\begin{axis}[width=10cm,height=5cm,xmin=0,xmax=6.2832,ymin=-0.05,ymax=1.22,
  xtick={0,1.5708,3.1416,4.7124,6.2832},xticklabels={$0$,$\pi/2$,$\pi$,$3\pi/2$,$2\pi$},
  xlabel={phase relative $\varphi$},ylabel={probabilité},legend pos=outer north east,
  legend style={font=\small,draw=none},tick label style={font=\small}]
\addplot[cmuted,thick,dashed,domain=0:6.2832] {0.5};
\addlegendentry{$P(0)$ sans $H$}
\addplot[cblue,very thick,domain=0:6.2832,samples=100] {cos(deg(x)/2)^2};
\addlegendentry{$P(0)$ après $H$}
\addplot[corange,very thick,domain=0:6.2832,samples=100] {sin(deg(x)/2)^2};
\addlegendentry{$P(1)$ après $H$}
\node[font=\scriptsize,anchor=south] at (axis cs:0.25,1.0) {$\ket+$};
\node[font=\scriptsize,anchor=south] at (axis cs:1.5708,0.5) {$\ket{+i}$};
\node[font=\scriptsize,anchor=south] at (axis cs:3.1416,1.0) {$\ket-$};
\end{axis}
\end{tikzpicture}
\caption{Pour $\tfrac1{\sqrt2}(\ket0+e^{i\varphi}\ket1)$, mesurer directement donne toujours 50/50 (tirets) : la phase est cachée. Après $H$, les deux chemins interfèrent et la phase se lit dans les probabilités.}\label{fig:interference}
\end{figure}

\begin{resultat}
\begin{itemize}
\item Phase globale : un choix d'horloge, invisible pour toute mesure. On l'ignore.
\item Phase relative : une propriété physique (forme de la polarisation, direction du spin). Elle est invisible en base de calcul, mais une porte qui fait interférer les chemins, comme $H$, la convertit en probabilités.
\item Un calcul quantique \emph{écrit} l'information dans les phases, puis la fait \emph{ressortir} par interférence avant la mesure.
\end{itemize}
\end{resultat}

\paragraph{Lien avec QAOA, sur un qubit.} Prenons un coût $E(0)=0$, $E(1)=1$. QAOA part de $\ket+$, puis multiplie chaque composante par $e^{-i\gamma E(x)}$ :
\[
\tfrac1{\sqrt2}\big(\ket0+e^{-i\gamma}\ket1\big).
\]
Le coût est maintenant écrit dans la phase relative, $\varphi=-\gamma$, mais une mesure immédiate donnerait 50/50 quel que soit $\gamma$ : l'opération n'a rien appris de visible. Il faut ensuite une étape d'interférence, le \emph{mélangeur}, qui joue le rôle de $H$ ici. Avec $H$, on obtiendrait $P(0)=\cos^2(\gamma/2)$ : la probabilité de sortir la bonne réponse dépend maintenant du coût. Le chapitre~\ref{ch:qaoa} fait la même chose sur les $2^9=512$ profils à la fois.

\section{La sphère de Bloch}
En retirant la phase globale, tout état d'un qubit s'écrit
\[
\ket\psi=\cos\tfrac\theta2\,\ket0+e^{i\varphi}\sin\tfrac\theta2\,\ket1,\qquad \theta\in[0,\pi],\ \varphi\in[0,2\pi[,
\]
c'est-à-dire un point de la sphère unité, de coordonnées $\mathbf r=(\sin\theta\cos\varphi,\ \sin\theta\sin\varphi,\ \cos\theta)$.
\begin{itemize}
\item $\ket0$ et $\ket1$ sont aux pôles.
\item $\ket\pm$ sont sur l'axe $x$, et $\ket{\pm i}=\tfrac1{\sqrt2}(\ket0\pm i\ket1)$ sur l'axe $y$.
\item Deux états \emph{orthogonaux} sont \emph{antipodaux}.
\item Les portes à un qubit sont des \textbf{rotations} de la sphère (ch.~\ref{ch:algebre}).
\end{itemize}

\begin{figure}[htbp]
\centering
\begin{tikzpicture}[x={(-1.147cm,-0.491cm)},y={(1.638cm,-0.344cm)},z={(0cm,1.91cm)},>={Stealth[length=2mm]}]
\draw[cink,thin] (0,0,0) circle (2cm);
\draw[cmuted,thin] plot[domain=0:360,samples=90,smooth] ({cos(\x)},{sin(\x)},0);
\draw[->,cmuted] (0,0,0) -- (1.3,0,0) node[anchor=north east,cink] {$x$};
\draw[->,cmuted] (0,0,0) -- (0,1.3,0) node[anchor=west,cink] {$y$};
\draw[->,cmuted] (0,0,0) -- (0,0,1.3) node[anchor=south west,cink] {$z$};
\fill[cink] (0,0,1) circle (1.5pt) node[anchor=south east] {$\ket0$};
\fill[cink] (0,0,-1) circle (1.5pt) node[anchor=north] {$\ket1$};
\fill[cink] (1,0,0) circle (1.5pt) node[anchor=south east] {$\ket+$};
\fill[cink] (-1,0,0) circle (1.5pt) node[anchor=south west] {$\ket-$};
\fill[cink] (0,1,0) circle (1.5pt) node[anchor=north west] {$\ket{+i}$};
\fill[cink] (0,-1,0) circle (1.5pt) node[anchor=east] {$\ket{-i}$};
% état (theta = 50°, phi = 60°)
\coordinate (P) at ({sin(50)*cos(60)},{sin(50)*sin(60)},{cos(50)});
\coordinate (Pe) at ({sin(50)*cos(60)},{sin(50)*sin(60)},0);
\draw[->,cblue,very thick] (0,0,0) -- (P) node[anchor=south west] {$\ket\psi$};
\draw[cblue,dashed] (P) -- (Pe); \draw[cblue,dashed] (0,0,0) -- (Pe);
\draw[corange,thick] plot[domain=0:50,samples=20] ({0.45*sin(\x)*cos(60)},{0.45*sin(\x)*sin(60)},{0.45*cos(\x)});
\node[corange] at ({0.62*sin(25)*cos(60)},{0.62*sin(25)*sin(60)},{0.62*cos(25)}) {$\theta$};
\draw[corange,thick] plot[domain=0:60,samples=20] ({0.4*cos(\x)},{0.4*sin(\x)},0);
\node[corange] at ({0.58*cos(30)},{0.58*sin(30)},0) {$\varphi$};
\end{tikzpicture}
\caption{Sphère de Bloch. La latitude $\theta$ fixe les probabilités de mesure ($\cos^2\frac\theta2$ pour $\ket0$) ; la longitude $\varphi$ est la phase relative, invisible en base de calcul mais essentielle pour les interférences.}\label{fig:bloch}
\end{figure}

\section{Mesurer dans une autre base}\label{sec:autre-base}
\paragraph{Une base de mesure, c'est une question posée.} Physiquement, mesurer consiste à faire passer le système dans un appareil \emph{orienté} : un polariseur tourné d'un certain angle, ou l'aimant de Stern-Gerlach orienté selon $z$, $x$ ou $y$. L'orientation de l'appareil choisit la \emph{question} posée au système (« es-tu horizontal ou vertical ? », « es-tu à $45^\circ$ ou à $135^\circ$ ? »). Chaque orientation correspond à une base orthonormée $(\ket{u_0},\ket{u_1})$ : ce sont les deux états qui donnent chacun une réponse certaine à cette question.

\begin{definition}[ (mesure dans une base quelconque)]
Mesurer $\ket\psi$ dans la base orthonormée $(\ket{u_0},\ket{u_1})$ donne le résultat $k$ avec probabilité $|\braket{u_k}\psi|^2$, et laisse le système dans l'état $\ket{u_k}$.
\end{definition}
La règle de Born de la section précédente est le cas particulier $\ket{u_k}=\ket k$.

\paragraph{Exemple chiffré : un même état, trois questions.} Reprenons $\ket\psi=\tfrac35\ket0+\tfrac{4i}5\ket1$ et posons-lui les trois questions du tableau précédent :
\begin{itemize}
\item base $\{\ket0,\ket1\}$ : $P(0)=\tfrac9{25}=0{,}36$ et $P(1)=0{,}64$ ;
\item base $\{\ket+,\ket-\}$ : $\braket+\psi=\tfrac1{\sqrt2}\big(\tfrac35+\tfrac{4i}5\big)$, de module carré $\tfrac12\cdot\tfrac{9+16}{25}=\tfrac12$. Résultat : 50/50 ;
\item base $\{\ket{+i},\ket{-i}\}$ : le bra $\bra{+i}=\tfrac1{\sqrt2}(\bra0-i\bra1)$ est conjugué, donc $\braket{+i}\psi=\tfrac1{\sqrt2}\big(\tfrac35-i\cdot\tfrac{4i}5\big)=\tfrac1{\sqrt2}\cdot\tfrac75$, et $P(+i)=\tfrac{49}{50}=0{,}98$.
\end{itemize}
Chaque question révèle une facette différente de l'état. Dans la base $\{\ket{\pm i}\}$, cet état est presque un état de base, alors qu'il paraît très « superposé » dans les deux autres. En répétant les trois mesures sur de nombreux shots, on reconstruit complètement l'état : c'est la \emph{tomographie}. Les trois réponses correspondent aux trois coordonnées $(x,y,z)$ du point sur la sphère de Bloch, ici $(0;\ 0{,}96;\ -0{,}28)$.

\paragraph{Lecture géométrique : une probabilité, c'est une distance au pôle.} Si $\vec r$ est le point de l'état sur la sphère de Bloch et $\vec u$ celui de l'état de base $\ket u$, alors
\[
P(u)=|\braket u\psi|^2=\frac{1+\vec r\cdot\vec u}{2}=\cos^2\frac\Theta2,
\]
où $\Theta$ est l'angle entre $\vec r$ et $\vec u$ sur la sphère. Ici, $P(+i)=0{,}98$ donne $\vec r\cdot\vec u=y=0{,}96$, donc $\Theta=\arccos0{,}96\approx16^\circ$ : l'état est à $16^\circ$ du pôle $\ket{+i}$. Mais tous les points situés à $16^\circ$ de ce pôle donnent la même probabilité. Une seule base ne fixe donc qu'un \emph{cercle} de positions possibles, pas un point. Les deux autres bases ($x=0$, $z=-0{,}28$) choisissent le point sur ce cercle (figure~\ref{fig:cercle-mesure}).

\begin{figure}[htbp]
\centering
\begin{tikzpicture}
% ---------------- gauche : vue 3D
\begin{scope}[x={(-1.147cm,-0.491cm)},y={(1.638cm,-0.344cm)},z={(0cm,1.91cm)},>={Stealth[length=2mm]}]
\draw[cink,thin] (0,0,0) circle (2cm);
\draw[cmuted,thin] plot[domain=0:360,samples=90,smooth] ({cos(\x)},{sin(\x)},0);
\draw[->,cmuted] (0,0,0) -- (1.3,0,0) node[anchor=north east,cink] {$x$};
\draw[->,cmuted] (0,0,0) -- (0,0,1.3) node[anchor=south west,cink] {$z$};
\draw[->,cmuted] (0,0,0) -- (0,1.55,0) node[anchor=west,cink] {$y$};
\fill[cink] (0,0,1) circle (1.5pt) node[anchor=south east] {$\ket0$};
\fill[cink] (0,0,-1) circle (1.5pt) node[anchor=north] {$\ket1$};
\fill[cink] (0,-1,0) circle (1.5pt) node[anchor=east] {$\ket{-i}$};
\fill[cblue] (0,1,0) circle (1.8pt);
\node[cblue,anchor=north,font=\small] at (0,1.32,-0.08) {$\ket{+i}$};
% cône et cercle des états à 16° du pôle +i (cos 16,26° = 0,96 ; sin = 0,28)
\draw[corange,dashed] (0,0,0) -- (0,0.96,0.28);
\draw[corange,dashed] (0,0,0) -- (0,0.96,-0.28);
\draw[corange,very thick] plot[domain=0:360,samples=60,smooth] ({0.28*cos(\x)},0.96,{0.28*sin(\x)});
\draw[->,cblue,very thick] (0,0,0) -- (0,0.96,-0.28);
\fill[cblue] (0,0.96,-0.28) circle (2pt);
\node[cblue,anchor=north east,font=\small] at (0,0.92,-0.32) {$\ket\psi$};
\end{scope}
\node[corange,font=\small,align=center] at (4.0,1.75) {tous ces états\\ donnent $P(+i)=0{,}98$};
\draw[corange,-{Stealth[length=2mm]}] (3.7,1.3) -- (1.95,0.05);
\node[font=\small] at (0,-2.75) {(a) ce que fixe une seule mesure};
% ---------------- droite : coupe dans le plan (y, z)
\begin{scope}[xshift=8.6cm]
\def\R{2.1}
\draw[cink,thin] (0,0) circle (\R);
\draw[->,cmuted] (-\R-0.3,0) -- (\R+1.0,0) node[right,cink] {$y$};
\draw[->,cmuted] (0,-\R-0.3) -- (0,\R+0.35) node[above,cink] {$z$};
\fill[cink] (0,\R) circle (1.5pt) node[anchor=south west] {$\ket0$};
\fill[cink] (0,-\R) circle (1.5pt) node[anchor=north west] {$\ket1$};
\fill[cink] (-\R,0) circle (1.5pt) node[anchor=south east] {$\ket{-i}$};
\fill[cblue] (\R,0) circle (1.8pt) node[anchor=south west,cblue,xshift=2pt] {$\ket{+i}$};
% le cercle des candidats, vu de profil
\draw[cblue,very thick,->] (0,0) -- ({0.96*\R},{-0.28*\R});
\fill[cblue] ({0.96*\R},{-0.28*\R}) circle (2pt) node[anchor=north west,cblue,font=\small] {$\ket\psi$};
\draw[cmuted,dashed] (0,0) -- ({0.96*\R},{0.28*\R});
\draw[corange,line width=2.2pt] ({0.96*\R},{0.28*\R}) -- ({0.96*\R},{-0.28*\R});
\draw[cmuted,fill=white] ({0.96*\R},{0.28*\R}) circle (2pt);
\fill[cblue] ({0.96*\R},{-0.28*\R}) circle (2pt);
\node[corange,font=\small,align=left,anchor=west] (prof) at ({\R+0.1},{1.7}) {cercle orange\\ vu de profil};
\draw[corange,-{Stealth[length=2mm]}] (prof.south west) ++(0.2,0.05) -- ({0.96*\R+0.08},{0.3*\R});
\draw[corange,thick] (0.95,0) arc (0:-16.26:0.95);
\node[corange,font=\small] at (1.3,-0.2) {$\Theta$};
% projection et lecture de P
\draw[cblue,dashed] ({0.96*\R},{-0.28*\R}) -- ({0.96*\R},0);
\draw[decorate,decoration={brace,amplitude=5pt,mirror}] (-\R,-\R-0.55) -- ({0.96*\R},-\R-0.55)
  node[midway,below=6pt,font=\small] {$1+y=1{,}96$};
\draw[cmuted,thin] ({0.96*\R},0) -- ({0.96*\R},-\R-0.5);
\draw[cmuted,thin] (-\R,0) -- (-\R,-\R-0.5);
\end{scope}
\node[font=\small,align=center] at (8.6,-3.75) {(b) coupe dans le plan $(y,z)$ : $P(+i)=\dfrac{1+y}{2}=\dfrac{1{,}96}{2}=0{,}98$};
\end{tikzpicture}
\caption{Ce qu'apprend une mesure dans la base $\{\ket{+i},\ket{-i}\}$ pour $\ket\psi=\tfrac35\ket0+\tfrac{4i}5\ket1$. (a)~La probabilité $P(+i)=0{,}98$ fixe seulement l'angle $\Theta\approx16^\circ$ au pôle $\ket{+i}$ : l'état est quelque part sur le cercle orange. (b)~Vue en coupe : la probabilité est la projection $y$ du point sur l'axe, ramenée à $[0,1]$. Le cercle orange est vu de profil (segment vertical) ; le point en tirets donne la même probabilité. Les mesures en base $\{\ket0,\ket1\}$ ($z=-0{,}28$) et $\{\ket\pm\}$ ($x=0$) lèvent l'ambiguïté.}\label{fig:cercle-mesure}
\end{figure}

\paragraph{Analogie avec le polariseur, et sa limite.} Un filtre circulaire transmettrait bien 98~\% de l'intensité d'un faisceau dans cet état : c'est la loi de Malus, de même forme en $\cos^2$. À l'échelle d'\emph{un} photon, en revanche, c'est tout ou rien. Dans 98~\% des cas, le photon passe et ressort \emph{exactement} dans l'état $\ket{+i}$ ; dans 2~\% des cas, il est absorbé. La mesure ne conserve pas « 98~\% de $\ket\psi$ » : elle remplace l'état par l'état de base obtenu.

\paragraph{Complémentarité.} Aucun état ne donne une réponse certaine à la fois en base $\{\ket0,\ket1\}$ et en base $\{\ket+,\ket-\}$ : si $P(0)=1$, alors l'état est $\ket0$, et $P(+)=\tfrac12$. Sur la sphère de Bloch, être certain pour une question, c'est être sur un pôle de l'axe correspondant, donc sur l'équateur des deux autres axes. C'est pour cette raison que l'information de phase est « cachée » quand on ne mesure qu'en base de calcul : c'est la question complémentaire.

\paragraph{En pratique, on tourne l'état plutôt que l'appareil.} Un processeur quantique supraconducteur ne sait lire qu'une chose : le niveau d'énergie du qubit, donc la base $\{\ket0,\ket1\}$. Pour poser une autre question, on applique avant la mesure une porte $U$ qui envoie la base voulue sur la base de calcul, $U\ket{u_k}=\ket k$. Comme $U$ est unitaire, elle conserve les produits scalaires :
\[
|\braket{k}{U\psi}|^2=|\braket{U u_k}{U\psi}|^2=|\braket{u_k}\psi|^2 .
\]
Tourner l'état de $U$ revient à tourner l'appareil de $U^{-1}$ : c'est le même jeu qu'entre repère fixe et repère tournant en mécanique, où l'on peut faire tourner soit l'objet, soit l'observateur.

\begin{figure}[htbp]
\centering
\begin{tikzpicture}[x=1cm,y=1cm]
\begin{scope}
  \node[font=\small] at (-0.6,0) {$\ket\psi$};
  \draw[cink] (-0.2,0) -- (1.9,0);
  \meter{1.2}{0}
  \node[font=\small,align=center] at (0.9,-0.9) {base $\{\ket0,\ket1\}$\\ ($Z$)};
\end{scope}
\begin{scope}[xshift=4.2cm]
  \node[font=\small] at (-0.6,0) {$\ket\psi$};
  \draw[cink] (-0.2,0) -- (2.4,0);
  \node[gate] at (0.6,0) {$H$};
  \meter{1.7}{0}
  \node[font=\small,align=center] at (1.1,-0.9) {base $\{\ket+,\ket-\}$\\ ($X$)};
\end{scope}
\begin{scope}[xshift=9.0cm]
  \node[font=\small] at (-0.6,0) {$\ket\psi$};
  \draw[cink] (-0.2,0) -- (3.3,0);
  \node[gate] at (0.6,0) {$S^\dagger$};
  \node[gate] at (1.5,0) {$H$};
  \meter{2.6}{0}
  \node[font=\small,align=center] at (1.5,-0.9) {base $\{\ket{+i},\ket{-i}\}$\\ ($Y$)};
\end{scope}
\end{tikzpicture}
\caption{Mesurer dans trois bases avec un matériel qui ne lit que la base de calcul. $H$ envoie $\ket\pm$ sur $\ket0/\ket1$. $S^\dagger=\mathrm{diag}(1,-i)$ retire d'abord une phase relative $\pi/2$, ce qui envoie $\ket{\pm i}$ sur $\ket\pm$, puis $H$ termine le travail.}\label{fig:mesure-bases}
\end{figure}

\begin{projet}
À la fin d'un circuit QAOA, on mesure les 9 qubits en base de calcul : chaque shot renvoie une chaîne de 9 bits, c'est-à-dire \emph{un profil} $(m,p,t)$ de la grille. La règle de Born dit avec quelle probabilité chaque profil sort. QAOA est réussi si les profils de faible $C_D$ sortent souvent.

On ne mesure qu'en base de calcul, car c'est la question qui nous intéresse : « quel profil ? ». Les étapes du circuit qui précèdent la mesure ont précisément pour rôle de convertir l'information de coût, écrite dans les phases relatives, en différences de probabilités dans cette base. La fonction de coût $H_C$ étant diagonale dans la base de calcul, son espérance $\expval{H_C}$, que l'optimiseur classique minimise, s'estime aussi avec ces seules mesures.
\end{projet}

\exercices
\begin{exercice}
Soit $\ket\psi=\tfrac{1}{\sqrt2}(\ket0+i\ket1)$. (a) Probabilités de mesure en base de calcul. (b) Angles $(\theta,\varphi)$ et position sur la sphère de Bloch. (c) Probabilités de mesure en base $\{\ket+,\ket-\}$.
\end{exercice}
\begin{exercice}
On rappelle que $-\ket-$ désigne le vecteur $\ket-$ multiplié par le nombre $-1$ (voir l'encadré « Piège » du §\ref{sec:phases}). (a) Écrire $-\ket-$ en composantes. Montrer que $\ket-$ et $-\ket-$ sont physiquement identiques. (b) Montrer que $\ket+$ et $\ket-$ ne le sont pas. Proposer une procédure (porte + mesure) qui les distingue avec certitude. (c) En polarisation, à quelles lumières correspondent $\ket+$ et $\ket-$ ? Quel filtre les distingue ?
\end{exercice}

\chapter{Observables et évolution : l'algèbre linéaire du quantique}\label{ch:algebre}
\iterbanner{V0}

Toute la mécanique quantique de ce cours se joue en dimension finie ($\C^{2^n}$). Il n'y a pas d'analyse fonctionnelle, seulement des matrices.

\section{Espace de Hilbert, adjoint, bases}
L'expression « espace de Hilbert » fait peur, mais ici elle désigne simplement $\C^N$, avec $N=2^n$ pour $n$ qubits, muni d'un produit scalaire adapté aux nombres complexes. Cette section rassemble les quatre outils d'algèbre linéaire dont tout le reste du cours a besoin.

\paragraph{1. Le produit hermitien.} Pour deux vecteurs $\ket\phi$ et $\ket\psi$ de composantes $\phi_k$ et $\psi_k$,
\[
\braket\phi\psi=\sum_k\overline{\phi_k}\,\psi_k .
\]
On conjugue les composantes de gauche pour que la norme soit un réel positif : $\braket\psi\psi=\sum_k\overline{\psi_k}\psi_k=\sum_k|\psi_k|^2\ge0$. Sans conjugaison, le vecteur $(1,i)$ aurait une « norme au carré » $1+i^2=0$, ce qui n'a pas de sens. Contrepartie : l'ordre compte, $\braket\psi\phi=\overline{\braket\phi\psi}$. Deux vecteurs sont \textbf{orthogonaux} si $\braket\phi\psi=0$, et une base $(\ket{e_k})$ est \textbf{orthonormée} si $\braket{e_j}{e_k}=\delta_{jk}$.

\paragraph{2. Composantes et projecteurs.} Dans une base orthonormée, la composante de $\ket\psi$ selon $\ket{e_k}$ est $\braket{e_k}\psi$, exactement comme la composante d'une vitesse selon un axe est $\mathbf u\cdot\mathbf e_k$. La matrice
\[
P_k=\ket{e_k}\bra{e_k}
\]
est le \textbf{projecteur} sur $\ket{e_k}$ : $P_k\ket\psi=\ket{e_k}\,\braket{e_k}\psi$ garde la composante selon $\ket{e_k}$ et efface les autres. Projeter deux fois revient à projeter une fois ($P_k^2=P_k$). Avec ce langage, la règle de Born s'écrit
\[
P(k)=|\braket{e_k}\psi|^2=\bra\psi P_k\ket\psi .
\]

\paragraph{Pourquoi le \emph{carré} de la composante ?} Prenons un état à amplitudes réelles, $\ket\psi=\cos\theta\,\ket{e_0}+\sin\theta\,\ket{e_1}$ : un vecteur unitaire du plan, à l'angle $\theta$ de l'axe $\ket{e_0}$ (figure~\ref{fig:pythagore}). Ses composantes $\cos\theta$ et $\sin\theta$ ne peuvent pas être des probabilités : elles peuvent être négatives, et leur somme n'est pas 1 (elle vaut $\sqrt2$ à $45^\circ$). Leurs \emph{carrés}, eux, sont positifs et somment toujours à 1 : c'est le théorème de Pythagore, $\cos^2\theta+\sin^2\theta=1$. La règle de Born exploite exactement cela. Un vecteur d'état est de longueur 1, la mesure choisit des axes orthogonaux, et le carré de la projection sur chaque axe donne la part de probabilité de ce résultat. Avec des composantes complexes, on remplace le carré par le module au carré.

\begin{figure}[htbp]
\centering
\begin{tikzpicture}[scale=2.6]
\draw[cmuted,thin] (0,0) -- (1.2,0);
\draw[cmuted,thin] (0,0) -- (0,1.15);
\draw[cmuted,thin,dashed] (0:1) arc (0:90:1);
\draw[arr] (0,0) -- (1,0) node[below,font=\small] {$\ket{e_0}$};
\draw[arr] (0,0) -- (0,1) node[left,font=\small] {$\ket{e_1}$};
\draw[cblue,very thick,-{Stealth[length=2.5mm]}] (0,0) -- (35:1) node[above right,font=\small] {$\ket\psi$};
\draw[corange,dashed] (35:1) -- ({cos(35)},0);
\draw[corange,dashed] (35:1) -- (0,{sin(35)});
\draw[corange,line width=2pt] (0,0) -- ({cos(35)},0);
\draw[cgreen,line width=2pt] (0,0) -- (0,{sin(35)});
\node[corange,font=\small,below] at ({cos(35)/2},-0.02) {$\braket{e_0}\psi=\cos\theta$};
\node[cgreen,font=\small,left] at (-0.02,{sin(35)/2}) {$\braket{e_1}\psi=\sin\theta$};
\draw[cink] (0.3,0) arc (0:35:0.3); \node[font=\small] at (17:0.4) {$\theta$};
\node[font=\small,align=left,anchor=west] at (1.35,0.55) {longueur de $\ket\psi$ : $1$\\[2pt]
  $P(0)=\cos^2\theta$, $\ P(1)=\sin^2\theta$\\[2pt]
  Pythagore : $P(0)+P(1)=1$};
\end{tikzpicture}
\caption{La règle de Born est le théorème de Pythagore. Les probabilités sont les carrés des projections du vecteur d'état, de longueur 1, sur les axes de la base de mesure.}\label{fig:pythagore}
\end{figure}

\paragraph{3. La relation de fermeture.} La somme des projecteurs sur tous les vecteurs d'une base orthonormée redonne l'identité :
\[
\sum_k\ket{e_k}\bra{e_k}=\Id .
\]
Appliquée à $\ket\psi$, elle dit $\ket\psi=\sum_k\ket{e_k}\braket{e_k}\psi$ : un vecteur est la somme de ses composantes, comme $\mathbf u=\sum_k(\mathbf u\cdot\mathbf e_k)\,\mathbf e_k$. Son intérêt pratique est qu'on peut \emph{insérer} $\Id$ n'importe où dans un calcul pour changer de base. Deux exemples :
\begin{itemize}
\item les probabilités somment à 1 : $\sum_kP(k)=\sum_k\bra\psi\big(\ket{e_k}\bra{e_k}\big)\ket\psi=\bra\psi\Id\ket\psi=1$ ;
\item écrire $\ket0$ dans la base $\{\ket\pm\}$ : $\ket0=\ket+\braket+0+\ket-\braket-0=\tfrac1{\sqrt2}(\ket++\ket-)$.
\end{itemize}

\paragraph{4. L'adjoint.} L'\textbf{adjoint} d'une matrice est sa transposée conjuguée, $A^\dagger=\overline{A}^\top$. C'est l'analogue complexe de la transposée, et il vérifie
\[
\braket{\phi}{A\psi}=\braket{A^\dagger\phi}{\psi}.
\]
On l'utilise sans cesse pour passer d'un ket à un bra : $(\ket\psi)^\dagger=\bra\psi$, $(A\ket\psi)^\dagger=\bra\psi A^\dagger$, et $(AB)^\dagger=B^\dagger A^\dagger$ (l'ordre s'inverse, comme pour la transposée). Exemple :
\[
A=\begin{pmatrix}1&2i\\3&4\end{pmatrix}\quad\Longrightarrow\quad A^\dagger=\begin{pmatrix}1&3\\-2i&4\end{pmatrix}.
\]
Deux familles de matrices, définies par l'adjoint, structurent toute la mécanique quantique :
\begin{center}
\begin{tabularx}{\linewidth}{llXX}
\toprule
Famille & Définition & Propriété clé & Rôle physique \\
\midrule
hermitienne & $A^\dagger=A$ & valeurs propres réelles, vecteurs propres orthonormés & grandeur mesurable (§\ref{sec:observables}) \\
unitaire & $U^\dagger U=\Id$ & conserve les produits scalaires et la norme & évolution, porte (§\ref{sec:unitaires}) \\
\bottomrule
\end{tabularx}
\end{center}

\section{Observables : opérateurs hermitiens}\label{sec:observables}
\paragraph{Le problème.} En physique classique, une grandeur mesurable (énergie, position) est un nombre, fonction de l'état. En quantique, une même mesure peut donner plusieurs résultats, avec des probabilités. Il faut donc un objet qui contienne à la fois :
\begin{enumerate}
\item la \emph{question} posée, c'est-à-dire la base de mesure $(\ket{u_k})$, autrement dit l'orientation de l'appareil (§\ref{sec:autre-base}) ;
\item la \emph{valeur} affichée par l'appareil pour chaque réponse : un nombre réel $a_k$, par exemple $+1$ ou $-1$ pour un spin, une énergie en joules, ou un $C_D$.
\end{enumerate}
La matrice $A=\sum_ka_k\ket{u_k}\bra{u_k}$ fait exactement cela : ses vecteurs propres sont les états de la base, et ses valeurs propres sont les valeurs affichées. Une observable n'est rien d'autre que cet emballage.

\paragraph{Pourquoi cette forme, et pas une autre ?} On cherche une matrice qui \emph{encode} l'appareil : appliquée à l'état $\ket{u_k}$, pour lequel l'appareil affiche $a_k$ à coup sûr, elle doit simplement le multiplier par $a_k$. On impose donc $A\ket{u_k}=a_k\ket{u_k}$ pour tout $k$. Une application linéaire est entièrement déterminée par son action sur une base, donc cette condition fixe $A$ de façon unique. Pour l'écrire, on insère la relation de fermeture $\Id=\sum_k\ket{u_k}\bra{u_k}$ :
\[
A=A\,\Id=\sum_kA\ket{u_k}\bra{u_k}=\sum_ka_k\ket{u_k}\bra{u_k}.
\]
Inversement, on vérifie que cette matrice convient. Par orthonormalité, $\braket{u_k}{u_j}=\delta_{kj}$, donc $A\ket{u_j}=\sum_ka_k\ket{u_k}\braket{u_k}{u_j}=a_j\ket{u_j}$. La formule se lit comme une recette en trois temps :
\begin{enumerate}
\item projeter l'état sur chaque axe de mesure ($\bra{u_k}\cdot$) ;
\item multiplier chaque composante par la valeur affichée $a_k$ ;
\item recombiner ($\ket{u_k}\cdot$).
\end{enumerate}
C'est exactement une matrice diagonale $\mathrm{diag}(a_0,a_1,\dots)$, écrite dans la base de l'appareil plutôt que dans la base de calcul. Si l'appareil mesure en base de calcul ($\ket{u_k}=\ket k$), on retrouve littéralement $A=\mathrm{diag}(a_0,a_1,\dots)$. Pour $a_0=1$ et $a_1=-1$, c'est $Z$.

\begin{figure}[htbp]
\centering
\begin{tikzpicture}
\node[font=\small] (psi) at (0,0) {$\ket\psi$};
\node[flow,minimum width=2.6cm,minimum height=1.5cm] (app) at (2.8,0) {appareil orienté\\ selon la base $(\ket{u_k})$};
\draw[arr] (psi) -- (app);
\node[font=\small,anchor=west] (r0) at (5.3,0.55) {affiche $a_0$, avec probabilité $|\braket{u_0}\psi|^2$};
\node[font=\small,anchor=west] (r1) at (5.3,-0.55) {affiche $a_1$, avec probabilité $|\braket{u_1}\psi|^2$};
\draw[arr] (app.east) ++(0,0.2) -- (r0.west);
\draw[arr] (app.east) ++(0,-0.2) -- (r1.west);
\node[cblue,font=\small] at (2.8,-1.35) {$A=a_0\ket{u_0}\bra{u_0}+a_1\ket{u_1}\bra{u_1}$};
\end{tikzpicture}
\caption{Une observable regroupe une base de mesure (l'orientation de l'appareil) et les valeurs affichées pour chaque résultat.}\label{fig:observable}
\end{figure}

\paragraph{Pourquoi hermitienne ?} Le théorème spectral dit qu'une matrice est hermitienne si et seulement si elle s'écrit $\sum_ka_k\ket{u_k}\bra{u_k}$ avec des $a_k$ \emph{réels} et des $\ket{u_k}$ \emph{orthonormés}. Ce sont exactement les deux exigences physiques : un appareil affiche des nombres réels, et ses différentes réponses doivent correspondre à des états parfaitement distinguables, donc orthogonaux.

\begin{definition}[ (postulat des observables)]
Une grandeur mesurable est représentée par une matrice \textbf{hermitienne} $A=A^\dagger=\sum_k a_k\ket{u_k}\bra{u_k}$. Mesurer $A$ sur $\ket\psi$ donne la valeur $a_k$ avec probabilité $|\braket{u_k}\psi|^2$, et laisse le système dans $\ket{u_k}$. La \textbf{valeur moyenne} sur de nombreux shots vaut
\[\expval{A}_\psi=\bra\psi A\ket\psi=\sum_k a_k\,|\braket{u_k}\psi|^2 .\]
\end{definition}
\paragraph{Intuition.} Une \emph{observable} est une « question » qu'on pose au système. Ses valeurs propres sont les \emph{réponses possibles} ; ses vecteurs propres sont les états pour lesquels la réponse est certaine. Si l'état n'est pas un vecteur propre, la réponse est aléatoire, et l'on ne peut prédire que sa loi de probabilité, en particulier sa moyenne.

\subsection*{Valeur affichée et probabilité : deux choses à ne pas mélanger}
Une observable $A=\sum_ka_k\ket{u_k}\bra{u_k}$ contient deux ingrédients qui jouent des rôles \emph{complètement séparés} :
\begin{itemize}
\item la \textbf{base} $(\ket{u_k})$ décide \emph{quels} résultats sont possibles et \emph{avec quelle probabilité} ils sortent : $P(k)=|\braket{u_k}\psi|^2$. Les $a_k$ n'y apparaissent pas ;
\item les \textbf{valeurs} $a_k$ sont seulement les nombres que l'appareil \emph{inscrit} pour chaque résultat. Elles n'interviennent que lorsqu'on calcule une moyenne (ou une variance) : $\expval A=\sum_ka_k\,P(k)$.
\end{itemize}

\paragraph{L'analogie du dé.} Un dé équilibré tombe sur chaque face avec probabilité $\tfrac16$. Les nombres peints sur les faces sont les « valeurs affichées ». Si l'on repeint les faces avec 10, 20, \dots, 60, ou avec des 0 et des 1, aucune probabilité ne change : le dé roule exactement pareil. Seule la moyenne change. Il en va de même ici : la physique (l'état et l'orientation de l'appareil) fixe les probabilités, et les $a_k$ ne sont qu'une étiquette écrite sur le cadran.

\begin{figure}[htbp]
\centering
\begin{tikzpicture}[node distance=0.9cm]
\node[flow] (psi) {état\\ $\ket\psi$};
\node[flow,right=of psi] (amp) {amplitudes\\ $\braket{u_k}\psi\in\C$};
\node[flow,right=of amp] (prob) {probabilités\\ $P(k)=|\braket{u_k}\psi|^2$};
\node[flow,right=of prob] (moy) {moyenne\\ $\expval A=\sum_ka_kP(k)$};
\draw[arr] (psi) -- node[above,font=\scriptsize] {base} (amp);
\draw[arr] (amp) -- node[above,font=\scriptsize] {$|\cdot|^2$} (prob);
\draw[arr] (prob) -- node[above,font=\scriptsize] {$\times a_k$} (moy);
\draw[decorate,decoration={brace,amplitude=5pt,mirror},cblue] ([yshift=-3mm]psi.south west) -- ([yshift=-3mm]prob.south east)
  node[midway,below=6pt,font=\small,cblue] {physique : état et orientation de l'appareil};
\draw[decorate,decoration={brace,amplitude=5pt,mirror},corange] ([yshift=-3mm]moy.south west) -- ([yshift=-3mm]moy.south east)
  node[midway,below=6pt,font=\small,corange,align=center] {étiquettes $a_k$};
\end{tikzpicture}
\caption{Les valeurs $a_k$ n'entrent qu'à la dernière étape. Elles ne modifient jamais les probabilités.}\label{fig:chaine-mesure}
\end{figure}

\paragraph{Exemple : un même appareil, trois cadrans.} On mesure $\ket\psi=\tfrac35\ket0+\tfrac{4i}5\ket1$ en base de calcul avec trois cadrans différents :
\begin{center}
\begin{tabular}{lccccc}
\toprule
Cadran & $a_0$ & $a_1$ & $P(0)$, $P(1)$ & moyenne & formule fausse $|a_k\braket{u_k}\psi|^2$ \\
\midrule
spin, observable $Z$ & $+1$ & $-1$ & $0{,}36$ ; $0{,}64$ & $-0{,}28$ & $0{,}36$ ; $0{,}64$ (par chance) \\
bit, observable $\hat b=\ket1\bra1$ & $0$ & $1$ & $0{,}36$ ; $0{,}64$ & $0{,}64$ & $0$ ; $0{,}64$ (somme $\neq1$) \\
énergie en unités arbitraires & $10$ & $20$ & $0{,}36$ ; $0{,}64$ & $16{,}4$ & $36$ ; $256$ (absurde) \\
\bottomrule
\end{tabular}
\end{center}
Les probabilités sont les mêmes sur les trois lignes : c'est le même appareil, seul le cadran change. La formule $|a_k\braket{u_k}\psi|^2$ échoue de trois façons :
\begin{itemize}
\item elle dépend du choix d'unités ;
\item elle ne somme pas à 1 ;
\item elle donne une probabilité nulle à tout résultat affiché « 0 », alors qu'un polariseur qui affiche 0 quand le photon est absorbé voit bien des photons absorbés.
\end{itemize}
Avec $Z$, elle donne le bon résultat uniquement parce que $a_k^2=1$, ce qui peut tromper.

\begin{piege}[ : $A\ket\psi$ n'est ni l'état après la mesure, ni un vecteur de probabilités]
La quantité $a_k\braket{u_k}\psi$ existe bien : c'est la composante du vecteur $A\ket\psi$ selon $\ket{u_k}$, car $\braket{u_k}{A\psi}=a_k\braket{u_k}\psi$. Mais $A\ket\psi$ n'a pas de sens physique direct :
\begin{itemize}
\item ce n'est pas l'état après la mesure, qui est $\ket{u_k}$ ;
\item il n'est même pas normé. Sa norme au carré vaut $\|A\ket\psi\|^2=\sum_ka_k^2\,|\braket{u_k}\psi|^2=\expval{A^2}$, le moment d'ordre 2 de la mesure.
\end{itemize}
Les nombres $|a_k\braket{u_k}\psi|^2$ sont donc les termes de $\expval{A^2}$, pas des probabilités. Appliquer une observable à un état est un outil de calcul, pour obtenir $\bra\psi A\ket\psi$ ou $\expval{A^2}$. Ce n'est pas ce que fait l'appareil au système.
\end{piege}

\paragraph{D'où vient la formule de la moyenne ?} On remplace $A$ par sa décomposition spectrale :
\[
\bra\psi A\ket\psi=\sum_ka_k\,\braket\psi{u_k}\braket{u_k}\psi=\sum_ka_k\,\overline{\braket{u_k}\psi}\,\braket{u_k}\psi=\sum_ka_k\,|\braket{u_k}\psi|^2 .
\]
C'est la définition usuelle d'une espérance, $\sum(\text{valeur})\times(\text{probabilité})$. L'intérêt de la forme $\bra\psi A\ket\psi$ est qu'elle se calcule par un simple produit matriciel, \emph{sans} diagonaliser $A$.

\paragraph{Exemple : construire une observable à partir d'un appareil.} Un aimant de Stern-Gerlach orienté selon $x$ distingue les états $\ket+$ (déviés vers le haut, valeur $+1$) et $\ket-$ (vers le bas, valeur $-1$). L'observable correspondante est
\[
\ket+\bra+-\ket-\bra-=\frac12\begin{pmatrix}1&1\\1&1\end{pmatrix}-\frac12\begin{pmatrix}1&-1\\-1&1\end{pmatrix}=\begin{pmatrix}0&1\\1&0\end{pmatrix}=X .
\]
De même, l'aimant orienté selon $z$ donne $\ket0\bra0-\ket1\bra1=Z$. Les matrices de Pauli ci-dessous ne sont donc pas des objets abstraits : ce sont les observables « spin selon $x$, $y$, $z$ ». Un polariseur qui affiche 1 si le photon passe et 0 sinon correspond, lui, au projecteur $\ket{u_0}\bra{u_0}$.

\paragraph{Exemple chiffré.} Reprenons $\ket\psi=\tfrac35\ket0+\tfrac{4i}5\ket1$ et l'observable $Z=\mathrm{diag}(1,-1)$. Le résultat est $+1$ avec probabilité $\tfrac9{25}$ et $-1$ avec probabilité $\tfrac{16}{25}$, d'où une moyenne $\tfrac9{25}-\tfrac{16}{25}=-\tfrac7{25}$. Par le calcul matriciel : $Z\ket\psi=\tfrac35\ket0-\tfrac{4i}5\ket1$, et $\bra\psi Z\ket\psi=\tfrac35\cdot\tfrac35+\big(-\tfrac{4i}5\big)\big(-\tfrac{4i}5\big)=\tfrac9{25}+\tfrac{16i^2}{25}=-\tfrac7{25}$. Les deux calculs coïncident.

Le même calcul avec $X$ et $Y$ donne
\[
X\ket\psi=\tfrac{4i}5\ket0+\tfrac35\ket1\ \Rightarrow\ \expval X=\tfrac35\cdot\tfrac{4i}5+\big(-\tfrac{4i}5\big)\tfrac35=0,\qquad
Y\ket\psi=\tfrac45\ket0+\tfrac{3i}5\ket1\ \Rightarrow\ \expval Y=\tfrac{12}{25}+\tfrac{12}{25}=\tfrac{24}{25}.
\]
On retrouve les coordonnées $(0;\ 0{,}96;\ -0{,}28)$ du point de Bloch obtenues au §\ref{sec:autre-base} par les trois mesures.

\paragraph{La dispersion, et le nombre de shots.} Une moyenne ne suffit pas : il faut savoir à quel point les résultats fluctuent autour d'elle. L'écart-type de la mesure est
\[
\Delta A=\sqrt{\expval{A^2}-\expval{A}^2}.
\]
Il est nul si et seulement si $\ket\psi$ est un vecteur propre de $A$ (réponse certaine). Pour notre exemple, $Z^2=\Id$, donc $\expval{Z^2}=1$ et $\Delta Z=\sqrt{1-0{,}28^2}=0{,}96$ : les résultats $\pm1$ fluctuent beaucoup (figure~\ref{fig:histo-Z}). Avec $N$ shots, la moyenne empirique a une erreur de l'ordre de $\Delta A/\sqrt N$ : il faut environ $1\,000$ shots pour connaître $\expval Z$ à $\pm0{,}03$ près. C'est le coût de toute estimation sur machine réelle, et en particulier de chaque évaluation de la fonction de coût de QAOA.

\begin{figure}[htbp]
\centering
\begin{tikzpicture}
\begin{axis}[width=9cm,height=5.4cm,ybar,bar width=1.1cm,ymin=0,ymax=0.9,xmin=-1.8,xmax=1.8,
  enlargelimits=false,clip=false,
  xtick={-1,1},xticklabels={$-1$ (état $\ket1$),$+1$ (état $\ket0$)},ylabel={probabilité},
  ytick={0,0.2,0.4,0.6,0.8},yticklabel style={/pgf/number format/use comma},
  tick label style={font=\small},nodes near coords,
  nodes near coords style={font=\small,/pgf/number format/use comma}]
\addplot[fill=cblue!35,draw=cblue] coordinates {(-1,0.64) (1,0.36)};
\draw[corange,very thick] (axis cs:-0.28,0) -- (axis cs:-0.28,0.78);
\node[corange,font=\small,anchor=south] at (axis cs:-0.28,0.78) {$\expval Z=-0{,}28$};
\draw[corange,thick,{Stealth[length=2mm]}-{Stealth[length=2mm]}] (axis cs:-1.24,0.5) -- (axis cs:0.68,0.5);
\node[corange,font=\small,anchor=south west] at (axis cs:-0.2,0.51) {$\pm\Delta Z=\pm0{,}96$};
\end{axis}
\end{tikzpicture}
\caption{Loi des résultats de la mesure de $Z$ sur $\tfrac35\ket0+\tfrac{4i}5\ket1$. Chaque shot donne $+1$ ou $-1$, jamais $-0{,}28$ : la moyenne n'est pas un résultat possible, seulement la limite de la moyenne empirique sur de nombreux shots.}\label{fig:histo-Z}
\end{figure}

\paragraph{Le cas qui nous intéresse : les observables diagonales.} Toute fonction $f$ des chaînes de bits définit une observable diagonale dans la base de calcul :
\[
D_f=\sum_xf(x)\,\ket x\bra x=\mathrm{diag}\big(f(0),f(1),\dots,f(2^n-1)\big).
\]
La mesurer, c'est mesurer en base de calcul, lire la chaîne $x$ obtenue et afficher $f(x)$. Sa moyenne est
\[
\expval{D_f}_\psi=\sum_xf(x)\,|\braket x\psi|^2,
\]
c'est-à-dire \emph{la valeur moyenne de $f$ sur les chaînes tirées selon la règle de Born}. Pour notre projet, $f(x)$ est le $C_D$ prédit pour le profil $x$, et $D_f$ est l'hamiltonien de coût $H_C$ (ch.~\ref{ch:hamiltoniens}).

\paragraph{Minimiser une moyenne, c'est résoudre le problème.} Pour une observable hermitienne, $\bra\psi A\ket\psi$ est la forme quadratique de $A$, et sur les vecteurs normés c'est le quotient de Rayleigh. Son minimum sur tous les états vaut la plus petite valeur propre $a_{\min}$, atteinte au vecteur propre correspondant : c'est le \emph{principe variationnel}. La preuve tient en une ligne : $\bra\psi A\ket\psi=\sum_ka_k|\braket{u_k}\psi|^2\ge a_{\min}\sum_k|\braket{u_k}\psi|^2=a_{\min}$.

Pour $D_f$, les valeurs propres sont les $f(x)$ et les vecteurs propres sont les chaînes $\ket x$. Chercher l'état $\ket\psi$ qui minimise $\expval{H_C}$ revient donc exactement à chercher le profil $x^\star$ de $C_D$ minimal, et l'état optimal est $\ket{x^\star}$ lui-même. Toute l'idée de QAOA est là : remplacer une recherche discrète parmi $2^n$ chaînes par la minimisation continue d'une moyenne $\expval{H_C}$, sur une famille d'états paramétrés.

\paragraph{Matrices de Pauli.} Ce sont les briques de tout le cours :
\[
X=\begin{pmatrix}0&1\\1&0\end{pmatrix},\quad Y=\begin{pmatrix}0&-i\\i&0\end{pmatrix},\quad Z=\begin{pmatrix}1&0\\0&-1\end{pmatrix}.
\]
Elles sont hermitiennes, unitaires, de carré $\Id$, de valeurs propres $\pm1$.
\begin{center}
\begin{tabular}{lll}
\toprule
Observable & vecteur propre $+1$ & vecteur propre $-1$\\
\midrule
$Z$ & $\ket0$ & $\ket1$\\
$X$ & $\ket+$ & $\ket-$\\
$Y$ & $\ket{+i}$ & $\ket{-i}$\\
\bottomrule
\end{tabular}
\end{center}
Sur la sphère de Bloch, $\expval X,\expval Y,\expval Z$ sont exactement les coordonnées $(x,y,z)$ du point représentant l'état.

\begin{resultat}
$Z\ket b=(1-2b)\ket b$ pour $b\in\{0,1\}$ : la valeur propre de $Z$ est le \textbf{spin} $s=1-2b\in\{+1,-1\}$. Par conséquent, l'opérateur
\[\hat b=\frac{\Id-Z}{2}=\ket1\bra1\]
a pour valeur propre le bit $b$ lui-même. C'est le pont entre une variable binaire d'optimisation et un opérateur quantique (ch.~\ref{ch:qubo}).
\end{resultat}

\section{Évolution : opérateurs unitaires}\label{sec:unitaires}
\begin{definition}[ (postulat d'évolution)]
Un système isolé évolue par une matrice \textbf{unitaire} $U$ ($U^\dagger U=\Id$) : $\ket\psi\mapsto U\ket\psi$. L'unitarité garantit la conservation de la norme, donc la somme des probabilités égale à 1. Toute évolution est \emph{réversible} ($U^{-1}=U^\dagger$) ; seule la mesure ne l'est pas.
\end{definition}

\paragraph{Exponentielle de matrice et rotations.} Pour $P\in\{X,Y,Z\}$, $P^2=\Id$, donc la série de l'exponentielle se sépare en termes pairs et impairs :
\[
e^{-i\theta P}=\cos\theta\,\Id-i\sin\theta\,P .
\]
Les \textbf{rotations} $R_P(\theta)=e^{-i\theta P/2}$ font tourner la sphère de Bloch d'un angle $\theta$ autour de l'axe $P$ (le facteur $1/2$ vient de ce que l'état est paramétré par $\theta/2$). Par exemple, $R_Z(\theta)=\mathrm{diag}(e^{-i\theta/2},e^{i\theta/2})$ ne change que la phase relative : rotation autour de $z$, probabilités inchangées.

\section{Non-commutation}
$XZ=-ZX$ : les Pauli \emph{anticommutent}. Deux observables qui ne commutent pas n'ont pas de base propre commune, et on ne peut pas leur attribuer simultanément une valeur précise. Pour l'optimisation, la conséquence est directe :
\begin{itemize}
\item un hamiltonien de coût construit avec des $Z$ seulement est \emph{diagonal}. Son évolution $e^{-i\gamma H_C}$ ne fait que déphaser les états de la base de calcul, sans jamais changer les probabilités de mesure ;
\item pour déplacer de la probabilité d'un profil vers un autre, il faut un opérateur qui \emph{ne commute pas} avec $H_C$ : c'est le rôle du \textbf{mélangeur} $\sum_kX_k$ dans QAOA et du champ transverse dans le recuit quantique.
\end{itemize}

\exercices
\begin{exercice}
(a) Vérifier que $\ket\pm$ sont vecteurs propres de $X$, de valeurs propres $\pm1$. (b) Calculer $\bra+Z\ket+$ et $\bra+X\ket+$, et interpréter sur la sphère de Bloch.
\end{exercice}
\begin{exercice}
(a) Vérifier que $H$ est unitaire et hermitienne. (b) Montrer que $HZH=X$ et $HXH=Z$ ; en déduire que $H$ échange les bases $\{\ket0,\ket1\}$ et $\{\ket+,\ket-\}$. (c) Établir $e^{-i\theta X}=\cos\theta\,\Id-i\sin\theta\,X$.
\end{exercice}

\chapter{Plusieurs qubits}\label{ch:multi}
\iterbanner{V0}

\section{Produit tensoriel}
\paragraph{Le point de départ : une amplitude par configuration.} Deux bits classiques ont quatre configurations possibles : 00, 01, 10, 11. Pour un qubit, on attache une amplitude complexe à chacun des deux résultats possibles. Pour deux qubits, la règle est la même : on attache une amplitude à chacun des \emph{quatre} résultats possibles de la mesure des deux qubits.
\[
\ket\psi=\psi_{00}\ket{00}+\psi_{01}\ket{01}+\psi_{10}\ket{10}+\psi_{11}\ket{11},\qquad |\psi_{00}|^2+|\psi_{01}|^2+|\psi_{10}|^2+|\psi_{11}|^2=1 .
\]
Le ket $\ket{ab}$ désigne l'état où le qubit~0 vaut $a$ et le qubit~1 vaut $b$, et $|\psi_{ab}|^2$ est la probabilité de lire $a$ sur le qubit~0 \emph{et} $b$ sur le qubit~1. Avec $n$ qubits, il y a $2^n$ chaînes de bits, donc $2^n$ amplitudes :
\begin{center}
\begin{tabular}{lcccccc}
\toprule
qubits $n$ & 1 & 2 & 3 & 9 (notre QUBO) & 30 & 50\\
\midrule
amplitudes $2^n$ & 2 & 4 & 8 & 512 & $\approx10^9$ & $\approx10^{15}$\\
\bottomrule
\end{tabular}
\end{center}
L'état de $n$ qubits n'est donc \emph{pas} une liste de $n$ états à un qubit, ce qui ferait $2n$ nombres. C'est un vecteur de $\C^{2^n}$. Pour nos 9 qubits, c'est une amplitude par profil $(m,p,t)$ de la grille.

\paragraph{Le produit tensoriel : combiner deux qubits indépendants.} Supposons que le qubit~0 soit dans l'état $\alpha\ket0+\beta\ket1$ et le qubit~1, préparé séparément, dans $\gamma\ket0+\delta\ket1$. L'amplitude de chaque configuration jointe est le \emph{produit} des amplitudes individuelles, comme la probabilité de deux événements indépendants est le produit des probabilités. On note cette opération $\otimes$ :
\[
(\alpha\ket0+\beta\ket1)\otimes(\gamma\ket0+\delta\ket1)=\alpha\gamma\ket{00}+\alpha\delta\ket{01}+\beta\gamma\ket{10}+\beta\delta\ket{11}.
\]
En pratique, on développe comme un produit ordinaire, en gardant l'ordre des facteurs : $\ket a\otimes\ket b=\ket{ab}$, et on écrit souvent $\ket a\ket b$. La normalisation est automatique : $\sum|\cdot|^2=(|\alpha|^2+|\beta|^2)(|\gamma|^2+|\delta|^2)=1$.

\paragraph{Exemples.}
\begin{itemize}
\item $\ket0\otimes\ket1=\ket{01}$ : un état de base, avec une seule amplitude non nulle.
\item $\ket0\otimes\ket+=\ket0\otimes\tfrac1{\sqrt2}(\ket0+\ket1)=\tfrac1{\sqrt2}(\ket{00}+\ket{01})$. Le qubit~0 vaut 0 à coup sûr, et le qubit~1 est à pile ou face.
\item $\ket+\otimes\ket+=\tfrac12(\ket{00}+\ket{01}+\ket{10}+\ket{11})$ : les quatre configurations ont la même probabilité $\tfrac14=\tfrac12\cdot\tfrac12$.
\item $\ket+^{\otimes3}=\ket+\otimes\ket+\otimes\ket+=\tfrac1{2\sqrt2}\big(\ket{000}+\ket{001}+\dots+\ket{111}\big)$ : 8 chaînes, chacune de probabilité $\tfrac18$. De même, $\ket+^{\otimes9}$ donne la même amplitude $2^{-9/2}$ aux 512 profils : c'est l'état de départ de QAOA.
\end{itemize}

\paragraph{Une porte qui agit sur un seul qubit.} L'opérateur $A\otimes B$ applique $A$ au qubit~0 et $B$ au qubit~1 : $(A\otimes B)(\ket u\otimes\ket v)=A\ket u\otimes B\ket v$. Pour ne toucher qu'un qubit, on met l'identité $\Id$ (« ne rien faire ») sur l'autre :
\[
(X\otimes\Id)\ket{01}=X\ket0\otimes\ket1=\ket{11},\qquad (H\otimes\Id)\ket{00}=H\ket0\otimes\ket0=\tfrac1{\sqrt2}(\ket{00}+\ket{10}).
\]
Sur un état qui n'est pas un produit, on applique la règle à chaque terme, par linéarité.

\paragraph{Dans la mémoire de l'ordinateur.} Un simulateur range les $2^n$ amplitudes dans un tableau de nombres complexes. Le produit de Kronecker (\texttt{np.kron}) calcule ce tableau à partir des facteurs :
\[
\begin{pmatrix}\alpha\\\beta\end{pmatrix}\otimes\begin{pmatrix}\gamma\\\delta\end{pmatrix}=\begin{pmatrix}\alpha\gamma\\\alpha\delta\\\beta\gamma\\\beta\delta\end{pmatrix}
\quad\text{rangé dans l'ordre}\quad 00,\ 01,\ 10,\ 11 .
\]
Avec \texttt{np.kron(u, v)}, le \emph{premier} facteur donne le bit de poids fort de l'indice.

\begin{piege}[ : l'ordre de rangement est une convention]
L'état physique ne dépend d'aucune convention. Seul l'\emph{ordre} dans lequel on range les amplitudes dans le tableau en dépend. Dans ce cours et dans notre code, la chaîne $\ket{b_0b_1\cdots b_{n-1}}$ (qubit~0 écrit à gauche) est rangée à l'indice $x=\sum_kb_k2^k$ : le qubit~0 est le bit de \emph{poids faible}. Avec \texttt{np.kron}, il faut donc écrire les facteurs à l'envers, \texttt{kron(q[n-1], ..., q[0])}. C'est aussi la convention d'indice de Qiskit, qui \emph{affiche} en revanche les chaînes de droite à gauche ($q_{n-1}\cdots q_0$, comme on écrit le nombre $x$ en binaire). Il faut toujours vérifier l'ordre sur un exemple à deux qubits avant d'interpréter des résultats.
\end{piege}

\section{États produits et états intriqués}\label{sec:intrique}
\paragraph{Définition.} Un état est \textbf{produit} s'il s'écrit $\ket{\phi_0}\otimes\ket{\phi_1}\otimes\cdots$ : chaque qubit a alors son propre état, indépendant des autres. Sinon, il est \textbf{intriqué} : l'état du système ne se décompose pas en états individuels.

\paragraph{Le test à deux qubits.} Rangeons les quatre amplitudes dans un tableau $2\times2$, avec le qubit~0 en ligne et le qubit~1 en colonne. Pour un état produit, $\psi_{ab}=u_a\,v_b$ : chaque case est le produit d'un facteur de ligne par un facteur de colonne (figure~\ref{fig:tableau-amplitudes}). On en déduit
\[
\psi_{00}\,\psi_{11}=u_0v_0\,u_1v_1=u_0v_1\,u_1v_0=\psi_{01}\,\psi_{10}.
\]
Un état produit vérifie donc $\psi_{00}\psi_{11}-\psi_{01}\psi_{10}=0$ : le déterminant du tableau est nul. Réciproquement, un tableau $2\times2$ de déterminant nul est de rang 1, donc s'écrit $u_av_b$. Ce critère est une équivalence.

\begin{figure}[htbp]
\centering
\begin{tikzpicture}[cell/.style={draw=cink,minimum width=1.25cm,minimum height=0.8cm,font=\small}]
\foreach \titre/\aa/\ab/\ba/\bb/\fa/\fb/\fc/\fd [count=\i] in
  {{$\ket+\ket+$ (produit)}/{$\tfrac12$}/{$\tfrac12$}/{$\tfrac12$}/{$\tfrac12$}/25/25/25/25,
   {$\ket0\ket+$ (produit)}/{$\tfrac1{\sqrt2}$}/{$\tfrac1{\sqrt2}$}/{$0$}/{$0$}/40/40/0/0,
   {$\ket{\Phi^+}$ (intriqué)}/{$\tfrac1{\sqrt2}$}/{$0$}/{$0$}/{$\tfrac1{\sqrt2}$}/40/0/0/40}{
  \begin{scope}[xshift={(\i-1)*4.6cm}]
    \node[cell,fill=cblue!\fa] at (0,0) {\aa};
    \node[cell,fill=cblue!\fb] at (1.25,0) {\ab};
    \node[cell,fill=cblue!\fc] at (0,-0.8) {\ba};
    \node[cell,fill=cblue!\fd] at (1.25,-0.8) {\bb};
    \node[font=\scriptsize,cmuted] at (0,0.62) {$q_1=0$};
    \node[font=\scriptsize,cmuted] at (1.25,0.62) {$q_1=1$};
    \node[font=\scriptsize,cmuted,anchor=east] at (-0.68,0) {$q_0=0$};
    \node[font=\scriptsize,cmuted,anchor=east] at (-0.68,-0.8) {$q_0=1$};
    \node[font=\small] at (0.62,-1.65) {\titre};
  \end{scope}}
\end{tikzpicture}
\caption{Les amplitudes de deux qubits rangées en tableau (ligne : qubit~0, colonne : qubit~1). Pour un état produit, les lignes sont proportionnelles (rang 1, déterminant nul). Pour $\ket{\Phi^+}$, le déterminant vaut $\tfrac12$ : impossible de factoriser.}\label{fig:tableau-amplitudes}
\end{figure}

\paragraph{Exemples.}
\begin{itemize}
\item $\tfrac12(\ket{00}+\ket{01}+\ket{10}+\ket{11})$ : $\tfrac12\cdot\tfrac12-\tfrac12\cdot\tfrac12=0$, donc produit. C'est $\ket+\ket+$.
\item $\tfrac1{\sqrt2}(\ket{00}+\ket{01})$ : $\tfrac1{\sqrt2}\cdot0-\tfrac1{\sqrt2}\cdot0=0$, donc produit. C'est $\ket0\ket+$.
\item L'\textbf{état de Bell} $\ket{\Phi^+}=\tfrac1{\sqrt2}(\ket{00}+\ket{11})$ : $\tfrac1{\sqrt2}\cdot\tfrac1{\sqrt2}-0\cdot0=\tfrac12\neq0$, donc intriqué. On peut le voir directement. Si l'on avait $\psi_{ab}=u_av_b$, alors $\psi_{00}\neq0$ imposerait $u_0\neq0$, et $\psi_{11}\neq0$ imposerait $v_1\neq0$. On aurait donc $\psi_{01}=u_0v_1\neq0$, alors que $\psi_{01}=0$. Contradiction.
\end{itemize}

\paragraph{Ce que l'intrication change aux mesures.} Comparons les probabilités jointes $P(ab)=|\psi_{ab}|^2$ :
\begin{itemize}
\item pour $\ket+\ket+$, $P(ab)=\tfrac14$ partout. Lire le qubit~0 n'apprend \emph{rien} sur le qubit~1. C'est vrai de tout état produit, pour lequel $P(ab)=|u_a|^2|v_b|^2$ : les deux résultats sont des variables aléatoires indépendantes ;
\item pour $\ket{\Phi^+}$, $P(00)=P(11)=\tfrac12$ et $P(01)=P(10)=0$. Chaque qubit, regardé seul, est à pile ou face, mais les deux résultats sont \emph{toujours égaux}. Lire le qubit~0 donne le qubit~1 à coup sûr.
\end{itemize}

\paragraph{Corrélation quantique ou simple corrélation classique ?} Deux pièces truquées, collées pour tomber toujours du même côté, donnent aussi « 00 ou 11, à 50/50 ». En base de calcul, $\ket{\Phi^+}$ ressemble donc à ce \emph{mélange classique}. La différence apparaît dans une autre base (§\ref{sec:autre-base}). Un calcul direct montre que
\[
\ket{\Phi^+}=\tfrac1{\sqrt2}\big(\ket{++}+\ket{--}\big).
\]
Mesurés tous deux en base $\{\ket+,\ket-\}$, les deux qubits donnent encore toujours le même résultat. Pour le mélange classique, en revanche, chaque qubit est un $\ket0$ ou un $\ket1$, qui donne 50/50 en base $\{\ket\pm\}$, indépendamment de l'autre : la corrélation disparaît. L'intrication est une corrélation qui tient \emph{dans toutes les bases à la fois}, ce qu'aucun mélange de bits classiques ne permet.

\paragraph{D'où vient l'intrication ?} Une porte à un qubit transforme un état produit en un autre état produit, puisque $(A\otimes B)(\ket u\otimes\ket v)=A\ket u\otimes B\ket v$. Pour intriquer, il faut une porte qui agit sur deux qubits \emph{à la fois} et dont l'effet sur l'un dépend de l'autre. C'est le cas de la CNOT et de $R_{ZZ}$ (ch.~\ref{ch:circuits}). Dans QAOA, ce sont précisément les termes de couplage $Q_{ij}b_ib_j$ du coût qui produisent ces portes.

\paragraph{Pourquoi c'est utile pour l'optimisation.} Prenons un QUBO à deux bits, $E(b_0,b_1)=b_0+b_1-3\,b_0b_1$, soit $E(00)=0$, $E(01)=E(10)=1$ et $E(11)=-1$. Le meilleur choix pour $b_0$ dépend de $b_1$ : si $b_1=0$, il faut $b_0=0$ ; si $b_1=1$, il faut $b_0=1$. Les bonnes chaînes sont celles où $b_0=b_1$.

Un état qui « a compris » cette structure met beaucoup de probabilité sur 00 et 11, et peu sur 01 et 10, par exemple $P(00)=P(11)=0{,}45$ et $P(01)=P(10)=0{,}05$. Aucun état produit ne peut faire cela. Pour un état produit, $P(00)\,P(11)=P(01)\,P(10)$, ce qui donnerait ici $0{,}2025=0{,}0025$, impossible. Il faut un état intriqué pour représenter la règle « $b_0=b_1$ » sans encore savoir quelle valeur choisir.

Au fil des couches de QAOA, l'état passe de $\ket+^{\otimes n}$, qui est produit et uniforme, à des états intriqués qui concentrent la probabilité sur les chaînes compatibles avec les couplages. L'optimum lui-même, $\ket{x^\star}=\ket{11}$ ici, est un état de base, donc produit : l'intrication sert pendant le chemin, pas à l'arrivée.

\section{Opérateurs sur $n$ qubits}
\paragraph{Agir sur un seul qubit.} Un opérateur qui agit sur le seul qubit $k$ s'écrit $Z_k=\Id\otimes\cdots\otimes Z\otimes\cdots\otimes\Id$, avec $Z$ sur le facteur du qubit $k$ et $\Id$ sur tous les autres. Sur une chaîne de bits, il suffit de regarder ce que fait chaque facteur :
\[
Z_k\ket{b_0\cdots b_k\cdots b_{n-1}}=\ket{b_0}\otimes\cdots\otimes Z\ket{b_k}\otimes\cdots\otimes\ket{b_{n-1}}=s_k\,\ket{b_0\cdots b_{n-1}},\qquad s_k=1-2b_k .
\]
$Z_k$ \textbf{lit} le bit $k$ et multiplie la chaîne par $+1$ ou $-1$ selon sa valeur, sans la modifier. Toutes les chaînes $\ket x$ sont donc vecteurs propres de $Z_k$ : sa matrice est \emph{diagonale} dans la base de calcul. Exemple à deux qubits :
\[
Z_0\ket{00}=+\ket{00},\quad Z_0\ket{10}=-\ket{10},\quad Z_0\ket{01}=+\ket{01},\quad Z_0\ket{11}=-\ket{11}.
\]
À l'inverse, $X_k$ \textbf{retourne} le bit $k$ : $X_0\ket{00}=\ket{10}$ et $X_0\ket{11}=\ket{01}$. Il fait passer d'une chaîne à une autre, donc sa matrice n'est pas diagonale. Toute la suite repose sur cette distinction :
\begin{center}
\begin{tabular}{llll}
\toprule
Opérateur & Effet sur $\ket x$ & Matrice & Rôle dans QAOA \\
\midrule
$Z_k$, $Z_iZ_j$, $\hat b_k$ & multiplie par un nombre & diagonale & coût $H_C$ \\
$X_k$ & change de chaîne (retourne un bit) & non diagonale & mélangeur $H_M$ \\
\bottomrule
\end{tabular}
\end{center}

\paragraph{Produits d'opérateurs diagonaux.} Les opérateurs diagonaux commutent, et un produit d'opérateurs diagonaux multiplie leurs valeurs propres :
\[
Z_iZ_j\ket x=s_is_j\ket x,\qquad \hat b_k=\frac{\Id-Z_k}2\ \Rightarrow\ \hat b_k\ket x=\frac{1-s_k}{2}\ket x=b_k\ket x,\qquad \hat b_i\hat b_j\ket x=b_ib_j\ket x .
\]
L'opérateur $\hat b_k$ « affiche » le bit $k$, et $\hat b_i\hat b_j$ affiche le produit $b_ib_j$, qui vaut 1 seulement si les deux bits valent 1. Par linéarité, n'importe quel polynôme en bits devient un opérateur qui affiche la valeur de ce polynôme.

\paragraph{Exemple complet : le QUBO à deux bits du §\ref{sec:intrique}.} Pour $E(b_0,b_1)=b_0+b_1-3\,b_0b_1$, on remplace chaque bit par son opérateur :
\[
H_C=\hat b_0+\hat b_1-3\,\hat b_0\hat b_1 .
\]
Sur chaque chaîne, $H_C\ket{b_0b_1}=(b_0+b_1-3b_0b_1)\ket{b_0b_1}=E(b_0,b_1)\ket{b_0b_1}$. Dans l'ordre $x=0,1,2,3$, c'est-à-dire $\ket{00},\ket{10},\ket{01},\ket{11}$ :
\[
H_C=\mathrm{diag}\big(E(00),E(10),E(01),E(11)\big)=\mathrm{diag}(0,\ 1,\ 1,\ -1).
\]
La fonction de coût est simplement \emph{écrite sur la diagonale}. En remplaçant $\hat b_k=(\Id-Z_k)/2$ et en développant, on obtient la même matrice exprimée avec des $Z$ :
\[
H_C=\tfrac14\,\Id+\tfrac14\,(Z_0+Z_1)-\tfrac34\,Z_0Z_1 .
\]
Sur $\ket{11}$, par exemple, $s_0=s_1=-1$, et on retrouve bien $\tfrac14-\tfrac12-\tfrac34=-1$. C'est la \emph{forme d'Ising} du coût, que le chapitre~\ref{ch:qubo} construit en général. Elle est utile parce que chaque terme correspond à une porte : $Z_k$ à une rotation $R_Z$, $Z_iZ_j$ à une $R_{ZZ}$ (ch.~\ref{ch:circuits}).

\begin{resultat}
Toute fonction de coût polynomiale en bits $E(b_0,\dots,b_{n-1})$ devient un opérateur \textbf{diagonal} dans la base de calcul en remplaçant chaque $b_k$ par $\hat b_k=(\Id-Z_k)/2$ :
\[H_C=\sum_xE(x)\ket x\bra x,\qquad H_C\ket x=E(x)\ket x.\]
Sa plus petite valeur propre est $\min_xE(x)$, et le vecteur propre associé est le ket de la meilleure chaîne de bits.
\end{resultat}

\paragraph{Ce que fait $H_C$ à une superposition.} Par linéarité, sur $\ket\psi=\sum_x\psi_x\ket x$ :
\begin{itemize}
\item $H_C\ket\psi=\sum_xE(x)\,\psi_x\ket x$ : chaque amplitude est multipliée par le coût de sa chaîne ;
\item $\expval{H_C}=\sum_xE(x)\,|\psi_x|^2$ : le coût moyen des chaînes tirées à la mesure (§\ref{sec:observables}) ;
\item $e^{-i\gamma H_C}\ket\psi=\sum_xe^{-i\gamma E(x)}\psi_x\ket x$ : chaque amplitude reçoit une phase proportionnelle à son coût, sans que son module change. C'est la première étape d'une couche de QAOA, et la raison pour laquelle on a tant insisté sur les phases au §\ref{sec:phases}.
\end{itemize}

\paragraph{On n'écrit jamais la matrice $2^n\times2^n$.} Pour 9 qubits, $H_C$ est une matrice $512\times512$, mais seule sa diagonale compte : un vecteur de 512 coûts, que notre code calcule à partir de la matrice $Q$. Sur un ordinateur quantique, on n'a même pas besoin de cette liste. On applique une porte par terme du polynôme : 9 termes linéaires et $\binom92=36$ couplages, soit 45 portes. Le nombre de termes croît comme $n^2$, alors que le nombre de chaînes croît comme $2^n$ : c'est ce qui rend le codage d'un QUBO sur machine quantique efficace.

\section{Le coût exponentiel de la simulation}
\paragraph{Mémoire.} Un simulateur par vecteur d'état stocke les $2^n$ amplitudes complexes, à 16~octets chacune en double précision (deux réels de 8~octets) :
\begin{center}
\begin{tabular}{rrrl}
\toprule
$n$ & amplitudes $2^n$ & mémoire $16\cdot2^n$ & machine \\
\midrule
9 & 512 & 8~ko & tout (notre QUBO) \\
20 & $\approx10^6$ & 16~Mo & tout \\
30 & $\approx10^9$ & 16~Go & portable puissant \\
40 & $\approx10^{12}$ & 16~To & grand cluster \\
50 & $\approx10^{15}$ & 16~Po & hors de portée \\
\bottomrule
\end{tabular}
\end{center}
Chaque qubit ajouté \emph{double} la mémoire (figure~\ref{fig:memoire}).

\begin{figure}[htbp]
\centering
\begin{tikzpicture}
\begin{semilogyaxis}[width=0.8\linewidth,height=5.6cm,axis lines=left,xmin=0,xmax=52,ymin=10,ymax=1e17,
  xlabel={nombre de qubits $n$},ylabel={mémoire du vecteur d'état (octets)},
  ytick={1e3,1e6,1e9,1e12,1e15},yticklabels={1~ko,1~Mo,1~Go,1~To,1~Po}]
\addplot[cblue,very thick,domain=0:52,samples=53] {16*2^x};
\addplot[only marks,mark=*,corange,mark size=2.4pt] coordinates {(9,8192) (30,1.718e10) (45,5.6e14)};
\node[anchor=west,font=\small] at (axis cs:10,8192) {$n=9$ : 8~ko (notre QUBO)};
\node[anchor=north west,font=\small] at (axis cs:30.5,1.7e10) {$n=30$ : 16~Go (portable)};
\node[anchor=south east,font=\small] at (axis cs:44,5.6e14) {$n\approx45$ : supercalculateur};
\end{semilogyaxis}
\end{tikzpicture}
\caption{Mémoire d'une simulation exacte : $16\cdot2^n$ octets. Chaque qubit ajouté double le coût.}\label{fig:memoire}
\end{figure}

\paragraph{Temps de calcul.} Appliquer une porte, même sur un seul qubit, modifie en général toutes les amplitudes. Par exemple, $X_0$ échange les amplitudes de $\ket{0\cdots}$ et $\ket{1\cdots}$ deux à deux. Chaque porte coûte donc environ $2^n$ opérations. Une couche de QAOA sur notre QUBO, soit 45 termes de coût et 9 portes de mélange, revient à quelques dizaines de passes sur 512 nombres : quelques microsecondes. À 40 qubits, chaque passe porte sur $10^{12}$ nombres.

\paragraph{Pourquoi on ne peut pas compresser.} Un état produit n'a besoin que de $2n$ amplitudes, deux par qubit. C'est l'intrication (§\ref{sec:intrique}) qui oblige à garder les $2^n$ nombres, car ils ne se factorisent plus. Des méthodes approchées, comme les réseaux de tenseurs, exploitent les états \emph{peu} intriqués pour aller plus loin. Mais pour un état fortement intriqué, aucune méthode classique connue n'évite le coût exponentiel.

\paragraph{Ce que change une vraie machine, et ce qu'elle ne change pas.} Sur un processeur quantique, les $2^n$ amplitudes ne sont stockées nulle part : elles \emph{sont} l'état physique des $n$ qubits, et une porte prend le même temps quel que soit $n$. Mais on ne peut pas lire ces amplitudes. Chaque shot ne renvoie que $n$ bits, une seule chaîne tirée selon la règle de Born. Le vecteur d'état exponentiel n'est donc pas une mémoire exponentielle consultable. L'avantage n'existe que si le circuit, par interférence, concentre la probabilité sur les bonnes chaînes, pour qu'un petit nombre de shots suffise à les trouver.

\begin{projet}
Nos 9 bits donnent $2^9=512$ amplitudes : \emph{exactement} la taille de la grille de profils. Simuler QAOA coûte donc au moins autant que d'énumérer le QUBO, et c'est voulu, puisque le cas doit rester vérifiable : on peut comparer la sortie de QAOA au vrai optimum. Deux ordres de grandeur à garder en tête pour les extensions :
\begin{itemize}
\item 3 variables à 16 niveaux (4 bits chacune) : 12 qubits, 4\,096 amplitudes. C'est encore trivial à simuler ;
\item 20 variables de forme à 4 bits : 80 qubits. Ni l'énumération ($2^{80}\approx10^{24}$ profils) ni la simulation ne sont possibles. C'est le régime où un ordinateur quantique pourrait servir.
\end{itemize}
L'intérêt d'une machine réelle n'apparaîtrait que pour $n\gtrsim50$. C'est pourquoi l'analyse des \emph{ressources} en fonction de $n$ (qubits, portes, shots) fait partie du livrable.
\end{projet}

\exercices
\begin{exercice}
Dire si chaque état est produit ou intriqué, et le factoriser le cas échéant : (a) $\tfrac12(\ket{00}-\ket{01}+\ket{10}-\ket{11})$ ; (b) $\tfrac1{\sqrt2}(\ket{01}-\ket{10})$.
\end{exercice}
\begin{exercice}
Pour $n=2$, écrire les matrices diagonales de $Z_0$, $Z_1$, $Z_0Z_1$ et $\hat b_0\hat b_1=\tfrac14(\Id-Z_0)(\Id-Z_1)$ dans la base $(\ket{00},\ket{10},\ket{01},\ket{11})$, c'est-à-dire dans l'ordre $x=0,1,2,3$ avec $x=b_0+2b_1$. Vérifier que $\hat b_0\hat b_1$ vaut 1 sur $\ket{11}$ seulement.
\end{exercice}

\chapter{Portes et circuits quantiques}\label{ch:circuits}
\iterbanner{V0}

Un \textbf{circuit quantique} est la façon standard d'écrire une suite d'opérations unitaires (des \emph{portes}) suivie de mesures. C'est l'équivalent d'un programme : QAOA est un circuit.

\section{Lire un circuit}
\begin{itemize}
\item Chaque \emph{ligne horizontale} est un qubit ; le temps s'écoule de gauche à droite.
\item Chaque \emph{boîte} est une porte, c'est-à-dire une matrice unitaire appliquée au(x) qubit(s) qu'elle touche.
\item Le symbole \emph{cadran} est une mesure en base de calcul ; il en sort un bit classique.
\item Par convention, tous les qubits commencent dans $\ket0$.
\end{itemize}
Une suite de portes $U_1$ puis $U_2$ correspond au produit matriciel $U_2U_1$ : l'ordre est inversé, car la première porte est la plus proche du vecteur.

\begin{figure}[htbp]
\centering
\begin{tikzpicture}[x=1.15cm,y=0.9cm]
\draw[cink] (0,0) -- (3.2,0);
\node[anchor=east,font=\small] at (0,0) {$\ket0$};
\node[gate] at (1,0) {$H$};
\meter{2.2}{0}
\draw[cink,double] (2.62,0) -- (3.3,0) node[right,font=\small] {bit : 0 ou 1, prob.\ $\tfrac12$};
\end{tikzpicture}
\caption{Le plus petit circuit utile : un générateur de bit aléatoire parfait. $H\ket0=\ket+$, puis la mesure donne 0 ou 1 avec probabilité $1/2$.}
\end{figure}

\section{Portes à un qubit}
\begin{center}
\renewcommand{\arraystretch}{1.25}
\begin{tabular}{llp{7.3cm}}
\toprule
Porte & Matrice & Effet\\
\midrule
$X$ (NOT) & $\begin{pmatrix}0&1\\1&0\end{pmatrix}$ & échange $\ket0\leftrightarrow\ket1$ ; demi-tour autour de $x$\\
$Z$ & $\begin{pmatrix}1&0\\0&-1\end{pmatrix}$ & change le signe de $\ket1$ ; demi-tour autour de $z$\\
$H$ (Hadamard) & $\tfrac1{\sqrt2}\begin{pmatrix}1&1\\1&-1\end{pmatrix}$ & $\ket0\mapsto\ket+$, $\ket1\mapsto\ket-$ ; crée la superposition\\
$R_X(\theta)$ & $e^{-i\theta X/2}$ & rotation d'angle $\theta$ autour de $x$ \textbf{(mélangeur de QAOA)}\\
$R_Z(\theta)$ & $\mathrm{diag}(e^{-i\theta/2},e^{i\theta/2})$ & rotation autour de $z$, déphasage \textbf{(termes linéaires du coût)}\\
\bottomrule
\end{tabular}
\end{center}
Remarque utile : $R_Z(\theta)=e^{-i\theta/2}\,\mathrm{diag}(1,e^{i\theta})$. Les conventions diffèrent d'une phase globale selon les bibliothèques, ce qui est sans conséquence physique (ch.~\ref{ch:qubit}).

\subsection*{Pourquoi les portes à un qubit sont des rotations}
\paragraph{L'argument géométrique.} Une porte $U$ est unitaire, donc elle conserve les produits scalaires : $|\braket{U\phi}{U\psi}|=|\braket\phi\psi|$. Or on a vu au §\ref{sec:autre-base} que $|\braket\phi\psi|^2=\tfrac12(1+\vec r_\phi\cdot\vec r_\psi)$, où $\vec r_\phi$ et $\vec r_\psi$ sont les points de Bloch des deux états. Une porte conserve donc le produit scalaire entre points de la sphère, c'est-à-dire les angles et les distances : c'est une \emph{isométrie} de la sphère, qui fixe son centre. Les seules isométries qu'on peut atteindre continûment depuis « ne rien faire », en tournant progressivement un bouton physique (durée d'une impulsion, champ magnétique), sont les \textbf{rotations}.

\paragraph{Vérification directe sur $R_Z$.} Partons d'un état quelconque, écrit avec ses angles de Bloch : $\ket\psi=\cos\tfrac\theta2\ket0+e^{i\varphi}\sin\tfrac\theta2\ket1$. Alors
\[
R_Z(\alpha)\ket\psi=e^{-i\alpha/2}\cos\tfrac\theta2\ket0+e^{i\alpha/2}e^{i\varphi}\sin\tfrac\theta2\ket1
=e^{-i\alpha/2}\Big(\cos\tfrac\theta2\ket0+e^{i(\varphi+\alpha)}\sin\tfrac\theta2\ket1\Big).
\]
Le facteur $e^{-i\alpha/2}$ est une phase globale, sans effet. Il reste $\theta\to\theta$ et $\varphi\to\varphi+\alpha$ : la latitude ne bouge pas, la longitude augmente de $\alpha$. C'est exactement une rotation d'angle $\alpha$ autour de l'axe $z$ (figure~\ref{fig:rotations}a). Les probabilités en base de calcul ne changent pas ; seule la phase relative tourne.

\paragraph{Vérification directe sur $R_X$.} Comme $X^2=\Id$, on a (exercice~4.2) $R_X(\theta)=e^{-i\theta X/2}=\cos\tfrac\theta2\,\Id-i\sin\tfrac\theta2\,X$, soit
\[
R_X(\theta)=\begin{pmatrix}\cos\frac\theta2&-i\sin\frac\theta2\\[2pt]-i\sin\frac\theta2&\cos\frac\theta2\end{pmatrix},\qquad
R_X(\theta)\ket0=\cos\tfrac\theta2\ket0+e^{-i\pi/2}\sin\tfrac\theta2\ket1 .
\]
On lit les angles de Bloch du résultat : latitude $\theta$, longitude $\varphi=-\pi/2$. Le pôle nord $\ket0$, de coordonnées $(0,0,1)$, arrive en $(0,-\sin\theta,\cos\theta)$ : il a tourné d'un angle $\theta$ dans le plan $(y,z)$, autour de l'axe $x$ (figure~\ref{fig:rotations}b). Contrairement à $R_Z$, cette rotation \emph{change les probabilités} : $P(1)$ passe de 0 à $\sin^2\tfrac\theta2$. Par exemple, $R_X(\pi/2)\ket0=\tfrac1{\sqrt2}(\ket0-i\ket1)=\ket{-i}$ : un quart de tour amène le pôle nord sur l'équateur.

\begin{figure}[htbp]
\centering
\begin{tikzpicture}[>={Stealth[length=2mm]}]
\foreach \pan/\xs in {a/0,b/8}{
\begin{scope}[xshift=\xs cm,x={(-0.92cm,-0.39cm)},y={(1.31cm,-0.275cm)},z={(0cm,1.53cm)}]
  \draw[cink,thin] (0,0,0) circle (1.6cm);
  \draw[cmuted,thin] plot[domain=0:360,samples=90,smooth] ({cos(\x)},{sin(\x)},0);
  \draw[->,cmuted] (0,0,0) -- (1.35,0,0) node[anchor=north east,cink,font=\small] {$x$};
  \draw[->,cmuted] (0,0,0) -- (0,1.35,0) node[anchor=west,cink,font=\small] {$y$};
  \draw[->,cmuted] (0,0,0) -- (0,0,1.3) node[anchor=south west,cink,font=\small] {$z$};
  \fill[cink] (0,0,1) circle (1.3pt) node[anchor=south east,font=\small] {$\ket0$};
  \fill[cink] (0,0,-1) circle (1.3pt) node[anchor=north,font=\small] {$\ket1$};
\end{scope}}
% (a) R_Z : |+> tourne sur l'équateur
\begin{scope}[x={(-0.92cm,-0.39cm)},y={(1.31cm,-0.275cm)},z={(0cm,1.53cm)}]
  \fill[cblue] (1,0,0) circle (1.8pt) node[anchor=north,font=\small,cblue,yshift=-2pt] {$\ket+$};
  \draw[corange,very thick,->] plot[domain=0:70,samples=30,smooth] ({cos(\x)},{sin(\x)},0);
  \fill[cblue] ({cos(70)},{sin(70)},0) circle (1.8pt);
  \draw[cblue,thick,->] (0,0,0) -- ({cos(70)},{sin(70)},0);
  \draw[corange,thick,->] plot[domain=20:330,samples=40,smooth] ({0.18*cos(\x)},{0.18*sin(\x)},0.7);
  \node[corange,font=\small] at (0.75,0.95,0) {$\alpha$};
\end{scope}
\node[font=\small,align=center,text width=7cm] at (0,-2.6) {(a) $R_Z(\alpha)$ : rotation autour de $z$.\\ La longitude $\varphi$ tourne, $P(0)$ et $P(1)$ restent fixes.};
% (b) R_X : |0> tourne dans le plan (y,z)
\begin{scope}[xshift=8cm,x={(-0.92cm,-0.39cm)},y={(1.31cm,-0.275cm)},z={(0cm,1.53cm)}]
  \draw[corange,very thick,->] plot[domain=0:70,samples=30,smooth] (0,{-sin(\x)},{cos(\x)});
  \fill[cblue] (0,{-sin(70)},{cos(70)}) circle (1.8pt);
  \draw[cblue,thick,->] (0,0,0) -- (0,{-sin(70)},{cos(70)});
  \draw[cmuted,dashed] (0,0,0) -- (0,-1,0) node[anchor=east,cink,font=\small] {$\ket{-i}$};
  \draw[corange,thick,->] plot[domain=20:330,samples=40,smooth] (1.2,{0.18*cos(\x)},{0.18*sin(\x)});
  \node[corange,font=\small] at (0,-0.75,0.95) {$\theta$};
\end{scope}
\node[font=\small,align=center,text width=7cm] at (8,-2.6) {(b) $R_X(\theta)$ : rotation autour de $x$.\\ Le pôle $\ket0$ descend vers l'équateur : $P(1)=\sin^2\frac\theta2$.};
\end{tikzpicture}
\caption{Les deux rotations de QAOA sur la sphère de Bloch, ici pour $\alpha=\theta=70^\circ$. $R_Z$ écrit une phase sans changer les probabilités (étape de coût). $R_X$ fait basculer l'état entre les pôles, ce qui transforme une phase en probabilité (étape de mélange).}\label{fig:rotations}
\end{figure}

\paragraph{Toutes les portes à un qubit.} Plus généralement, pour un axe unitaire $\vec n=(n_x,n_y,n_z)$, la porte $R_{\vec n}(\theta)=e^{-i\theta(n_xX+n_yY+n_zZ)/2}$ est la rotation d'angle $\theta$ autour de $\vec n$, et \emph{toute} porte à un qubit est une telle rotation, à une phase globale près. Les portes du tableau en sont des cas particuliers :
\begin{itemize}
\item $X=i\,R_X(\pi)$ et $Z=i\,R_Z(\pi)$ : des demi-tours autour de $x$ et de $z$ ;
\item $H=\tfrac1{\sqrt2}(X+Z)=i\,R_{\vec n}(\pi)$ avec $\vec n=\tfrac1{\sqrt2}(1,0,1)$ : un demi-tour autour de la diagonale entre $x$ et $z$. Il échange les axes $x$ et $z$, donc $\ket0\leftrightarrow\ket+$ et $\ket1\leftrightarrow\ket-$.
\end{itemize}

\paragraph{D'où vient le facteur $\tfrac12$ ?} Dans $R_Z(\alpha)=e^{-i\alpha Z/2}$, l'angle sur la sphère est $\alpha$, alors que les amplitudes ne reçoivent que $\pm\alpha/2$. C'est le même « angle moitié » que dans $\cos\tfrac\theta2$ (§3.4) : deux états orthogonaux ne font que $90^\circ$ dans l'espace des vecteurs, mais $180^\circ$ sur la sphère. Conséquence surprenante : $R_Z(2\pi)=-\Id$. Un tour complet ramène le point de Bloch à sa place, mais multiplie l'état par $-1$, une phase globale invisible.

\section{Portes à deux qubits}
\paragraph{Pourquoi une porte à deux qubits ?} Avec des portes à un qubit seulement, un état produit reste produit (§\ref{sec:intrique}) : les qubits évoluent chacun de leur côté et ne s'influencent jamais. Pour que la valeur d'un qubit agisse sur un autre, comme $b_ib_j$ dans un QUBO, il faut au moins une porte qui agisse sur deux qubits à la fois. La plus simple est la CNOT.

\paragraph{CNOT (NON contrôlé) : définition.} Elle a un qubit \emph{contrôle} $c$ et un qubit \emph{cible} $t$. Elle applique $X$ (NOT) à la cible si et seulement si le contrôle vaut 1 :
\[
\mathrm{CNOT}\ket{c,t}=\ket{c,\ t\oplus c},
\]
où $\oplus$ est l'addition modulo 2 ($1\oplus1=0$). Le contrôle n'est jamais modifié. Sur les quatre états de base :
\begin{center}
\begin{tabular}{ccc}
\toprule
entrée $\ket{c,t}$ & action sur la cible & sortie $\ket{c,t\oplus c}$\\
\midrule
$\ket{00}$ & contrôle à 0 : rien & $\ket{00}$\\
$\ket{01}$ & contrôle à 0 : rien & $\ket{01}$\\
$\ket{10}$ & contrôle à 1 : NOT & $\ket{11}$\\
$\ket{11}$ & contrôle à 1 : NOT & $\ket{10}$\\
\bottomrule
\end{tabular}
\end{center}
Ce tableau définit entièrement la porte, puisqu'une porte est linéaire. Sur une superposition, on l'applique terme à terme.

\paragraph{Sa matrice.} Dans l'ordre des étiquettes $\ket{c,t}=\ket{00},\ket{01},\ket{10},\ket{11}$, la matrice est une permutation : elle échange les deux dernières lignes. Elle se lit aussi par blocs :
\[
\mathrm{CNOT}=\begin{pmatrix}1&0&0&0\\0&1&0&0\\0&0&0&1\\0&0&1&0\end{pmatrix}=\begin{pmatrix}\Id&0\\0&X\end{pmatrix}
=\ket0\bra0\otimes\Id+\ket1\bra1\otimes X .
\]
Autrement dit, « si le contrôle est dans $\ket0$, ne rien faire ; s'il est dans $\ket1$, appliquer $X$ à la cible ». C'est une matrice de permutation réelle et symétrique, donc unitaire, et elle est sa propre inverse : $\mathrm{CNOT}^2=\Id$. Appliquer deux CNOT de suite ne fait rien, ce qui servira pour $R_{ZZ}$.

\paragraph{Deux usages classiques.} Sur des états de base, la CNOT est une porte logique ordinaire :
\begin{itemize}
\item elle calcule le \textbf{XOR} : la cible reçoit $t\oplus c$, qui vaut 1 si les deux bits diffèrent. C'est la \emph{parité} des deux bits ;
\item si la cible part de 0, elle \textbf{copie} le contrôle : $\ket{c,0}\mapsto\ket{c,c}$.
\end{itemize}

\paragraph{Sur une superposition : intrication, pas copie.} Prenons un contrôle dans un état quelconque $\alpha\ket0+\beta\ket1$ et une cible dans $\ket0$. Par linéarité, chaque terme est traité séparément :
\[
\mathrm{CNOT}\big[(\alpha\ket0+\beta\ket1)\otimes\ket0\big]=\alpha\,\mathrm{CNOT}\ket{00}+\beta\,\mathrm{CNOT}\ket{10}=\alpha\ket{00}+\beta\ket{11}.
\]
Ce n'est \emph{pas} une copie de l'état. Une copie serait $(\alpha\ket0+\beta\ket1)\otimes(\alpha\ket0+\beta\ket1)$, qui contient des termes $\alpha\beta\ket{01}$ et $\alpha\beta\ket{10}$. Le résultat est intriqué dès que $\alpha$ et $\beta$ sont tous deux non nuls : son déterminant $\psi_{00}\psi_{11}-\psi_{01}\psi_{10}=\alpha\beta$ n'est pas nul. La CNOT a copié la \emph{valeur de base} de chaque branche, pas l'état. C'est une illustration du théorème de non-clonage : aucune porte ne peut copier un état quantique inconnu. Avec $\alpha=\beta=\tfrac1{\sqrt2}$, c'est-à-dire un contrôle préparé par $H$, on obtient l'état de Bell (figure~\ref{fig:bell}).

\begin{figure}[htbp]
\centering
\begin{tikzpicture}[x=1.15cm,y=0.9cm]
\draw[cink] (0,0) -- (3.4,0); \draw[cink] (0,-1) -- (3.4,-1);
\node[anchor=east,font=\small] at (0,0) {$q_0:\ \ket0$};
\node[anchor=east,font=\small] at (0,-1) {$q_1:\ \ket0$};
\node[gate] at (0.9,0) {$H$};
\draw[cink] (1.9,0) -- (1.9,-1); \node[ctrl] at (1.9,0) {}; \targ{1.9}{-1}
\meter{2.9}{0} \meter{2.9}{-1}
\node[font=\small,anchor=west] at (3.6,-0.5) {$\to\ \tfrac1{\sqrt2}(\ket{00}+\ket{11})$ : on lit 00 ou 11, jamais 01 ni 10};
\end{tikzpicture}
\caption{Préparation de l'état de Bell : une superposition (porte $H$) puis une CNOT.}\label{fig:bell}
\end{figure}

\paragraph{Sur le circuit.} Le contrôle est représenté par un point plein ($\bullet$) et la cible par le symbole $\oplus$, reliés par un trait vertical. Sur le matériel supraconducteur, la porte native à deux qubits est souvent la $CZ$, qui applique $Z$ au lieu de $X$ (soit un signe $-1$ sur $\ket{11}$ seulement). On en tire la CNOT en entourant la cible de deux $H$, puisque $HZH=X$ : $\mathrm{CNOT}=(\Id\otimes H)\,CZ\,(\Id\otimes H)$.

\paragraph{La porte $R_{ZZ}$ : le couplage du QUBO.} Le terme $J\,Z_iZ_j$ d'un coût d'Ising demande d'appliquer $e^{-i\theta Z_iZ_j/2}$. C'est une matrice diagonale, qui multiplie $\ket{b_ib_j}$ par $e^{-i\theta s_is_j/2}$ : la phase est $e^{-i\theta/2}$ si les bits sont égaux, et $e^{+i\theta/2}$ s'ils diffèrent. Comme $(Z\otimes Z)^2=\Id$, on a aussi $e^{-i\theta Z_iZ_j/2}=\cos\tfrac\theta2\,\Id-i\sin\tfrac\theta2\,Z_iZ_j$ : c'est une « rotation » à deux qubits, dont l'axe est $Z\otimes Z$. Sur les états de base, elle ne fait qu'ajouter une phase qui dépend de la \emph{parité} des deux bits :
\begin{center}
\begin{tabular}{ccccc}
\toprule
$b_i$ & $b_j$ & parité $b_i\oplus b_j$ & $s_is_j$ & phase \\
\midrule
0 & 0 & 0 & $+1$ & $e^{-i\theta/2}$\\
0 & 1 & 1 & $-1$ & $e^{+i\theta/2}$\\
1 & 0 & 1 & $-1$ & $e^{+i\theta/2}$\\
1 & 1 & 0 & $+1$ & $e^{-i\theta/2}$\\
\bottomrule
\end{tabular}
\end{center}
Aucune porte de base ne fait directement cela, mais on la construit avec deux CNOT et une rotation (figure~\ref{fig:rzz}), en trois temps :
\begin{enumerate}
\item la première CNOT écrit la parité $b_i\oplus b_j$ dans le qubit $j$ : $\ket{b_i,b_j}\mapsto\ket{b_i,\,b_i\oplus b_j}$ ;
\item $R_Z(\theta)$ sur le qubit $j$ applique $e^{-i\theta/2}$ si ce qubit vaut 0 (bits égaux), et $e^{+i\theta/2}$ s'il vaut 1 (bits différents). C'est exactement la phase voulue ;
\item la seconde CNOT « décalcule » la parité. Comme $\mathrm{CNOT}^2=\Id$, le qubit $j$ retrouve sa valeur $b_j$, et il ne reste que la phase.
\end{enumerate}
Dans QAOA, l'angle vaut $\theta=2\gamma J_{ij}$, où $J_{ij}$ est le couplage d'Ising entre les bits $i$ et $j$ (ch.~\ref{ch:qaoa}). C'est par ces portes que la structure du QUBO entre dans le circuit, et elles sont la source de l'intrication.

\begin{figure}[htbp]
\centering
\begin{tikzpicture}[x=1.1cm,y=0.9cm]
\draw[cink] (0,0) -- (1.6,0); \draw[cink] (0,-1) -- (1.6,-1);
\node[anchor=east,font=\small] at (0,0) {$q_i$}; \node[anchor=east,font=\small] at (0,-1) {$q_j$};
\node[draw=cink,fill=white,minimum width=12mm,minimum height=15mm,font=\small] at (0.8,-0.5) {$R_{ZZ}(\theta)$};
\node at (2.1,-0.5) {$=$};
\draw[cink] (2.6,0) -- (6.2,0); \draw[cink] (2.6,-1) -- (6.2,-1);
\draw[cink] (3.2,0) -- (3.2,-1); \node[ctrl] at (3.2,0) {}; \targ{3.2}{-1}
\node[gate,minimum width=12mm] at (4.4,-1) {$R_Z(\theta)$};
\draw[cink] (5.6,0) -- (5.6,-1); \node[ctrl] at (5.6,0) {}; \targ{5.6}{-1}
\end{tikzpicture}
\caption{$R_{ZZ}(\theta)=e^{-i\theta Z\otimes Z/2}=\mathrm{CNOT}\,(\Id\otimes R_Z(\theta))\,\mathrm{CNOT}$. Chaque coefficient non nul $Q_{ij}$ du QUBO coûte ainsi \textbf{deux CNOT} par couche de QAOA.}\label{fig:rzz}
\end{figure}

\section{Universalité et coût réel d'un circuit}
Les rotations à un qubit et la CNOT forment un ensemble \textbf{universel} : tout unitaire sur $n$ qubits peut s'en approcher. Deux mesures du coût :
\begin{itemize}
\item la \textbf{profondeur} : le nombre de « tranches » de portes exécutables en parallèle. Elle mesure la durée du calcul face à la durée de vie de l'information quantique (décohérence, ch.~\ref{ch:outils}) ;
\item le \textbf{nombre de portes à deux qubits} : elles sont environ dix fois plus bruitées que les portes à un qubit sur le matériel actuel. C'est le coût qui compte.
\end{itemize}
Sur un vrai processeur, seuls certains couples de qubits sont physiquement reliés. Coupler deux qubits éloignés demande des portes SWAP supplémentaires, ce qui augmente la profondeur.

\section{Mesures, shots et statistique}
Exécuter un circuit une fois donne \emph{une} chaîne de bits. On l'exécute $N$ fois (shots) pour estimer :
\begin{itemize}
\item une probabilité $P(x)$ par la fréquence $\hat P$, d'écart-type $\sqrt{P(1-P)/N}$ ;
\item une valeur moyenne $\expval{H_C}$ par la moyenne empirique de $E(x)$ sur les shots, d'écart-type $\sigma_E/\sqrt N$.
\end{itemize}
Pour l'optimisation, on ne s'intéresse pas seulement à la moyenne : on garde \textbf{la meilleure chaîne observée}, dont on recalcule le coût classiquement et exactement.

\begin{projet}
Le circuit QAOA de notre QUBO n'utilise que $H$ (initialisation), $R_Z$ (termes linéaires), $R_{ZZ}$ (couplages) et $R_X$ (mélangeur). Notre QUBO aéro a des couplages entre \emph{tous} les couples de bits : $\binom92=36$ portes $R_{ZZ}$, soit 72 CNOT par couche.
\end{projet}

\exercices
\begin{exercice}
On modifie le circuit de la figure~\ref{fig:bell} en partant de $\ket{q_0q_1}=\ket{01}$ au lieu de $\ket{00}$. Calculer pas à pas l'état obtenu : d'abord $(H\otimes\Id)\ket{01}$, puis l'effet de la CNOT. Quelles sont les probabilités des quatre résultats de mesure ? L'état final est-il intriqué ?
\end{exercice}
\begin{exercice}
Vérifier l'identité de la figure~\ref{fig:rzz} sur les quatre états de base : suivre $\ket{b_ib_j}$ à travers CNOT, $R_Z(\theta)$ sur le second qubit, puis CNOT, et montrer que la phase obtenue est $e^{-i\theta s_is_j/2}$.
\end{exercice}

\chapter{Hamiltoniens : l'énergie comme fonction de coût}\label{ch:hamiltoniens}
\iterbanner{V0}

\section{Qu'est-ce qu'un hamiltonien ?}
En mécanique classique, l'énergie totale d'un système (cinétique + potentielle) gouverne son mouvement. En quantique, l'énergie est une \emph{observable}, donc une matrice hermitienne $H$ appelée \textbf{hamiltonien}.
\begin{itemize}
\item Ses valeurs propres $E_0\le E_1\le\cdots$ sont les \textbf{niveaux d'énergie} possibles (ch.~\ref{ch:bases}).
\item Ses vecteurs propres sont les états d'énergie bien définie.
\item Le plus bas, $\ket{E_0}$, est l'\textbf{état fondamental}.
\end{itemize}
Le hamiltonien joue un double rôle : il \emph{définit l'énergie} et il \emph{génère l'évolution dans le temps}.

\section{L'équation de Schrödinger}
\begin{definition}[ (évolution temporelle)]
\[ i\frac{d}{dt}\ket{\psi(t)}=H\ket{\psi(t)}\qquad\Longrightarrow\qquad \ket{\psi(t)}=e^{-iHt}\ket{\psi(0)}\quad(H\text{ constant}). \]
\end{definition}
C'est une EDO linéaire du premier ordre, comme $\dot y=Ay$ en dynamique des systèmes, mais avec $-iH$. Comme $H$ est hermitienne, $e^{-iHt}$ est \emph{unitaire} : on retrouve le postulat d'évolution. En décomposant $\ket{\psi(0)}=\sum_kc_k\ket{E_k}$ :
\[\ket{\psi(t)}=\sum_kc_k\,e^{-iE_kt}\ket{E_k}.\]
Chaque composante \emph{tourne en phase} à la vitesse de son énergie ; les modules $|c_k|$ ne bougent pas. Un état propre ne fait que tourner d'une phase globale : il est \emph{stationnaire}.

\begin{resultat}
Faire évoluer sous un hamiltonien diagonal $H_C$ pendant un « temps » $\gamma$ multiplie chaque état de base par $e^{-i\gamma E(x)}$. \textbf{La fonction de coût est inscrite dans les phases.} Les probabilités ne changent pas : il faudra un second hamiltonien, qui ne commute pas avec $H_C$, pour transformer ces phases en différences de probabilité.
\end{resultat}

\section{État fondamental et principe variationnel}
Pour tout état normé $\ket\psi$ :
\[\bra\psi H\ket\psi=\sum_kE_k|c_k|^2\ \ge\ E_0,\]
avec égalité si et seulement si $\ket\psi$ est un état fondamental. C'est le quotient de Rayleigh d'une matrice symétrique, que vous connaissez en algèbre linéaire. Conséquence : \textbf{minimiser $\expval{H}$ sur une famille d'états paramétrés donne une borne supérieure de $E_0$}, et l'atteint si la famille est assez riche. C'est le principe des algorithmes \emph{variationnels}, dont QAOA (ch.~\ref{ch:qaoa}).

\section{Un hamiltonien diagonal est une fonction de coût}
Si $H_C=\sum_xE(x)\ket x\bra x$ (ch.~\ref{ch:multi}), ses niveaux d'énergie sont les valeurs $E(x)$ sur les $2^n$ chaînes de bits, et son état fondamental est la meilleure chaîne. La figure~\ref{fig:niveaux} montre le problème « jouet » à 3 bits du chapitre suivant.

\begin{figure}[htbp]
\centering
\begin{tikzpicture}[y=1.2cm]
\draw[->,cmuted] (-0.3,-2.4) -- (-0.3,1.5) node[above,cink,font=\small] {énergie $E(x)$};
\foreach \e/\l/\xo in {0/000/0,0/110/1.4,0/011/2.8,-1/100/0,-1/010/1.4,-1/001/2.8,1/111/1.4}{
  \draw[cink,thick] (\xo,\e) -- ++(1,0) node[right,font=\scriptsize] {$\ket{\l}$};}
\draw[cblue,very thick] (1.4,-2) -- ++(1,0) node[right,font=\scriptsize] {$\ket{101}$};
\foreach \e in {-2,-1,0,1}{\node[font=\small,anchor=east] at (-0.35,\e) {$\e$};}
\node[cblue,font=\small,anchor=west] at (4.3,-2) {état fondamental = solution};
\end{tikzpicture}
\caption{Niveaux d'énergie du hamiltonien diagonal associé au QUBO jouet du chapitre~\ref{ch:qubo} (pas deux bits voisins à 1). Plusieurs chaînes partagent le même niveau : on dit que le niveau est \emph{dégénéré}.}\label{fig:niveaux}
\end{figure}

\section{Deux hamiltoniens qui ne commutent pas : la formule de Trotter}
Si $[A,B]=AB-BA\neq0$, alors $e^{-i(A+B)t}\neq e^{-iAt}e^{-iBt}$. Mais en découpant le temps en $N$ petits pas :
\[e^{-i(A+B)t}=\lim_{N\to\infty}\big(e^{-iAt/N}e^{-iBt/N}\big)^N,\qquad \text{erreur d'un pas}\sim\tfrac{t^2}{2N^2}\|[A,B]\|.\]
C'est l'analogue exact du \emph{splitting d'opérateurs} (Lie-Trotter, Strang) utilisé en CFD pour séparer convection et diffusion. QAOA est une telle alternance, entre $e^{-i\gamma H_C}$ et $e^{-i\beta H_M}$, mais où l'on laisse \emph{libres} les durées de chaque pas.

\exercices
\begin{exercice}
Calculer $e^{-i\gamma Z}\ket+$. Montrer que c'est, à une phase globale près, $\tfrac1{\sqrt2}(\ket0+e^{2i\gamma}\ket1)$. Quelle rotation de la sphère de Bloch a-t-on effectuée ? Les probabilités en base de calcul ont-elles changé ?
\end{exercice}
\begin{exercice}
Montrer que l'état fondamental de $H_M=-\sum_{k=0}^{n-1}X_k$ est $\ket+^{\otimes n}=H^{\otimes n}\ket{0\cdots0}$, d'énergie $-n$. Pourquoi cet état est-il appelé « superposition uniforme » ?
\end{exercice}

\chapter{QUBO et modèle d'Ising}\label{ch:qubo}
\iterbanner{V0}

C'est le chapitre charnière : il traduit un problème d'optimisation d'ingénieur en un hamiltonien que le recuit ou QAOA savent traiter.

\section{Optimisation combinatoire et QUBO}
\begin{definition}[ (QUBO)]
Un problème \textbf{QUBO} (\emph{Quadratic Unconstrained Binary Optimization}) consiste à
\[\min_{\mathbf b\in\{0,1\}^n}\ E(\mathbf b)=\mathbf b^\top Q\,\mathbf b+c=\sum_iQ_{ii}b_i+\sum_{i<j}Q_{ij}b_ib_j+c,\]
avec $Q$ triangulaire supérieure. Les termes linéaires sont sur la diagonale, car $b_i^2=b_i$.
\end{definition}
\begin{itemize}
\item \emph{Quadratic} : pas de produit de trois bits ou plus.
\item \emph{Unconstrained} : toutes les $2^n$ chaînes sont admissibles. Les contraintes doivent être transformées en pénalités (§\ref{sec:penalites}).
\item \emph{Binary} : variables 0/1. Les entiers doivent être encodés (§\ref{sec:encodage}).
\end{itemize}
Le QUBO général est \textbf{NP-difficile} : il contient Max-Cut, le voyageur de commerce, l'ordonnancement… On ne connaît pas d'algorithme, classique ou quantique, qui le résolve en temps polynomial dans le pire cas. Recuit et QAOA sont des heuristiques.

\section{Deux exemples}
\paragraph{Le QUBO jouet du cours.} Trois bits sur une chaîne $0-1-2$ ; on veut en activer le plus possible sans que deux voisins soient actifs (« ensemble indépendant maximal ») :
\[E=-b_0-b_1-b_2+2b_0b_1+2b_1b_2,\qquad Q=\begin{pmatrix}-1&2&0\\0&-1&2\\0&0&-1\end{pmatrix}.\]
Chaque bit actif rapporte $-1$, chaque paire de voisins actifs coûte $+2$. Énumération : $E(101)=-2$ est le minimum unique (figure~\ref{fig:niveaux}). Ce problème à 3 qubits sert de fil rouge pour le recuit et QAOA, car tout s'y calcule à la main ou presque.

\paragraph{Max-Cut.} Couper un graphe en deux groupes en maximisant le nombre d'arêtes coupées. Avec $b_i$ le groupe du nœud $i$, l'arête $(i,j)$ est coupée si $b_i\neq b_j$, ce qui s'écrit $b_i+b_j-2b_ib_j=1$. On minimise donc $E=-\sum_{(i,j)}(b_i+b_j-2b_ib_j)$. C'est l'exemple historique de QAOA (Farhi et al., 2014).

\section{Du QUBO au modèle d'Ising}
Le matériel quantique manipule des \emph{spins} $s_i=\pm1$, valeurs propres de $Z_i$, plutôt que des bits. On substitue $b_i=(1-s_i)/2$ :
\begin{definition}[ (modèle d'Ising)]
\[E(\mathbf s)=\sum_ih_is_i+\sum_{i<j}J_{ij}s_is_j+\text{cte},\]
\[J_{ij}=\frac{Q_{ij}}4,\qquad h_i=-\frac{Q_{ii}}2-\frac14\sum_{j\neq i}Q_{\min(i,j)\max(i,j)},\qquad \text{cte}=c+\frac12\sum_iQ_{ii}+\frac14\sum_{i<j}Q_{ij}.\]
L'\textbf{hamiltonien de coût} s'obtient en remplaçant $s_i$ par l'opérateur $Z_i$ :
\[H_C=\sum_ih_iZ_i+\sum_{i<j}J_{ij}Z_iZ_j+\text{cte}\cdot\Id .\]
\end{definition}
\begin{itemize}
\item $h_i$ est un \emph{champ local}, qui favorise un spin particulier.
\item $J_{ij}>0$ favorise des spins opposés (antiferromagnétique) ; $J_{ij}<0$ favorise des spins égaux (ferromagnétique).
\item Le vocabulaire vient du magnétisme. Ising (1925) modélisait des aimants ; c'est le même objet mathématique.
\end{itemize}
Pour le jouet : $h=(0,\,-\tfrac12,\,0)$, $J_{01}=J_{12}=\tfrac12$, constante $-\tfrac12$ (exercice~\ref{ch:qubo}.1).

\begin{figure}[htbp]
\centering
\begin{tikzpicture}[node distance=22mm,sp/.style={circle,draw=cink,thick,minimum size=9mm,font=\small}]
\node[sp] (0) {$s_0$}; \node[sp,right=of 0] (1) {$s_1$}; \node[sp,right=of 1] (2) {$s_2$};
\draw[corange,very thick] (0) -- node[above,font=\small] {$J_{01}=+\frac12$} (1);
\draw[corange,very thick] (1) -- node[above,font=\small] {$J_{12}=+\frac12$} (2);
\node[below=2mm of 0,font=\small] {$h_0=0$}; \node[below=2mm of 1,font=\small] {$h_1=-\frac12$}; \node[below=2mm of 2,font=\small] {$h_2=0$};
\end{tikzpicture}
\caption{Le QUBO jouet comme graphe d'Ising : un nœud par qubit, une arête par couplage non nul. Le nombre d'arêtes fixe le nombre de portes $R_{ZZ}$ (QAOA) ou de coupleurs (recuit quantique).}\label{fig:graphe-jouet}
\end{figure}

\section{Contraintes : les pénalités}\label{sec:penalites}
Un QUBO n'a pas de contraintes. On les ajoute à l'énergie sous forme de \textbf{pénalité}, nulle si la contrainte est satisfaite et positive sinon :
\begin{itemize}
\item \textbf{égalité} $g(\mathbf b)=0$ : ajouter $P\,g(\mathbf b)^2$. Exemple, « exactement un bit actif parmi $k$ » (one-hot) : $P\big(\sum_kb_k-1\big)^2$ ;
\item \textbf{inégalité} $g(\mathbf b)\le0$ avec $g$ entier : introduire des bits d'écart $a_l$ et imposer $g(\mathbf b)+\sum_l2^la_l=0$ par pénalité. Cela coûte des qubits supplémentaires.
\end{itemize}
\begin{piege}[ (degré)]
La pénalité est un \emph{carré}. Si $g$ est linéaire en bits, $g^2$ est quadratique : on reste dans le QUBO. Si $g$ est déjà quadratique (par exemple $C_L$ approché par un modèle quadratique), $g^2$ est de \textbf{degré 4} : ce n'est plus un QUBO. Il faut alors linéariser $g$, ou reformuler le problème pour que la contrainte disparaisse.
\end{piege}
\paragraph{Réglage de $P$.} $P$ doit être assez grand pour qu'aucune violation ne soit rentable, mais pas trop : un $P$ énorme écrase les écarts d'énergie utiles. QAOA devient alors moins sensible au coût, et le recuit quantique perd en précision (coefficients bornés et quantifiés sur le matériel).

\section{Encoder des entiers en bits}\label{sec:encodage}
Une variable entière $x\in\{0,\dots,L-1\}$ peut s'écrire de plusieurs façons :
\begin{center}
\begin{tabular}{lll c}
\toprule
Encodage & $x=$ & contrainte & bits pour $L=8$\\
\midrule
binaire & $\sum_{j}2^jb_j$ & aucune ($L=2^k$) & 3\\
one-hot & $\sum_jj\,b_j$ & $\sum_jb_j=1$ (pénalité) & 8\\
unaire (« thermomètre ») & $\sum_jb_j$ & ordre $b_{j+1}\le b_j$, ou aucune & 7\\
\bottomrule
\end{tabular}
\end{center}
Le binaire est le plus économe en qubits, mais deux niveaux voisins peuvent différer de plusieurs bits : passer de 3 (011) à 4 (100) en change trois. Cela rend le paysage d'énergie plus « rugueux » pour un algorithme qui modifie un bit à la fois, comme le recuit simulé.

\section{Notre QUBO : du modèle de traînée à la matrice $Q$}\label{sec:notre-qubo}
\paragraph{Variables.} $m$, $p$, $t$ prennent chacune 8 niveaux d'indice $x_i\in\{0,\dots,7\}$, codés en binaire sur 3 bits : $x_i=b_{i0}+2b_{i1}+4b_{i2}$. Cela fait 9 bits au total, donc \textbf{9 qubits}.

\paragraph{Modèle.} On ajuste par moindres carrés, sur des profils simulés par la phase~1, un polynôme quadratique en les variables normalisées $z_i=2x_i/7-1\in[-1,1]$ :
\[\widehat{C_D}(z)=c_0+\mathbf g^\top\mathbf z+\mathbf z^\top M\mathbf z\qquad(1+3+6=10\text{ coefficients}).\]

\paragraph{Passage en bits.} Chaque $z_i$ est \emph{affine} en bits, $\mathbf z=\mathbf a_0+A\mathbf b$, avec $\mathbf a_0=-\mathbf 1$ et $A_{i,(i,j)}=2\cdot2^j/7$. En substituant :
\[\widehat{C_D}=\underbrace{c_0+\mathbf g^\top\mathbf a_0+\mathbf a_0^\top M\mathbf a_0}_{c}+\underbrace{(\mathbf g+2M\mathbf a_0)^\top A\,\mathbf b}_{\text{linéaire}\to\text{diag}(Q)}+\mathbf b^\top\underbrace{(A^\top MA)}_{\text{quadratique}}\mathbf b .\]
Un polynôme de degré 2 en $z$ reste de degré 2 en bits : \textbf{c'est un QUBO, sans approximation}. Le code \texttt{quantique/qubo.py} (\texttt{surrogate\_to\_qubo}) le vérifie : énergie QUBO et prédiction du modèle coïncident à $10^{-14}$ près sur les 512 profils.

\begin{figure}[htbp]
\centering
% figure générée par tools/make_figures.py : matrice_Q
\begin{tikzpicture}[scale=0.62]
\fill[corange!10] (0,0) rectangle ++(1,-1);
\node[font=\tiny] at (0.5,-0.5) {1.3};
\fill[corange!14] (1,0) rectangle ++(1,-1);
\node[font=\tiny] at (1.5,-0.5) {1.8};
\fill[corange!28] (2,0) rectangle ++(1,-1);
\node[font=\tiny] at (2.5,-0.5) {3.6};
\fill[cblue!4] (3,0) rectangle ++(1,-1);
\node[font=\tiny] at (3.5,-0.5) {-0.5};
\fill[cblue!8] (4,0) rectangle ++(1,-1);
\node[font=\tiny] at (4.5,-0.5) {-1.1};
\fill[cblue!17] (5,0) rectangle ++(1,-1);
\node[font=\tiny] at (5.5,-0.5) {-2.1};
\fill[corange!2] (6,0) rectangle ++(1,-1);
\node[font=\tiny] at (6.5,-0.5) {0.2};
\fill[corange!4] (7,0) rectangle ++(1,-1);
\node[font=\tiny] at (7.5,-0.5) {0.5};
\fill[corange!8] (8,0) rectangle ++(1,-1);
\node[font=\tiny] at (8.5,-0.5) {1.0};
\fill[corange!28] (1,-1) rectangle ++(1,-1);
\node[font=\tiny] at (1.5,-1.5) {3.6};
\fill[corange!55] (2,-1) rectangle ++(1,-1);
\node[font=\tiny] at (2.5,-1.5) {7.1};
\fill[cblue!8] (3,-1) rectangle ++(1,-1);
\node[font=\tiny] at (3.5,-1.5) {-1.1};
\fill[cblue!17] (4,-1) rectangle ++(1,-1);
\node[font=\tiny] at (4.5,-1.5) {-2.1};
\fill[cblue!33] (5,-1) rectangle ++(1,-1);
\node[font=\tiny] at (5.5,-1.5) {-4.3};
\fill[corange!4] (6,-1) rectangle ++(1,-1);
\node[font=\tiny] at (6.5,-1.5) {0.5};
\fill[corange!8] (7,-1) rectangle ++(1,-1);
\node[font=\tiny] at (7.5,-1.5) {1.0};
\fill[corange!15] (8,-1) rectangle ++(1,-1);
\node[font=\tiny] at (8.5,-1.5) {2.0};
\fill[corange!83] (2,-2) rectangle ++(1,-1);
\node[font=\tiny] at (2.5,-2.5) {10.7};
\fill[cblue!17] (3,-2) rectangle ++(1,-1);
\node[font=\tiny] at (3.5,-2.5) {-2.1};
\fill[cblue!33] (4,-2) rectangle ++(1,-1);
\node[font=\tiny] at (4.5,-2.5) {-4.3};
\fill[cblue!66] (5,-2) rectangle ++(1,-1);
\node[font=\tiny] at (5.5,-2.5) {-8.5};
\fill[corange!8] (6,-2) rectangle ++(1,-1);
\node[font=\tiny] at (6.5,-2.5) {1.0};
\fill[corange!15] (7,-2) rectangle ++(1,-1);
\node[font=\tiny] at (7.5,-2.5) {2.0};
\fill[corange!31] (8,-2) rectangle ++(1,-1);
\node[font=\tiny] at (8.5,-2.5) {3.9};
\fill[cblue!7] (3,-3) rectangle ++(1,-1);
\node[font=\tiny] at (3.5,-3.5) {-0.9};
\fill[corange!15] (4,-3) rectangle ++(1,-1);
\node[font=\tiny] at (4.5,-3.5) {1.9};
\fill[corange!29] (5,-3) rectangle ++(1,-1);
\node[font=\tiny] at (5.5,-3.5) {3.8};
\fill[cblue!1] (6,-3) rectangle ++(1,-1);
\node[font=\tiny] at (6.5,-3.5) {-0.2};
\fill[cblue!3] (7,-3) rectangle ++(1,-1);
\node[font=\tiny] at (7.5,-3.5) {-0.3};
\fill[cblue!5] (8,-3) rectangle ++(1,-1);
\node[font=\tiny] at (8.5,-3.5) {-0.7};
\fill[cblue!6] (4,-4) rectangle ++(1,-1);
\node[font=\tiny] at (4.5,-4.5) {-0.8};
\fill[corange!59] (5,-4) rectangle ++(1,-1);
\node[font=\tiny] at (5.5,-4.5) {7.6};
\fill[cblue!3] (6,-4) rectangle ++(1,-1);
\node[font=\tiny] at (6.5,-4.5) {-0.3};
\fill[cblue!5] (7,-4) rectangle ++(1,-1);
\node[font=\tiny] at (7.5,-4.5) {-0.7};
\fill[cblue!11] (8,-4) rectangle ++(1,-1);
\node[font=\tiny] at (8.5,-4.5) {-1.4};
\fill[corange!17] (5,-5) rectangle ++(1,-1);
\node[font=\tiny] at (5.5,-5.5) {2.1};
\fill[cblue!5] (6,-5) rectangle ++(1,-1);
\node[font=\tiny] at (6.5,-5.5) {-0.7};
\fill[cblue!11] (7,-5) rectangle ++(1,-1);
\node[font=\tiny] at (7.5,-5.5) {-1.4};
\fill[cblue!22] (8,-5) rectangle ++(1,-1);
\node[font=\tiny] at (8.5,-5.5) {-2.8};
\fill[corange!22] (6,-6) rectangle ++(1,-1);
\node[font=\tiny] at (6.5,-6.5) {2.8};
\fill[cblue!1] (7,-6) rectangle ++(1,-1);
\node[font=\tiny] at (7.5,-6.5) {-0.1};
\fill[cblue!1] (8,-6) rectangle ++(1,-1);
\node[font=\tiny] at (8.5,-6.5) {-0.2};
\fill[corange!43] (7,-7) rectangle ++(1,-1);
\node[font=\tiny] at (7.5,-7.5) {5.5};
\fill[cblue!3] (8,-7) rectangle ++(1,-1);
\node[font=\tiny] at (8.5,-7.5) {-0.3};
\fill[corange!85] (8,-8) rectangle ++(1,-1);
\node[font=\tiny] at (8.5,-8.5) {10.9};
\draw[cmuted,thin] (0,0) grid (9,-9);
\draw[cink,thick] (0,0) rectangle (3,-3) (3,-3) rectangle (6,-6) (6,-6) rectangle (9,-9);
\node[font=\scriptsize] at (0.5,0.35) {$m_0$};
\node[font=\scriptsize] at (1.5,0.35) {$m_1$};
\node[font=\scriptsize] at (2.5,0.35) {$m_2$};
\node[font=\scriptsize] at (3.5,0.35) {$p_0$};
\node[font=\scriptsize] at (4.5,0.35) {$p_1$};
\node[font=\scriptsize] at (5.5,0.35) {$p_2$};
\node[font=\scriptsize] at (6.5,0.35) {$t_0$};
\node[font=\scriptsize] at (7.5,0.35) {$t_1$};
\node[font=\scriptsize] at (8.5,0.35) {$t_2$};
\node[font=\scriptsize] at (-0.45,-0.5) {$m_0$};
\node[font=\scriptsize] at (-0.45,-1.5) {$m_1$};
\node[font=\scriptsize] at (-0.45,-2.5) {$m_2$};
\node[font=\scriptsize] at (-0.45,-3.5) {$p_0$};
\node[font=\scriptsize] at (-0.45,-4.5) {$p_1$};
\node[font=\scriptsize] at (-0.45,-5.5) {$p_2$};
\node[font=\scriptsize] at (-0.45,-6.5) {$t_0$};
\node[font=\scriptsize] at (-0.45,-7.5) {$t_1$};
\node[font=\scriptsize] at (-0.45,-8.5) {$t_2$};
\end{tikzpicture}

\caption{Matrice $Q$ (en counts de traînée) du QUBO aéro, modèle ajusté sur les 512 profils à $\alpha=2^\circ$, $Re=3\cdot10^6$. Bleu : négatif, orange : positif. Les blocs encadrés couplent les 3 bits d'une même variable (termes $z_i^2$) ; les blocs hors diagonale viennent des termes croisés $z_iz_j$. Tous les coefficients hors diagonale sont non nuls : le graphe d'Ising est \textbf{complet} (36 arêtes).}\label{fig:matrice-q}
\end{figure}

\begin{projet}[ : ce que dit le QUBO, et ce qu'il ne dit pas]
\begin{itemize}
\item Minimum du QUBO : $57{,}16$ counts pour le profil $(m,p,t)=(0,\ 0{,}25,\ 6\,\%)$. Le second niveau, $57{,}22$, est le profil $(0,\ 0{,}30,\ 6\,\%)$. Ces deux profils sont \emph{physiquement identiques}, puisque $p$ n'a pas d'effet quand $m=0$ ! Le modèle les sépare de $0{,}06$ count seulement : le fondamental est \emph{quasi dégénéré}.
\item Le vrai $C_D$ de ce profil vaut $61{,}1$ counts : il n'est que \textbf{38\textsuperscript{e} sur 512}. Le vrai minimum, $50{,}5$, est ailleurs. Le QUBO résout \emph{exactement} un problème \emph{approché} : l'erreur vient du modèle quadratique, pas de l'optimiseur. C'est pourquoi la note d'intention prévoit une \textbf{boucle itérative} qui simule les candidats et réajuste le modèle (ch.~\ref{ch:projet}).
\item Le graphe complet a une conséquence matérielle : 36 portes $R_{ZZ}$ par couche de QAOA, et un plongement coûteux sur un recuit quantique (ch.~\ref{ch:recuit}).
\end{itemize}
\end{projet}

\exercices
\begin{exercice}
Convertir le QUBO jouet en modèle d'Ising : calculer $h$, $J$ et la constante avec les formules du cours. Vérifier que $E(\mathbf s)=-2$ pour la chaîne $101$, c'est-à-dire $\mathbf s=(-1,+1,-1)$.
\end{exercice}
\begin{exercice}
On encode $m$, $p$, $t$ (8 niveaux chacun) en one-hot. Combien faut-il de qubits ? Écrire la pénalité qui impose « un seul niveau actif » pour $m$, et dire combien de couplages $Z_iZ_j$ elle ajoute. Conclure sur le choix d'encodage pour notre projet.
\end{exercice}

\chapter{Recuit simulé et recuit quantique}\label{ch:recuit}
\iterbanner{V0}

Deux façons de chercher l'état fondamental en « refroidissant » progressivement : avec du hasard thermique (recuit \emph{simulé}, classique) ou avec des fluctuations quantiques (recuit \emph{quantique}).

\section{Le paysage d'énergie}
On imagine les $2^n$ chaînes comme les points d'un paysage, où deux chaînes sont voisines si elles diffèrent d'un bit, et l'altitude est l'énergie. Une descente « gloutonne », qui retourne un bit dès que l'énergie baisse, s'arrête dans le premier \textbf{minimum local} venu.

\begin{figure}[htbp]
\centering
\begin{tikzpicture}[x=1.3cm,y=1.1cm]
\draw[->,cmuted] (0,-0.2) -- (0,3.1) node[above,cink,font=\small] {énergie};
\draw[->,cmuted] (0,-0.2) -- (9.3,-0.2) node[below left,cink,font=\small] {configurations (chaînes de bits)};
\draw[cink,very thick] plot[smooth,tension=0.7] coordinates {(0.2,2.3) (1.2,1.6) (2.2,1.2) (3.2,1.7) (4.3,2.7) (5.3,1.8) (6.4,0.25) (7.4,1.0) (8.4,2.4) (9.0,2.6)};
\fill[cblue] (2.2,1.2) circle (2.2pt);
\node[font=\small,anchor=north east] at (2.1,1.1) {minimum local};
\node[font=\small,anchor=west] at (6.6,0.2) {minimum global};
\draw[-{Stealth},corange,thick] (2.3,1.35) .. controls (3.4,3.6) and (4.9,3.6) .. (5.7,1.6);
\node[corange,font=\small,above] at (4.0,3.15) {saut thermique (recuit simulé)};
\draw[-{Stealth},cblue,thick,dashed] (2.4,1.25) -- (5.75,1.25);
\node[cblue,font=\small,align=center] at (4.0,0.75) {effet tunnel\\ (recuit quantique)};
\end{tikzpicture}
\caption{Deux façons de quitter un minimum local : passer \emph{par-dessus} la barrière grâce à l'agitation thermique, ou \emph{à travers} grâce aux fluctuations quantiques.}\label{fig:paysage}
\end{figure}

\section{Recuit simulé (classique)}
Inspiré de la métallurgie : un métal refroidi lentement cristallise dans un état de basse énergie. On introduit une \textbf{température} $T$ et on fait une marche aléatoire (algorithme de Metropolis) :
\begin{enumerate}
\item proposer de retourner un bit, ce qui change l'énergie de $\Delta E$ ;
\item accepter si $\Delta E\le0$, sinon accepter avec probabilité $e^{-\Delta E/T}$ ;
\item diminuer lentement $T$ (programme de température, par exemple géométrique).
\end{enumerate}
À $T$ fixée, la marche échantillonne la distribution de Boltzmann $\propto e^{-E(x)/T}$, qui se concentre sur le minimum global quand $T\to0$. Si $T$ baisse trop vite, on reste piégé ; la théorie garantit la convergence pour un refroidissement logarithmique, beaucoup trop lent en pratique.

Retourner le bit $i$ coûte $\Delta E_i=(1-2b_i)\big(Q_{ii}+\sum_{j\neq i}Q_{\min,\max}b_j\big)$. C'est un calcul en $O(n)$, ce qui rend le recuit simulé très rapide (\texttt{quantique/recuit.py}).

\begin{figure}[htbp]
\centering
% figure générée par tools/make_figures.py : recuit_succes
\begin{tikzpicture}
\begin{semilogxaxis}[width=0.7\linewidth,height=5.2cm,axis lines=left,ymin=0,ymax=1.02,
  xlabel={nombre de balayages (pas de température)},ylabel={fraction des 200 lectures},
  legend style={draw=none,font=\small,at={(0.98,0.03)},anchor=south east}]
\addplot[cblue,thick,mark=*] coordinates {(2,0.16) (5,0.33) (10,0.37) (20,0.51) (50,0.445) (100,0.49)};\addlegendentry{minimum exact du QUBO}
\addplot[corange,thick,mark=square*] coordinates {(2,0.435) (5,0.79) (10,0.89) (20,0.94) (50,0.965) (100,0.98)};\addlegendentry{à moins de 0{,}1 count du minimum}
\end{semilogxaxis}
\end{tikzpicture}

\caption{Recuit simulé sur notre QUBO aéro à 9 bits, 200 lectures indépendantes par point. La réussite « exacte » plafonne vers 50~\% : le fondamental et le premier niveau excité sont deux profils physiquement identiques ($m=0$, $p=0{,}25$ ou $0{,}30$), séparés de $0{,}06$ count. Le recuit tombe au hasard dans l'un des deux. À $0{,}1$ count près, la réussite dépasse 95~\% dès 50 balayages.}\label{fig:recuit}
\end{figure}

\section{Le calcul quantique adiabatique}
\paragraph{Idée.} On prépare un système dans l'état fondamental, \emph{facile}, d'un hamiltonien simple $H_B$, puis on le transforme \emph{lentement} en notre $H_C$ :
\[H(s)=(1-s)\,H_B+s\,H_C,\qquad s=t/T\in[0,1],\qquad H_B=-\sum_kX_k .\]
L'état fondamental de $H_B$ est $\ket+^{\otimes n}$, la superposition uniforme de toutes les chaînes (exercice~\ref{ch:hamiltoniens}.2).

\paragraph{Théorème adiabatique.} Si l'évolution est assez lente, le système reste dans l'état fondamental instantané de $H(s)$, et arrive donc dans celui de $H_C$ : la solution. « Assez lent » dépend du plus petit \textbf{écart spectral} (\emph{gap}) $g_{\min}$ entre les deux niveaux les plus bas le long du chemin : il faut typiquement $T\gg1/g_{\min}^2$.

\begin{figure}[htbp]
\centering
% figure générée par tools/make_figures.py : spectre_adiabatique
\begin{tikzpicture}
\begin{axis}[width=0.8\linewidth,height=6.2cm,axis lines=left,xmin=0,xmax=1,ymin=-3.3,ymax=1.3,
  xlabel={$s=t/T$},ylabel={niveaux d'énergie de $H(s)$},
  legend style={draw=none,font=\small,at={(0.02,0.02)},anchor=south west,fill=white,fill opacity=0.85,text opacity=1}]
\addplot[cblue,very thick] coordinates {(0,-3) (0.02,-2.95) (0.04,-2.9) (0.06,-2.851) (0.08,-2.802) (0.1,-2.753) (0.12,-2.704) (0.14,-2.656) (0.16,-2.608) (0.18,-2.56) (0.2,-2.513) (0.22,-2.466) (0.24,-2.419) (0.26,-2.373) (0.28,-2.327) (0.3,-2.282) (0.32,-2.238) (0.34,-2.194) (0.36,-2.152) (0.38,-2.109) (0.4,-2.068) (0.42,-2.028) (0.44,-1.989) (0.46,-1.952) (0.48,-1.916) (0.5,-1.881) (0.52,-1.849) (0.54,-1.818) (0.56,-1.79) (0.58,-1.765) (0.6,-1.742) (0.62,-1.723) (0.64,-1.707) (0.66,-1.695) (0.68,-1.686) (0.7,-1.682) (0.72,-1.681) (0.74,-1.684) (0.76,-1.69) (0.78,-1.7) (0.8,-1.714) (0.82,-1.73) (0.84,-1.75) (0.86,-1.773) (0.88,-1.798) (0.9,-1.826) (0.92,-1.856) (0.94,-1.889) (0.96,-1.924) (0.98,-1.961) (1,-2)};
\addplot[corange,thick] coordinates {(0,-1) (0.02,-1.004) (0.04,-1.008) (0.06,-1.013) (0.08,-1.017) (0.1,-1.021) (0.12,-1.026) (0.14,-1.03) (0.16,-1.035) (0.18,-1.039) (0.2,-1.044) (0.22,-1.048) (0.24,-1.053) (0.26,-1.057) (0.28,-1.062) (0.3,-1.066) (0.32,-1.071) (0.34,-1.075) (0.36,-1.079) (0.38,-1.083) (0.4,-1.087) (0.42,-1.091) (0.44,-1.094) (0.46,-1.097) (0.48,-1.099) (0.5,-1.101) (0.52,-1.101) (0.54,-1.101) (0.56,-1.1) (0.58,-1.098) (0.6,-1.095) (0.62,-1.089) (0.64,-1.083) (0.66,-1.075) (0.68,-1.065) (0.7,-1.054) (0.72,-1.042) (0.74,-1.029) (0.76,-1.015) (0.78,-1.002) (0.8,-0.9892) (0.82,-0.9774) (0.84,-0.9673) (0.86,-0.9593) (0.88,-0.954) (0.9,-0.9518) (0.92,-0.9533) (0.94,-0.9587) (0.96,-0.9682) (0.98,-0.982) (1,-1)};
\addplot[cmuted,thin] coordinates {(0,-1) (0.02,-0.9901) (0.04,-0.9802) (0.06,-0.9705) (0.08,-0.9609) (0.1,-0.9514) (0.12,-0.942) (0.14,-0.9328) (0.16,-0.9238) (0.18,-0.9149) (0.2,-0.9062) (0.22,-0.8977) (0.24,-0.8894) (0.26,-0.8813) (0.28,-0.8735) (0.3,-0.8659) (0.32,-0.8586) (0.34,-0.8515) (0.36,-0.8448) (0.38,-0.8385) (0.4,-0.8325) (0.42,-0.8268) (0.44,-0.8217) (0.46,-0.8169) (0.48,-0.8127) (0.5,-0.809) (0.52,-0.8059) (0.54,-0.8034) (0.56,-0.8015) (0.58,-0.8004) (0.6,-0.8) (0.62,-0.8004) (0.64,-0.8017) (0.66,-0.8038) (0.68,-0.8069) (0.7,-0.811) (0.72,-0.8161) (0.74,-0.8222) (0.76,-0.8294) (0.78,-0.8378) (0.8,-0.8472) (0.82,-0.8578) (0.84,-0.8694) (0.86,-0.8822) (0.88,-0.8961) (0.9,-0.911) (0.92,-0.9269) (0.94,-0.9438) (0.96,-0.9617) (0.98,-0.9804) (1,-1)};
\addplot[cmuted,thin] coordinates {(0,-1) (0.02,-0.9759) (0.04,-0.9518) (0.06,-0.9278) (0.08,-0.9039) (0.1,-0.8801) (0.12,-0.8563) (0.14,-0.8327) (0.16,-0.8092) (0.18,-0.7859) (0.2,-0.7627) (0.22,-0.7397) (0.24,-0.7169) (0.26,-0.6945) (0.28,-0.6723) (0.3,-0.6505) (0.32,-0.6291) (0.34,-0.6082) (0.36,-0.588) (0.38,-0.5685) (0.4,-0.55) (0.42,-0.5326) (0.44,-0.5166) (0.46,-0.5023) (0.48,-0.4903) (0.5,-0.4811) (0.52,-0.4752) (0.54,-0.4734) (0.56,-0.4762) (0.58,-0.484) (0.6,-0.4966) (0.62,-0.5137) (0.64,-0.5346) (0.66,-0.5585) (0.68,-0.5845) (0.7,-0.6121) (0.72,-0.6406) (0.74,-0.6696) (0.76,-0.6988) (0.78,-0.7278) (0.8,-0.7564) (0.82,-0.7845) (0.84,-0.8119) (0.86,-0.8385) (0.88,-0.8643) (0.9,-0.8892) (0.92,-0.9132) (0.94,-0.9362) (0.96,-0.9583) (0.98,-0.9796) (1,-1)};
\addplot[cmuted,thin] coordinates {(0,1) (0.02,0.9559) (0.04,0.9118) (0.06,0.8678) (0.08,0.8239) (0.1,0.7801) (0.12,0.7363) (0.14,0.6927) (0.16,0.6492) (0.18,0.6059) (0.2,0.5627) (0.22,0.5197) (0.24,0.4769) (0.26,0.4345) (0.28,0.3923) (0.3,0.3505) (0.32,0.3091) (0.34,0.2682) (0.36,0.228) (0.38,0.1885) (0.4,0.15) (0.42,0.1126) (0.44,0.07657) (0.46,0.04234) (0.48,0.01033) (0.5,-0.01891) (0.52,-0.04477) (0.54,-0.06657) (0.56,-0.08376) (0.58,-0.09601) (0.6,-0.1034) (0.62,-0.1063) (0.64,-0.1054) (0.66,-0.1015) (0.68,-0.09548) (0.7,-0.08792) (0.72,-0.07941) (0.74,-0.0704) (0.76,-0.06125) (0.78,-0.05224) (0.8,-0.0436) (0.82,-0.03551) (0.84,-0.0281) (0.86,-0.02146) (0.88,-0.01568) (0.9,-0.0108) (0.92,-0.006836) (0.94,-0.003794) (0.96,-0.00166) (0.98,-0.0004078) (1,0)};
\addplot[cmuted,thin] coordinates {(0,1) (0.02,0.9701) (0.04,0.9402) (0.06,0.9105) (0.08,0.8809) (0.1,0.8514) (0.12,0.822) (0.14,0.7928) (0.16,0.7638) (0.18,0.7349) (0.2,0.7062) (0.22,0.6777) (0.24,0.6494) (0.26,0.6213) (0.28,0.5935) (0.3,0.5659) (0.32,0.5386) (0.34,0.5115) (0.36,0.4848) (0.38,0.4585) (0.4,0.4325) (0.42,0.4068) (0.44,0.3817) (0.46,0.3569) (0.48,0.3327) (0.5,0.309) (0.52,0.2859) (0.54,0.2634) (0.56,0.2415) (0.58,0.2204) (0.6,0.2) (0.62,0.1804) (0.64,0.1617) (0.66,0.1438) (0.68,0.1269) (0.7,0.111) (0.72,0.09607) (0.74,0.08222) (0.76,0.06944) (0.78,0.05777) (0.8,0.04721) (0.82,0.03777) (0.84,0.02944) (0.86,0.02222) (0.88,0.01607) (0.9,0.01098) (0.92,0.006905) (0.94,0.003814) (0.96,0.001664) (0.98,0.000408) (1,0)};
\addplot[cmuted,thin] coordinates {(0,1) (0.02,0.9842) (0.04,0.9684) (0.06,0.9526) (0.08,0.937) (0.1,0.9213) (0.12,0.9057) (0.14,0.8901) (0.16,0.8746) (0.18,0.8591) (0.2,0.8436) (0.22,0.8282) (0.24,0.8127) (0.26,0.7973) (0.28,0.7818) (0.3,0.7662) (0.32,0.7506) (0.34,0.735) (0.36,0.7192) (0.38,0.7032) (0.4,0.687) (0.42,0.6706) (0.44,0.6538) (0.46,0.6366) (0.48,0.6189) (0.5,0.6005) (0.52,0.5814) (0.54,0.5614) (0.56,0.5404) (0.58,0.5181) (0.6,0.4945) (0.62,0.4694) (0.64,0.4428) (0.66,0.4146) (0.68,0.3849) (0.7,0.3538) (0.72,0.3217) (0.74,0.2887) (0.76,0.2553) (0.78,0.222) (0.8,0.1892) (0.82,0.1574) (0.84,0.1273) (0.86,0.09931) (0.88,0.07399) (0.9,0.05184) (0.92,0.0333) (0.94,0.01869) (0.96,0.008247) (0.98,0.002036) (1,0)};
\addplot[cmuted,thin] coordinates {(0,3) (0.02,2.93) (0.04,2.86) (0.06,2.791) (0.08,2.722) (0.1,2.653) (0.12,2.584) (0.14,2.516) (0.16,2.448) (0.18,2.38) (0.2,2.313) (0.22,2.246) (0.24,2.179) (0.26,2.113) (0.28,2.047) (0.3,1.982) (0.32,1.918) (0.34,1.854) (0.36,1.792) (0.38,1.729) (0.4,1.668) (0.42,1.608) (0.44,1.549) (0.46,1.492) (0.48,1.436) (0.5,1.381) (0.52,1.329) (0.54,1.278) (0.56,1.23) (0.58,1.185) (0.6,1.142) (0.62,1.103) (0.64,1.067) (0.66,1.035) (0.68,1.006) (0.7,0.9817) (0.72,0.9608) (0.74,0.9437) (0.76,0.9303) (0.78,0.9203) (0.8,0.9137) (0.82,0.9103) (0.84,0.91) (0.86,0.9125) (0.88,0.9179) (0.9,0.9258) (0.92,0.9362) (0.94,0.9489) (0.96,0.9639) (0.98,0.981) (1,1)};
\draw[cink,<->,thick] (axis cs:0.670,-1.690) -- (axis cs:0.670,-1.070)
  node[midway,right,font=\small,fill=white] {$g_{\min}=0.62$};
\node[font=\small,anchor=west] at (axis cs:0.02,-2.75) {$H_B=-\sum_k X_k$ : fondamental $|{+}{+}{+}\rangle$};
\node[font=\small,anchor=east] at (axis cs:0.98,-2.25) {$H_C$ : fondamental $|101\rangle$, $E=-2$};
\end{axis}
\end{tikzpicture}

\caption{Les 8 niveaux d'énergie de $H(s)$ pour le QUBO jouet à 3 qubits (diagonalisation exacte). Le fondamental (bleu) relie $\ket{+++}$ à $\ket{101}$. L'écart minimal avec le premier excité (orange) fixe la durée nécessaire du recuit.}\label{fig:spectre}
\end{figure}

\paragraph{Pourquoi $H_B$ doit-il ne pas commuter avec $H_C$ ?} Si $[H_B,H_C]=0$, les deux auraient les mêmes vecteurs propres, et le système resterait dans $\ket+^{\otimes n}$ sans jamais « choisir » de chaîne. Le terme $X_k$, appelé \emph{champ transverse}, permet aux amplitudes de passer d'une chaîne à ses voisines. C'est l'effet tunnel de la figure~\ref{fig:paysage}.

\paragraph{La limite fondamentale.} Pour les problèmes difficiles, $g_{\min}$ peut décroître \emph{exponentiellement} avec $n$, ce qui impose un temps exponentiel : l'adiabatique ne fait pas de miracle sur les problèmes NP-difficiles.

\section{Le recuit quantique en pratique (D-Wave)}
Les machines D-Wave réalisent physiquement ce programme avec des milliers de qubits supraconducteurs, mais :
\begin{itemize}
\item le système n'est pas isolé (température non nulle, bruit). Ce n'est pas un calcul adiabatique idéal, et le résultat est un \emph{échantillon} de chaînes de basse énergie. On lance de nombreuses lectures (\emph{reads}) ;
\item les qubits ne sont couplés qu'à quelques voisins sur une grille fixe. Un graphe complet comme le nôtre (36 arêtes) exige un \textbf{plongement} (\emph{embedding}) : chaque qubit logique est représenté par une \emph{chaîne} de qubits physiques, maintenus égaux par de forts couplages ferromagnétiques. Les chaînes peuvent se « casser » ;
\item les coefficients $h_i$, $J_{ij}$ ont une plage et une précision limitées. Un QUBO dont les écarts d'énergie utiles sont minuscules devant le plus grand coefficient (pénalités trop fortes, quasi-dégénérescences comme la nôtre) est difficile à résoudre.
\end{itemize}
Le SDK \emph{Ocean} de D-Wave fournit aussi un recuit simulé classique (\texttt{SimulatedAnnealingSampler}), avec la même interface que le matériel : c'est ce que la note d'intention prévoit d'utiliser.

\exercices
\begin{exercice}
Au cours d'un recuit simulé, on propose un mouvement qui augmente l'énergie de $\Delta E=1$. Quelle est la probabilité de l'accepter à $T=2$ ? À $T=0{,}5$ ? À $T=0{,}1$ ? Expliquer le rôle de la température.
\end{exercice}
\begin{exercice}
Pour une famille de problèmes, on mesure $g_{\min}\approx2^{-n/2}$. Comment croît le temps de recuit adiabatique nécessaire avec $n$ ? Comparer au coût de l'énumération exhaustive ($2^n$ évaluations). Qu'en conclure ?
\end{exercice}

\chapter{QAOA : l'algorithme d'optimisation approximative quantique}\label{ch:qaoa}
\iterbanner{V0}

QAOA (\emph{Quantum Approximate Optimization Algorithm}, Farhi, Goldstone et Gutmann, 2014) est l'algorithme que la note d'intention prévoit d'exécuter sur simulateur. Il réunit tout ce qui précède : un hamiltonien de coût diagonal (ch.~\ref{ch:qubo}), un mélangeur qui ne commute pas avec lui (ch.~\ref{ch:algebre}), une alternance à la Trotter (ch.~\ref{ch:hamiltoniens}) et un principe variationnel réglé par un ordinateur classique.

\section{De l'adiabatique à QAOA}
Discrétisons le recuit adiabatique $H(s)=(1-s)H_B+sH_C$ en $p$ pas, et séparons chaque pas par Trotter :
\[e^{-iH(s_k)\Delta t}\approx e^{-i(1-s_k)\Delta t\,H_B}\ e^{-is_k\Delta t\,H_C}.\]
On obtient une alternance de deux types d'opérations, dont les « durées » suivent un programme fixé : $s_k\Delta t$ croissant pour le coût, $(1-s_k)\Delta t$ décroissant pour le mélangeur. \textbf{L'idée de QAOA : ne pas imposer ces durées, mais les optimiser.}

\section{Définition}
\begin{definition}[ (QAOA de profondeur $p$)]
Avec $H_C$ l'hamiltonien d'Ising du problème et le mélangeur $H_M=\sum_kX_k$, on prépare
\[\ket{\boldsymbol\gamma,\boldsymbol\beta}=e^{-i\beta_pH_M}e^{-i\gamma_pH_C}\cdots e^{-i\beta_1H_M}e^{-i\gamma_1H_C}\ \ket+^{\otimes n},\]
puis on cherche les $2p$ angles qui minimisent
\[F_p(\boldsymbol\gamma,\boldsymbol\beta)=\bra{\boldsymbol\gamma,\boldsymbol\beta}H_C\ket{\boldsymbol\gamma,\boldsymbol\beta}.\]
On mesure ensuite l'état optimal en base de calcul : chaque shot donne une chaîne de bits candidate.
\end{definition}
\paragraph{Ce que fait chaque étage, intuitivement.}
\begin{itemize}
\item $\ket+^{\otimes n}$ : toutes les chaînes ont la même amplitude, $2^{-n/2}$ ;
\item $e^{-i\gamma H_C}$ : chaque chaîne $x$ reçoit une phase $e^{-i\gamma E(x)}$. Le coût est « écrit » dans les phases, les probabilités sont inchangées ;
\item $e^{-i\beta H_M}$ : chaque qubit tourne autour de $x$. Les amplitudes de chaînes voisines (un bit de différence) se mélangent et \textbf{interfèrent}. Selon leurs phases, elles se renforcent ou s'annulent. Bien réglé, ce mélange concentre la probabilité sur les chaînes de basse énergie.
\end{itemize}
\paragraph{Garanties.} $F_p$ décroît (au sens large) avec $p$, et quand $p\to\infty$, QAOA contient le recuit adiabatique : il atteint l'optimum. À $p$ fini, c'est une heuristique.

\section{Le circuit}
On utilise $e^{-i\gamma h_iZ_i}=R_Z(2\gamma h_i)$, $e^{-i\gamma J_{ij}Z_iZ_j}=R_{ZZ}(2\gamma J_{ij})$ et $e^{-i\beta X_k}=R_X(2\beta)$. Tous les termes de $H_C$ commutent entre eux (ils sont tous diagonaux) : $e^{-i\gamma H_C}$ se factorise donc \emph{exactement}, sans erreur de Trotter.

\begin{figure}[htbp]
\centering
\begin{tikzpicture}[x=1cm,y=0.95cm]
\foreach \q in {0,1,2}{\draw[cink] (0,-\q) -- (11.6,-\q); \node[anchor=east,font=\small] at (0,-\q) {$q_\q:\ \ket0$};}
\foreach \q in {0,1,2}{\node[gate] at (0.6,-\q) {$H$};}
\draw[cmuted,dashed,rounded corners] (1.15,0.55) rectangle (7.55,-2.55);
\node[font=\scriptsize,cmuted] at (4.35,0.78) {coût $e^{-i\gamma_kH_C}$};
\node[gate,font=\footnotesize,minimum width=17mm] at (2.25,-1) {$R_Z(2\gamma_kh_1)$};
\node[draw=cink,fill=white,minimum width=20mm,minimum height=14mm,font=\footnotesize] at (4.3,-0.5) {$R_{ZZ}(2\gamma_kJ_{01})$};
\node[draw=cink,fill=white,minimum width=20mm,minimum height=14mm,font=\footnotesize] at (6.45,-1.5) {$R_{ZZ}(2\gamma_kJ_{12})$};
\draw[cmuted,dashed,rounded corners] (7.75,0.55) rectangle (9.65,-2.55);
\node[font=\scriptsize,cmuted] at (8.7,0.78) {mélangeur};
\foreach \q in {0,1,2}{\node[gate,font=\footnotesize,minimum width=14mm] at (8.7,-\q) {$R_X(2\beta_k)$};}
\draw[decorate,decoration={brace,amplitude=5pt,mirror}] (1.15,-2.75) -- (9.65,-2.75) node[midway,below=6pt,font=\small] {répété pour $k=1,\dots,p$};
\foreach \q in {0,1,2}{\meter{10.7}{-\q}}
\end{tikzpicture}
\caption{Circuit QAOA pour le QUBO jouet à 3 qubits ($h_0=h_2=0$, $h_1=-\frac12$, $J_{01}=J_{12}=\frac12$). Chaque $R_{ZZ}$ se décompose en 2 CNOT (figure~\ref{fig:rzz}).}\label{fig:circuit-qaoa}
\end{figure}

\paragraph{Ressources par couche.} Pour $n$ qubits et $|E|$ couplages non nuls : $n$ portes $R_Z$, $|E|$ portes $R_{ZZ}$ (donc $2|E|$ CNOT) et $n$ portes $R_X$. Au total, $p$ couches, $n$ portes $H$ initiales et $n$ mesures. Le nombre de qubits vaut $n$, le nombre de bits du QUBO.

\section{La boucle hybride quantique--classique}
\begin{figure}[htbp]
\centering
\begin{tikzpicture}[node distance=10mm and 16mm]
\node[flow,fill=cblue!10,minimum width=4.2cm,minimum height=1.5cm] (q) {\textbf{Processeur quantique}\\ (ou simulateur)\\ prépare $\ket{\boldsymbol\gamma,\boldsymbol\beta}$, mesure $N$ shots};
\node[flow,right=of q,minimum width=4.2cm,minimum height=1.5cm] (c) {\textbf{Optimiseur classique}\\ (COBYLA, Nelder-Mead)\\ propose de nouveaux $(\boldsymbol\gamma,\boldsymbol\beta)$};
\draw[arr] (q.north east) to[bend left=20] node[above,font=\small] {estimation de $F_p$} (c.north west);
\draw[arr] (c.south west) to[bend left=20] node[below,font=\small] {$(\boldsymbol\gamma,\boldsymbol\beta)$} (q.south east);
\node[flow,below=14mm of q,minimum width=4.2cm] (s) {échantillonnage final\\ $\to$ meilleure chaîne $x^*$};
\draw[arr] (q) -- node[left,font=\small] {à convergence} (s);
\end{tikzpicture}
\caption{QAOA est un algorithme \emph{variationnel} : le circuit quantique évalue la fonction, l'ordinateur classique l'optimise. Les $2p$ angles sont les seuls paramètres.}\label{fig:boucle-hybride}
\end{figure}

\section{Exemple entièrement calculé : un qubit}
Prenons $H_C=Z$ : les énergies sont $+1$ pour $\ket0$ et $-1$ pour $\ket1$, donc l'optimum est $\ket1$. Avec $p=1$, suivons le point sur la sphère de Bloch :
\begin{enumerate}
\item $\ket+$ est le point $(1,0,0)$ ;
\item $e^{-i\gamma Z}=R_Z(2\gamma)$ le fait tourner de $2\gamma$ autour de $z$ : $(\cos2\gamma,\ \sin2\gamma,\ 0)$ ;
\item $e^{-i\beta X}=R_X(2\beta)$ le fait tourner de $2\beta$ autour de $x$, ce qui donne une composante $z=\sin2\gamma\,\sin2\beta$.
\end{enumerate}
Comme $\expval Z$ est la coordonnée $z$ :
\[F_1(\gamma,\beta)=\sin2\gamma\ \sin2\beta,\qquad \min F_1=-1\ \text{ en }(\gamma,\beta)=(\tfrac\pi4,-\tfrac\pi4)\ \Rightarrow\ \text{on mesure }\ket1\text{ avec certitude}.\]
Le coût a d'abord fait tourner l'état \emph{dans l'équateur} (une phase), puis le mélangeur a converti cette phase en déplacement vers le pôle de basse énergie. C'est le mécanisme général, sur $2^n$ amplitudes au lieu d'une.

\section{Résultats sur le QUBO jouet (3 qubits)}
\begin{figure}[htbp]
\centering
% figure générée par tools/make_figures.py : paysage_qaoa_p1
\begin{tikzpicture}
\begin{axis}[width=0.78\linewidth,height=6cm,view={0}{90},xmin=0,xmax=6.2832,ymin=0,ymax=3.1416,
  xlabel={$\gamma$},ylabel={$\beta$},xtick={0,1.5708,3.1416,4.7124,6.2832},
  xticklabels={$0$,$\pi/2$,$\pi$,$3\pi/2$,$2\pi$},ytick={0,1.5708,3.1416},yticklabels={$0$,$\pi/2$,$\pi$},
  colormap={seq}{color=(cblue!85!black) color=(cblue!45) color=(white)},
  colorbar,colorbar style={ylabel={$F_1(\gamma,\beta)$}},axis on top,enlargelimits=false]
\addplot3[surf,shader=interp,mesh/rows=37] coordinates {
(0.0000,0.0000,-0.5000) (0.0873,0.0000,-0.5000) (0.1745,0.0000,-0.5000) (0.2618,0.0000,-0.5000) (0.3491,0.0000,-0.5000) (0.4363,0.0000,-0.5000) (0.5236,0.0000,-0.5000) (0.6109,0.0000,-0.5000) (0.6981,0.0000,-0.5000) (0.7854,0.0000,-0.5000) (0.8727,0.0000,-0.5000) (0.9599,0.0000,-0.5000) (1.0472,0.0000,-0.5000) (1.1345,0.0000,-0.5000) (1.2217,0.0000,-0.5000) (1.3090,0.0000,-0.5000) (1.3963,0.0000,-0.5000) (1.4835,0.0000,-0.5000) (1.5708,0.0000,-0.5000) (1.6581,0.0000,-0.5000) (1.7453,0.0000,-0.5000) (1.8326,0.0000,-0.5000) (1.9199,0.0000,-0.5000) (2.0071,0.0000,-0.5000) (2.0944,0.0000,-0.5000) (2.1817,0.0000,-0.5000) (2.2689,0.0000,-0.5000) (2.3562,0.0000,-0.5000) (2.4435,0.0000,-0.5000) (2.5307,0.0000,-0.5000) (2.6180,0.0000,-0.5000) (2.7053,0.0000,-0.5000) (2.7925,0.0000,-0.5000) (2.8798,0.0000,-0.5000) (2.9671,0.0000,-0.5000) (3.0543,0.0000,-0.5000) (3.1416,0.0000,-0.5000) (3.2289,0.0000,-0.5000) (3.3161,0.0000,-0.5000) (3.4034,0.0000,-0.5000) (3.4907,0.0000,-0.5000) (3.5779,0.0000,-0.5000) (3.6652,0.0000,-0.5000) (3.7525,0.0000,-0.5000) (3.8397,0.0000,-0.5000) (3.9270,0.0000,-0.5000) (4.0143,0.0000,-0.5000) (4.1015,0.0000,-0.5000) (4.1888,0.0000,-0.5000) (4.2761,0.0000,-0.5000) (4.3633,0.0000,-0.5000) (4.4506,0.0000,-0.5000) (4.5379,0.0000,-0.5000) (4.6251,0.0000,-0.5000) (4.7124,0.0000,-0.5000) (4.7997,0.0000,-0.5000) (4.8869,0.0000,-0.5000) (4.9742,0.0000,-0.5000) (5.0615,0.0000,-0.5000) (5.1487,0.0000,-0.5000) (5.2360,0.0000,-0.5000) (5.3233,0.0000,-0.5000) (5.4105,0.0000,-0.5000) (5.4978,0.0000,-0.5000) (5.5851,0.0000,-0.5000) (5.6723,0.0000,-0.5000) (5.7596,0.0000,-0.5000) (5.8469,0.0000,-0.5000) (5.9341,0.0000,-0.5000) (6.0214,0.0000,-0.5000) (6.1087,0.0000,-0.5000) (6.1959,0.0000,-0.5000) (6.2832,0.0000,-0.5000)
(0.0000,0.0873,-0.5000) (0.0873,0.0873,-0.4628) (0.1745,0.0873,-0.4269) (0.2618,0.0873,-0.3935) (0.3491,0.0873,-0.3636) (0.4363,0.0873,-0.3382) (0.5236,0.0873,-0.3178) (0.6109,0.0873,-0.3027) (0.6981,0.0873,-0.2928) (0.7854,0.0873,-0.2879) (0.8727,0.0873,-0.2874) (0.9599,0.0873,-0.2904) (1.0472,0.0873,-0.2961) (1.1345,0.0873,-0.3033) (1.2217,0.0873,-0.3110) (1.3090,0.0873,-0.3181) (1.3963,0.0873,-0.3239) (1.4835,0.0873,-0.3277) (1.5708,0.0873,-0.3290) (1.6581,0.0873,-0.3277) (1.7453,0.0873,-0.3239) (1.8326,0.0873,-0.3181) (1.9199,0.0873,-0.3110) (2.0071,0.0873,-0.3033) (2.0944,0.0873,-0.2961) (2.1817,0.0873,-0.2904) (2.2689,0.0873,-0.2874) (2.3562,0.0873,-0.2879) (2.4435,0.0873,-0.2928) (2.5307,0.0873,-0.3027) (2.6180,0.0873,-0.3178) (2.7053,0.0873,-0.3382) (2.7925,0.0873,-0.3636) (2.8798,0.0873,-0.3935) (2.9671,0.0873,-0.4269) (3.0543,0.0873,-0.4628) (3.1416,0.0873,-0.5000) (3.2289,0.0873,-0.5372) (3.3161,0.0873,-0.5731) (3.4034,0.0873,-0.6065) (3.4907,0.0873,-0.6364) (3.5779,0.0873,-0.6618) (3.6652,0.0873,-0.6822) (3.7525,0.0873,-0.6973) (3.8397,0.0873,-0.7072) (3.9270,0.0873,-0.7121) (4.0143,0.0873,-0.7126) (4.1015,0.0873,-0.7096) (4.1888,0.0873,-0.7039) (4.2761,0.0873,-0.6967) (4.3633,0.0873,-0.6890) (4.4506,0.0873,-0.6819) (4.5379,0.0873,-0.6761) (4.6251,0.0873,-0.6723) (4.7124,0.0873,-0.6710) (4.7997,0.0873,-0.6723) (4.8869,0.0873,-0.6761) (4.9742,0.0873,-0.6819) (5.0615,0.0873,-0.6890) (5.1487,0.0873,-0.6967) (5.2360,0.0873,-0.7039) (5.3233,0.0873,-0.7096) (5.4105,0.0873,-0.7126) (5.4978,0.0873,-0.7121) (5.5851,0.0873,-0.7072) (5.6723,0.0873,-0.6973) (5.7596,0.0873,-0.6822) (5.8469,0.0873,-0.6618) (5.9341,0.0873,-0.6364) (6.0214,0.0873,-0.6065) (6.1087,0.0873,-0.5731) (6.1959,0.0873,-0.5372) (6.2832,0.0873,-0.5000)
(0.0000,0.1745,-0.5000) (0.0873,0.1745,-0.4294) (0.1745,0.1745,-0.3613) (0.2618,0.1745,-0.2979) (0.3491,0.1745,-0.2414) (0.4363,0.1745,-0.1932) (0.5236,0.1745,-0.1547) (0.6109,0.1745,-0.1261) (0.6981,0.1745,-0.1077) (0.7854,0.1745,-0.0986) (0.8727,0.1745,-0.0979) (0.9599,0.1745,-0.1040) (1.0472,0.1745,-0.1151) (1.1345,0.1745,-0.1290) (1.2217,0.1745,-0.1439) (1.3090,0.1745,-0.1577) (1.3963,0.1745,-0.1689) (1.4835,0.1745,-0.1761) (1.5708,0.1745,-0.1786) (1.6581,0.1745,-0.1761) (1.7453,0.1745,-0.1689) (1.8326,0.1745,-0.1577) (1.9199,0.1745,-0.1439) (2.0071,0.1745,-0.1290) (2.0944,0.1745,-0.1151) (2.1817,0.1745,-0.1040) (2.2689,0.1745,-0.0979) (2.3562,0.1745,-0.0986) (2.4435,0.1745,-0.1077) (2.5307,0.1745,-0.1261) (2.6180,0.1745,-0.1547) (2.7053,0.1745,-0.1932) (2.7925,0.1745,-0.2414) (2.8798,0.1745,-0.2979) (2.9671,0.1745,-0.3613) (3.0543,0.1745,-0.4294) (3.1416,0.1745,-0.5000) (3.2289,0.1745,-0.5706) (3.3161,0.1745,-0.6387) (3.4034,0.1745,-0.7021) (3.4907,0.1745,-0.7586) (3.5779,0.1745,-0.8068) (3.6652,0.1745,-0.8453) (3.7525,0.1745,-0.8739) (3.8397,0.1745,-0.8923) (3.9270,0.1745,-0.9014) (4.0143,0.1745,-0.9021) (4.1015,0.1745,-0.8960) (4.1888,0.1745,-0.8849) (4.2761,0.1745,-0.8710) (4.3633,0.1745,-0.8561) (4.4506,0.1745,-0.8423) (4.5379,0.1745,-0.8311) (4.6251,0.1745,-0.8239) (4.7124,0.1745,-0.8214) (4.7997,0.1745,-0.8239) (4.8869,0.1745,-0.8311) (4.9742,0.1745,-0.8423) (5.0615,0.1745,-0.8561) (5.1487,0.1745,-0.8710) (5.2360,0.1745,-0.8849) (5.3233,0.1745,-0.8960) (5.4105,0.1745,-0.9021) (5.4978,0.1745,-0.9014) (5.5851,0.1745,-0.8923) (5.6723,0.1745,-0.8739) (5.7596,0.1745,-0.8453) (5.8469,0.1745,-0.8068) (5.9341,0.1745,-0.7586) (6.0214,0.1745,-0.7021) (6.1087,0.1745,-0.6387) (6.1959,0.1745,-0.5706) (6.2832,0.1745,-0.5000)
(0.0000,0.2618,-0.5000) (0.0873,0.2618,-0.4032) (0.1745,0.2618,-0.3098) (0.2618,0.2618,-0.2230) (0.3491,0.2618,-0.1456) (0.4363,0.2618,-0.0799) (0.5236,0.2618,-0.0274) (0.6109,0.2618,0.0112) (0.6981,0.2618,0.0360) (0.7854,0.2618,0.0477) (0.8727,0.2618,0.0479) (0.9599,0.2618,0.0388) (1.0472,0.2618,0.0229) (1.1345,0.2618,0.0030) (1.2217,0.2618,-0.0180) (1.3090,0.2618,-0.0375) (1.3963,0.2618,-0.0533) (1.4835,0.2618,-0.0635) (1.5708,0.2618,-0.0670) (1.6581,0.2618,-0.0635) (1.7453,0.2618,-0.0533) (1.8326,0.2618,-0.0375) (1.9199,0.2618,-0.0180) (2.0071,0.2618,0.0030) (2.0944,0.2618,0.0229) (2.1817,0.2618,0.0388) (2.2689,0.2618,0.0479) (2.3562,0.2618,0.0477) (2.4435,0.2618,0.0360) (2.5307,0.2618,0.0112) (2.6180,0.2618,-0.0274) (2.7053,0.2618,-0.0799) (2.7925,0.2618,-0.1456) (2.8798,0.2618,-0.2230) (2.9671,0.2618,-0.3098) (3.0543,0.2618,-0.4032) (3.1416,0.2618,-0.5000) (3.2289,0.2618,-0.5968) (3.3161,0.2618,-0.6902) (3.4034,0.2618,-0.7770) (3.4907,0.2618,-0.8544) (3.5779,0.2618,-0.9201) (3.6652,0.2618,-0.9726) (3.7525,0.2618,-1.0112) (3.8397,0.2618,-1.0360) (3.9270,0.2618,-1.0477) (4.0143,0.2618,-1.0479) (4.1015,0.2618,-1.0388) (4.1888,0.2618,-1.0229) (4.2761,0.2618,-1.0030) (4.3633,0.2618,-0.9820) (4.4506,0.2618,-0.9625) (4.5379,0.2618,-0.9467) (4.6251,0.2618,-0.9365) (4.7124,0.2618,-0.9330) (4.7997,0.2618,-0.9365) (4.8869,0.2618,-0.9467) (4.9742,0.2618,-0.9625) (5.0615,0.2618,-0.9820) (5.1487,0.2618,-1.0030) (5.2360,0.2618,-1.0229) (5.3233,0.2618,-1.0388) (5.4105,0.2618,-1.0479) (5.4978,0.2618,-1.0477) (5.5851,0.2618,-1.0360) (5.6723,0.2618,-1.0112) (5.7596,0.2618,-0.9726) (5.8469,0.2618,-0.9201) (5.9341,0.2618,-0.8544) (6.0214,0.2618,-0.7770) (6.1087,0.2618,-0.6902) (6.1959,0.2618,-0.5968) (6.2832,0.2618,-0.5000)
(0.0000,0.3491,-0.5000) (0.0873,0.3491,-0.3867) (0.1745,0.3491,-0.2774) (0.2618,0.3491,-0.1760) (0.3491,0.3491,-0.0858) (0.4363,0.3491,-0.0094) (0.5236,0.3491,0.0514) (0.6109,0.3491,0.0956) (0.6981,0.3491,0.1235) (0.7854,0.3491,0.1359) (0.8727,0.3491,0.1348) (0.9599,0.3491,0.1227) (1.0472,0.3491,0.1026) (1.1345,0.3491,0.0780) (1.2217,0.3491,0.0522) (1.3090,0.3491,0.0283) (1.3963,0.3491,0.0091) (1.4835,0.3491,-0.0033) (1.5708,0.3491,-0.0076) (1.6581,0.3491,-0.0033) (1.7453,0.3491,0.0091) (1.8326,0.3491,0.0283) (1.9199,0.3491,0.0522) (2.0071,0.3491,0.0780) (2.0944,0.3491,0.1026) (2.1817,0.3491,0.1227) (2.2689,0.3491,0.1348) (2.3562,0.3491,0.1359) (2.4435,0.3491,0.1235) (2.5307,0.3491,0.0956) (2.6180,0.3491,0.0514) (2.7053,0.3491,-0.0094) (2.7925,0.3491,-0.0858) (2.8798,0.3491,-0.1760) (2.9671,0.3491,-0.2774) (3.0543,0.3491,-0.3867) (3.1416,0.3491,-0.5000) (3.2289,0.3491,-0.6133) (3.3161,0.3491,-0.7226) (3.4034,0.3491,-0.8240) (3.4907,0.3491,-0.9142) (3.5779,0.3491,-0.9906) (3.6652,0.3491,-1.0514) (3.7525,0.3491,-1.0956) (3.8397,0.3491,-1.1235) (3.9270,0.3491,-1.1359) (4.0143,0.3491,-1.1348) (4.1015,0.3491,-1.1227) (4.1888,0.3491,-1.1026) (4.2761,0.3491,-1.0780) (4.3633,0.3491,-1.0522) (4.4506,0.3491,-1.0283) (4.5379,0.3491,-1.0091) (4.6251,0.3491,-0.9967) (4.7124,0.3491,-0.9924) (4.7997,0.3491,-0.9967) (4.8869,0.3491,-1.0091) (4.9742,0.3491,-1.0283) (5.0615,0.3491,-1.0522) (5.1487,0.3491,-1.0780) (5.2360,0.3491,-1.1026) (5.3233,0.3491,-1.1227) (5.4105,0.3491,-1.1348) (5.4978,0.3491,-1.1359) (5.5851,0.3491,-1.1235) (5.6723,0.3491,-1.0956) (5.7596,0.3491,-1.0514) (5.8469,0.3491,-0.9906) (5.9341,0.3491,-0.9142) (6.0214,0.3491,-0.8240) (6.1087,0.3491,-0.7226) (6.1959,0.3491,-0.6133) (6.2832,0.3491,-0.5000)
(0.0000,0.4363,-0.5000) (0.0873,0.4363,-0.3814) (0.1745,0.4363,-0.2671) (0.2618,0.4363,-0.1612) (0.3491,0.4363,-0.0672) (0.4363,0.4363,0.0120) (0.5236,0.4363,0.0745) (0.6109,0.4363,0.1194) (0.6981,0.4363,0.1467) (0.7854,0.4363,0.1577) (0.8727,0.4363,0.1543) (0.9599,0.4363,0.1393) (1.0472,0.4363,0.1160) (1.1345,0.4363,0.0880) (1.2217,0.4363,0.0589) (1.3090,0.4363,0.0323) (1.3963,0.4363,0.0109) (1.4835,0.4363,-0.0028) (1.5708,0.4363,-0.0076) (1.6581,0.4363,-0.0028) (1.7453,0.4363,0.0109) (1.8326,0.4363,0.0323) (1.9199,0.4363,0.0589) (2.0071,0.4363,0.0880) (2.0944,0.4363,0.1160) (2.1817,0.4363,0.1393) (2.2689,0.4363,0.1543) (2.3562,0.4363,0.1577) (2.4435,0.4363,0.1467) (2.5307,0.4363,0.1194) (2.6180,0.4363,0.0745) (2.7053,0.4363,0.0120) (2.7925,0.4363,-0.0672) (2.8798,0.4363,-0.1612) (2.9671,0.4363,-0.2671) (3.0543,0.4363,-0.3814) (3.1416,0.4363,-0.5000) (3.2289,0.4363,-0.6186) (3.3161,0.4363,-0.7329) (3.4034,0.4363,-0.8388) (3.4907,0.4363,-0.9328) (3.5779,0.4363,-1.0120) (3.6652,0.4363,-1.0745) (3.7525,0.4363,-1.1194) (3.8397,0.4363,-1.1467) (3.9270,0.4363,-1.1577) (4.0143,0.4363,-1.1543) (4.1015,0.4363,-1.1393) (4.1888,0.4363,-1.1160) (4.2761,0.4363,-1.0880) (4.3633,0.4363,-1.0589) (4.4506,0.4363,-1.0323) (4.5379,0.4363,-1.0109) (4.6251,0.4363,-0.9972) (4.7124,0.4363,-0.9924) (4.7997,0.4363,-0.9972) (4.8869,0.4363,-1.0109) (4.9742,0.4363,-1.0323) (5.0615,0.4363,-1.0589) (5.1487,0.4363,-1.0880) (5.2360,0.4363,-1.1160) (5.3233,0.4363,-1.1393) (5.4105,0.4363,-1.1543) (5.4978,0.4363,-1.1577) (5.5851,0.4363,-1.1467) (5.6723,0.4363,-1.1194) (5.7596,0.4363,-1.0745) (5.8469,0.4363,-1.0120) (5.9341,0.4363,-0.9328) (6.0214,0.4363,-0.8388) (6.1087,0.4363,-0.7329) (6.1959,0.4363,-0.6186) (6.2832,0.4363,-0.5000)
(0.0000,0.5236,-0.5000) (0.0873,0.5236,-0.3874) (0.1745,0.5236,-0.2790) (0.2618,0.5236,-0.1788) (0.3491,0.5236,-0.0904) (0.4363,0.5236,-0.0164) (0.5236,0.5236,0.0413) (0.6109,0.5236,0.0817) (0.6981,0.5236,0.1050) (0.7854,0.5236,0.1124) (0.8727,0.5236,0.1058) (0.9599,0.5236,0.0881) (1.0472,0.5236,0.0625) (1.1345,0.5236,0.0326) (1.2217,0.5236,0.0021) (1.3090,0.5236,-0.0257) (1.3963,0.5236,-0.0478) (1.4835,0.5236,-0.0621) (1.5708,0.5236,-0.0670) (1.6581,0.5236,-0.0621) (1.7453,0.5236,-0.0478) (1.8326,0.5236,-0.0257) (1.9199,0.5236,0.0021) (2.0071,0.5236,0.0326) (2.0944,0.5236,0.0625) (2.1817,0.5236,0.0881) (2.2689,0.5236,0.1058) (2.3562,0.5236,0.1124) (2.4435,0.5236,0.1050) (2.5307,0.5236,0.0817) (2.6180,0.5236,0.0413) (2.7053,0.5236,-0.0164) (2.7925,0.5236,-0.0904) (2.8798,0.5236,-0.1788) (2.9671,0.5236,-0.2790) (3.0543,0.5236,-0.3874) (3.1416,0.5236,-0.5000) (3.2289,0.5236,-0.6126) (3.3161,0.5236,-0.7210) (3.4034,0.5236,-0.8212) (3.4907,0.5236,-0.9096) (3.5779,0.5236,-0.9836) (3.6652,0.5236,-1.0413) (3.7525,0.5236,-1.0817) (3.8397,0.5236,-1.1050) (3.9270,0.5236,-1.1124) (4.0143,0.5236,-1.1058) (4.1015,0.5236,-1.0881) (4.1888,0.5236,-1.0625) (4.2761,0.5236,-1.0326) (4.3633,0.5236,-1.0021) (4.4506,0.5236,-0.9743) (4.5379,0.5236,-0.9522) (4.6251,0.5236,-0.9379) (4.7124,0.5236,-0.9330) (4.7997,0.5236,-0.9379) (4.8869,0.5236,-0.9522) (4.9742,0.5236,-0.9743) (5.0615,0.5236,-1.0021) (5.1487,0.5236,-1.0326) (5.2360,0.5236,-1.0625) (5.3233,0.5236,-1.0881) (5.4105,0.5236,-1.1058) (5.4978,0.5236,-1.1124) (5.5851,0.5236,-1.1050) (5.6723,0.5236,-1.0817) (5.7596,0.5236,-1.0413) (5.8469,0.5236,-0.9836) (5.9341,0.5236,-0.9096) (6.0214,0.5236,-0.8212) (6.1087,0.5236,-0.7210) (6.1959,0.5236,-0.6126) (6.2832,0.5236,-0.5000)
(0.0000,0.6109,-0.5000) (0.0873,0.6109,-0.4036) (0.1745,0.6109,-0.3109) (0.2618,0.6109,-0.2257) (0.3491,0.6109,-0.1511) (0.4363,0.6109,-0.0895) (0.5236,0.6109,-0.0426) (0.6109,0.6109,-0.0111) (0.6981,0.6109,0.0050) (0.7854,0.6109,0.0070) (0.8727,0.6109,-0.0034) (0.9599,0.6109,-0.0235) (1.0472,0.6109,-0.0504) (1.1345,0.6109,-0.0806) (1.2217,0.6109,-0.1110) (1.3090,0.6109,-0.1384) (1.3963,0.6109,-0.1600) (1.4835,0.6109,-0.1738) (1.5708,0.6109,-0.1786) (1.6581,0.6109,-0.1738) (1.7453,0.6109,-0.1600) (1.8326,0.6109,-0.1384) (1.9199,0.6109,-0.1110) (2.0071,0.6109,-0.0806) (2.0944,0.6109,-0.0504) (2.1817,0.6109,-0.0235) (2.2689,0.6109,-0.0034) (2.3562,0.6109,0.0070) (2.4435,0.6109,0.0050) (2.5307,0.6109,-0.0111) (2.6180,0.6109,-0.0426) (2.7053,0.6109,-0.0895) (2.7925,0.6109,-0.1511) (2.8798,0.6109,-0.2257) (2.9671,0.6109,-0.3109) (3.0543,0.6109,-0.4036) (3.1416,0.6109,-0.5000) (3.2289,0.6109,-0.5964) (3.3161,0.6109,-0.6891) (3.4034,0.6109,-0.7743) (3.4907,0.6109,-0.8489) (3.5779,0.6109,-0.9105) (3.6652,0.6109,-0.9574) (3.7525,0.6109,-0.9889) (3.8397,0.6109,-1.0050) (3.9270,0.6109,-1.0070) (4.0143,0.6109,-0.9966) (4.1015,0.6109,-0.9765) (4.1888,0.6109,-0.9496) (4.2761,0.6109,-0.9194) (4.3633,0.6109,-0.8890) (4.4506,0.6109,-0.8616) (4.5379,0.6109,-0.8400) (4.6251,0.6109,-0.8262) (4.7124,0.6109,-0.8214) (4.7997,0.6109,-0.8262) (4.8869,0.6109,-0.8400) (4.9742,0.6109,-0.8616) (5.0615,0.6109,-0.8890) (5.1487,0.6109,-0.9194) (5.2360,0.6109,-0.9496) (5.3233,0.6109,-0.9765) (5.4105,0.6109,-0.9966) (5.4978,0.6109,-1.0070) (5.5851,0.6109,-1.0050) (5.6723,0.6109,-0.9889) (5.7596,0.6109,-0.9574) (5.8469,0.6109,-0.9105) (5.9341,0.6109,-0.8489) (6.0214,0.6109,-0.7743) (6.1087,0.6109,-0.6891) (6.1959,0.6109,-0.5964) (6.2832,0.6109,-0.5000)
(0.0000,0.6981,-0.5000) (0.0873,0.6981,-0.4277) (0.1745,0.6981,-0.3586) (0.2618,0.6981,-0.2955) (0.3491,0.6981,-0.2412) (0.4363,0.6981,-0.1974) (0.5236,0.6981,-0.1657) (0.6109,0.6981,-0.1466) (0.6981,0.6981,-0.1398) (0.7854,0.6981,-0.1445) (0.8727,0.6981,-0.1590) (0.9599,0.6981,-0.1811) (1.0472,0.6981,-0.2083) (1.1345,0.6981,-0.2376) (1.2217,0.6981,-0.2664) (1.3090,0.6981,-0.2919) (1.3963,0.6981,-0.3119) (1.4835,0.6981,-0.3246) (1.5708,0.6981,-0.3290) (1.6581,0.6981,-0.3246) (1.7453,0.6981,-0.3119) (1.8326,0.6981,-0.2919) (1.9199,0.6981,-0.2664) (2.0071,0.6981,-0.2376) (2.0944,0.6981,-0.2083) (2.1817,0.6981,-0.1811) (2.2689,0.6981,-0.1590) (2.3562,0.6981,-0.1445) (2.4435,0.6981,-0.1398) (2.5307,0.6981,-0.1466) (2.6180,0.6981,-0.1657) (2.7053,0.6981,-0.1974) (2.7925,0.6981,-0.2412) (2.8798,0.6981,-0.2955) (2.9671,0.6981,-0.3586) (3.0543,0.6981,-0.4277) (3.1416,0.6981,-0.5000) (3.2289,0.6981,-0.5723) (3.3161,0.6981,-0.6414) (3.4034,0.6981,-0.7045) (3.4907,0.6981,-0.7588) (3.5779,0.6981,-0.8026) (3.6652,0.6981,-0.8343) (3.7525,0.6981,-0.8534) (3.8397,0.6981,-0.8602) (3.9270,0.6981,-0.8555) (4.0143,0.6981,-0.8410) (4.1015,0.6981,-0.8189) (4.1888,0.6981,-0.7917) (4.2761,0.6981,-0.7624) (4.3633,0.6981,-0.7336) (4.4506,0.6981,-0.7081) (4.5379,0.6981,-0.6881) (4.6251,0.6981,-0.6754) (4.7124,0.6981,-0.6710) (4.7997,0.6981,-0.6754) (4.8869,0.6981,-0.6881) (4.9742,0.6981,-0.7081) (5.0615,0.6981,-0.7336) (5.1487,0.6981,-0.7624) (5.2360,0.6981,-0.7917) (5.3233,0.6981,-0.8189) (5.4105,0.6981,-0.8410) (5.4978,0.6981,-0.8555) (5.5851,0.6981,-0.8602) (5.6723,0.6981,-0.8534) (5.7596,0.6981,-0.8343) (5.8469,0.6981,-0.8026) (5.9341,0.6981,-0.7588) (6.0214,0.6981,-0.7045) (6.1087,0.6981,-0.6414) (6.1959,0.6981,-0.5723) (6.2832,0.6981,-0.5000)
(0.0000,0.7854,-0.5000) (0.0873,0.7854,-0.4568) (0.1745,0.7854,-0.4158) (0.2618,0.7854,-0.3793) (0.3491,0.7854,-0.3490) (0.4363,0.7854,-0.3264) (0.5236,0.7854,-0.3125) (0.6109,0.7854,-0.3076) (0.6981,0.7854,-0.3114) (0.7854,0.7854,-0.3232) (0.8727,0.7854,-0.3417) (0.9599,0.7854,-0.3653) (1.0472,0.7854,-0.3917) (1.1345,0.7854,-0.4191) (1.2217,0.7854,-0.4450) (1.3090,0.7854,-0.4676) (1.3963,0.7854,-0.4852) (1.4835,0.7854,-0.4962) (1.5708,0.7854,-0.5000) (1.6581,0.7854,-0.4962) (1.7453,0.7854,-0.4852) (1.8326,0.7854,-0.4676) (1.9199,0.7854,-0.4450) (2.0071,0.7854,-0.4191) (2.0944,0.7854,-0.3917) (2.1817,0.7854,-0.3653) (2.2689,0.7854,-0.3417) (2.3562,0.7854,-0.3232) (2.4435,0.7854,-0.3114) (2.5307,0.7854,-0.3076) (2.6180,0.7854,-0.3125) (2.7053,0.7854,-0.3264) (2.7925,0.7854,-0.3490) (2.8798,0.7854,-0.3793) (2.9671,0.7854,-0.4158) (3.0543,0.7854,-0.4568) (3.1416,0.7854,-0.5000) (3.2289,0.7854,-0.5432) (3.3161,0.7854,-0.5842) (3.4034,0.7854,-0.6207) (3.4907,0.7854,-0.6510) (3.5779,0.7854,-0.6736) (3.6652,0.7854,-0.6875) (3.7525,0.7854,-0.6924) (3.8397,0.7854,-0.6886) (3.9270,0.7854,-0.6768) (4.0143,0.7854,-0.6583) (4.1015,0.7854,-0.6347) (4.1888,0.7854,-0.6083) (4.2761,0.7854,-0.5809) (4.3633,0.7854,-0.5550) (4.4506,0.7854,-0.5324) (4.5379,0.7854,-0.5148) (4.6251,0.7854,-0.5038) (4.7124,0.7854,-0.5000) (4.7997,0.7854,-0.5038) (4.8869,0.7854,-0.5148) (4.9742,0.7854,-0.5324) (5.0615,0.7854,-0.5550) (5.1487,0.7854,-0.5809) (5.2360,0.7854,-0.6083) (5.3233,0.7854,-0.6347) (5.4105,0.7854,-0.6583) (5.4978,0.7854,-0.6768) (5.5851,0.7854,-0.6886) (5.6723,0.7854,-0.6924) (5.7596,0.7854,-0.6875) (5.8469,0.7854,-0.6736) (5.9341,0.7854,-0.6510) (6.0214,0.7854,-0.6207) (6.1087,0.7854,-0.5842) (6.1959,0.7854,-0.5432) (6.2832,0.7854,-0.5000)
(0.0000,0.8727,-0.5000) (0.0873,0.8727,-0.4871) (0.1745,0.8727,-0.4756) (0.2618,0.8727,-0.4667) (0.3491,0.8727,-0.4614) (0.4363,0.8727,-0.4607) (0.5236,0.8727,-0.4650) (0.6109,0.8727,-0.4744) (0.6981,0.8727,-0.4887) (0.7854,0.8727,-0.5073) (0.8727,0.8727,-0.5293) (0.9599,0.8727,-0.5535) (1.0472,0.8727,-0.5785) (1.1345,0.8727,-0.6030) (1.2217,0.8727,-0.6254) (1.3090,0.8727,-0.6444) (1.3963,0.8727,-0.6589) (1.4835,0.8727,-0.6679) (1.5708,0.8727,-0.6710) (1.6581,0.8727,-0.6679) (1.7453,0.8727,-0.6589) (1.8326,0.8727,-0.6444) (1.9199,0.8727,-0.6254) (2.0071,0.8727,-0.6030) (2.0944,0.8727,-0.5785) (2.1817,0.8727,-0.5535) (2.2689,0.8727,-0.5293) (2.3562,0.8727,-0.5073) (2.4435,0.8727,-0.4887) (2.5307,0.8727,-0.4744) (2.6180,0.8727,-0.4650) (2.7053,0.8727,-0.4607) (2.7925,0.8727,-0.4614) (2.8798,0.8727,-0.4667) (2.9671,0.8727,-0.4756) (3.0543,0.8727,-0.4871) (3.1416,0.8727,-0.5000) (3.2289,0.8727,-0.5129) (3.3161,0.8727,-0.5244) (3.4034,0.8727,-0.5333) (3.4907,0.8727,-0.5386) (3.5779,0.8727,-0.5393) (3.6652,0.8727,-0.5350) (3.7525,0.8727,-0.5256) (3.8397,0.8727,-0.5113) (3.9270,0.8727,-0.4927) (4.0143,0.8727,-0.4707) (4.1015,0.8727,-0.4465) (4.1888,0.8727,-0.4215) (4.2761,0.8727,-0.3970) (4.3633,0.8727,-0.3746) (4.4506,0.8727,-0.3556) (4.5379,0.8727,-0.3411) (4.6251,0.8727,-0.3321) (4.7124,0.8727,-0.3290) (4.7997,0.8727,-0.3321) (4.8869,0.8727,-0.3411) (4.9742,0.8727,-0.3556) (5.0615,0.8727,-0.3746) (5.1487,0.8727,-0.3970) (5.2360,0.8727,-0.4215) (5.3233,0.8727,-0.4465) (5.4105,0.8727,-0.4707) (5.4978,0.8727,-0.4927) (5.5851,0.8727,-0.5113) (5.6723,0.8727,-0.5256) (5.7596,0.8727,-0.5350) (5.8469,0.8727,-0.5393) (5.9341,0.8727,-0.5386) (6.0214,0.8727,-0.5333) (6.1087,0.8727,-0.5244) (6.1959,0.8727,-0.5129) (6.2832,0.8727,-0.5000)
(0.0000,0.9599,-0.5000) (0.0873,0.9599,-0.5152) (0.1745,0.9599,-0.5308) (0.2618,0.9599,-0.5473) (0.3491,0.9599,-0.5651) (0.4363,0.9599,-0.5843) (0.5236,0.9599,-0.6050) (0.6109,0.9599,-0.6272) (0.6981,0.9599,-0.6506) (0.7854,0.9599,-0.6748) (0.8727,0.9599,-0.6992) (0.9599,0.9599,-0.7233) (1.0472,0.9599,-0.7462) (1.1345,0.9599,-0.7673) (1.2217,0.9599,-0.7857) (1.3090,0.9599,-0.8008) (1.3963,0.9599,-0.8121) (1.4835,0.9599,-0.8190) (1.5708,0.9599,-0.8214) (1.6581,0.9599,-0.8190) (1.7453,0.9599,-0.8121) (1.8326,0.9599,-0.8008) (1.9199,0.9599,-0.7857) (2.0071,0.9599,-0.7673) (2.0944,0.9599,-0.7462) (2.1817,0.9599,-0.7233) (2.2689,0.9599,-0.6992) (2.3562,0.9599,-0.6748) (2.4435,0.9599,-0.6506) (2.5307,0.9599,-0.6272) (2.6180,0.9599,-0.6050) (2.7053,0.9599,-0.5843) (2.7925,0.9599,-0.5651) (2.8798,0.9599,-0.5473) (2.9671,0.9599,-0.5308) (3.0543,0.9599,-0.5152) (3.1416,0.9599,-0.5000) (3.2289,0.9599,-0.4848) (3.3161,0.9599,-0.4692) (3.4034,0.9599,-0.4527) (3.4907,0.9599,-0.4349) (3.5779,0.9599,-0.4157) (3.6652,0.9599,-0.3950) (3.7525,0.9599,-0.3728) (3.8397,0.9599,-0.3494) (3.9270,0.9599,-0.3252) (4.0143,0.9599,-0.3008) (4.1015,0.9599,-0.2767) (4.1888,0.9599,-0.2538) (4.2761,0.9599,-0.2327) (4.3633,0.9599,-0.2143) (4.4506,0.9599,-0.1992) (4.5379,0.9599,-0.1879) (4.6251,0.9599,-0.1810) (4.7124,0.9599,-0.1786) (4.7997,0.9599,-0.1810) (4.8869,0.9599,-0.1879) (4.9742,0.9599,-0.1992) (5.0615,0.9599,-0.2143) (5.1487,0.9599,-0.2327) (5.2360,0.9599,-0.2538) (5.3233,0.9599,-0.2767) (5.4105,0.9599,-0.3008) (5.4978,0.9599,-0.3252) (5.5851,0.9599,-0.3494) (5.6723,0.9599,-0.3728) (5.7596,0.9599,-0.3950) (5.8469,0.9599,-0.4157) (5.9341,0.9599,-0.4349) (6.0214,0.9599,-0.4527) (6.1087,0.9599,-0.4692) (6.1959,0.9599,-0.4848) (6.2832,0.9599,-0.5000)
(0.0000,1.0472,-0.5000) (0.0873,1.0472,-0.5377) (0.1745,1.0472,-0.5752) (0.2618,1.0472,-0.6121) (0.3491,1.0472,-0.6481) (0.4363,1.0472,-0.6830) (0.5236,1.0472,-0.7165) (0.6109,1.0472,-0.7484) (0.6981,1.0472,-0.7783) (0.7854,1.0472,-0.8062) (0.8727,1.0472,-0.8317) (0.9599,1.0472,-0.8547) (1.0472,1.0472,-0.8750) (1.1345,1.0472,-0.8924) (1.2217,1.0472,-0.9069) (1.3090,1.0472,-0.9183) (1.3963,1.0472,-0.9264) (1.4835,1.0472,-0.9314) (1.5708,1.0472,-0.9330) (1.6581,1.0472,-0.9314) (1.7453,1.0472,-0.9264) (1.8326,1.0472,-0.9183) (1.9199,1.0472,-0.9069) (2.0071,1.0472,-0.8924) (2.0944,1.0472,-0.8750) (2.1817,1.0472,-0.8547) (2.2689,1.0472,-0.8317) (2.3562,1.0472,-0.8062) (2.4435,1.0472,-0.7783) (2.5307,1.0472,-0.7484) (2.6180,1.0472,-0.7165) (2.7053,1.0472,-0.6830) (2.7925,1.0472,-0.6481) (2.8798,1.0472,-0.6121) (2.9671,1.0472,-0.5752) (3.0543,1.0472,-0.5377) (3.1416,1.0472,-0.5000) (3.2289,1.0472,-0.4623) (3.3161,1.0472,-0.4248) (3.4034,1.0472,-0.3879) (3.4907,1.0472,-0.3519) (3.5779,1.0472,-0.3170) (3.6652,1.0472,-0.2835) (3.7525,1.0472,-0.2516) (3.8397,1.0472,-0.2217) (3.9270,1.0472,-0.1938) (4.0143,1.0472,-0.1683) (4.1015,1.0472,-0.1453) (4.1888,1.0472,-0.1250) (4.2761,1.0472,-0.1076) (4.3633,1.0472,-0.0931) (4.4506,1.0472,-0.0817) (4.5379,1.0472,-0.0736) (4.6251,1.0472,-0.0686) (4.7124,1.0472,-0.0670) (4.7997,1.0472,-0.0686) (4.8869,1.0472,-0.0736) (4.9742,1.0472,-0.0817) (5.0615,1.0472,-0.0931) (5.1487,1.0472,-0.1076) (5.2360,1.0472,-0.1250) (5.3233,1.0472,-0.1453) (5.4105,1.0472,-0.1683) (5.4978,1.0472,-0.1938) (5.5851,1.0472,-0.2217) (5.6723,1.0472,-0.2516) (5.7596,1.0472,-0.2835) (5.8469,1.0472,-0.3170) (5.9341,1.0472,-0.3519) (6.0214,1.0472,-0.3879) (6.1087,1.0472,-0.4248) (6.1959,1.0472,-0.4623) (6.2832,1.0472,-0.5000)
(0.0000,1.1345,-0.5000) (0.0873,1.1345,-0.5524) (0.1745,1.1345,-0.6039) (0.2618,1.1345,-0.6539) (0.3491,1.1345,-0.7014) (0.4363,1.1345,-0.7461) (0.5236,1.1345,-0.7872) (0.6109,1.1345,-0.8245) (0.6981,1.1345,-0.8578) (0.7854,1.1345,-0.8869) (0.8727,1.1345,-0.9118) (0.9599,1.1345,-0.9328) (1.0472,1.1345,-0.9501) (1.1345,1.1345,-0.9640) (1.2217,1.1345,-0.9747) (1.3090,1.1345,-0.9827) (1.3963,1.1345,-0.9882) (1.4835,1.1345,-0.9914) (1.5708,1.1345,-0.9924) (1.6581,1.1345,-0.9914) (1.7453,1.1345,-0.9882) (1.8326,1.1345,-0.9827) (1.9199,1.1345,-0.9747) (2.0071,1.1345,-0.9640) (2.0944,1.1345,-0.9501) (2.1817,1.1345,-0.9328) (2.2689,1.1345,-0.9118) (2.3562,1.1345,-0.8869) (2.4435,1.1345,-0.8578) (2.5307,1.1345,-0.8245) (2.6180,1.1345,-0.7872) (2.7053,1.1345,-0.7461) (2.7925,1.1345,-0.7014) (2.8798,1.1345,-0.6539) (2.9671,1.1345,-0.6039) (3.0543,1.1345,-0.5524) (3.1416,1.1345,-0.5000) (3.2289,1.1345,-0.4476) (3.3161,1.1345,-0.3961) (3.4034,1.1345,-0.3461) (3.4907,1.1345,-0.2986) (3.5779,1.1345,-0.2539) (3.6652,1.1345,-0.2128) (3.7525,1.1345,-0.1755) (3.8397,1.1345,-0.1422) (3.9270,1.1345,-0.1131) (4.0143,1.1345,-0.0882) (4.1015,1.1345,-0.0672) (4.1888,1.1345,-0.0499) (4.2761,1.1345,-0.0360) (4.3633,1.1345,-0.0253) (4.4506,1.1345,-0.0173) (4.5379,1.1345,-0.0118) (4.6251,1.1345,-0.0086) (4.7124,1.1345,-0.0076) (4.7997,1.1345,-0.0086) (4.8869,1.1345,-0.0118) (4.9742,1.1345,-0.0173) (5.0615,1.1345,-0.0253) (5.1487,1.1345,-0.0360) (5.2360,1.1345,-0.0499) (5.3233,1.1345,-0.0672) (5.4105,1.1345,-0.0882) (5.4978,1.1345,-0.1131) (5.5851,1.1345,-0.1422) (5.6723,1.1345,-0.1755) (5.7596,1.1345,-0.2128) (5.8469,1.1345,-0.2539) (5.9341,1.1345,-0.2986) (6.0214,1.1345,-0.3461) (6.1087,1.1345,-0.3961) (6.1959,1.1345,-0.4476) (6.2832,1.1345,-0.5000)
(0.0000,1.2217,-0.5000) (0.0873,1.2217,-0.5577) (0.1745,1.2217,-0.6143) (0.2618,1.2217,-0.6687) (0.3491,1.2217,-0.7201) (0.4363,1.2217,-0.7675) (0.5236,1.2217,-0.8103) (0.6109,1.2217,-0.8482) (0.6981,1.2217,-0.8810) (0.7854,1.2217,-0.9086) (0.8727,1.2217,-0.9313) (0.9599,1.2217,-0.9494) (1.0472,1.2217,-0.9635) (1.1345,1.2217,-0.9740) (1.2217,1.2217,-0.9815) (1.3090,1.2217,-0.9867) (1.3963,1.2217,-0.9900) (1.4835,1.2217,-0.9918) (1.5708,1.2217,-0.9924) (1.6581,1.2217,-0.9918) (1.7453,1.2217,-0.9900) (1.8326,1.2217,-0.9867) (1.9199,1.2217,-0.9815) (2.0071,1.2217,-0.9740) (2.0944,1.2217,-0.9635) (2.1817,1.2217,-0.9494) (2.2689,1.2217,-0.9313) (2.3562,1.2217,-0.9086) (2.4435,1.2217,-0.8810) (2.5307,1.2217,-0.8482) (2.6180,1.2217,-0.8103) (2.7053,1.2217,-0.7675) (2.7925,1.2217,-0.7201) (2.8798,1.2217,-0.6687) (2.9671,1.2217,-0.6143) (3.0543,1.2217,-0.5577) (3.1416,1.2217,-0.5000) (3.2289,1.2217,-0.4423) (3.3161,1.2217,-0.3857) (3.4034,1.2217,-0.3313) (3.4907,1.2217,-0.2799) (3.5779,1.2217,-0.2325) (3.6652,1.2217,-0.1897) (3.7525,1.2217,-0.1518) (3.8397,1.2217,-0.1190) (3.9270,1.2217,-0.0914) (4.0143,1.2217,-0.0687) (4.1015,1.2217,-0.0506) (4.1888,1.2217,-0.0365) (4.2761,1.2217,-0.0260) (4.3633,1.2217,-0.0185) (4.4506,1.2217,-0.0133) (4.5379,1.2217,-0.0100) (4.6251,1.2217,-0.0082) (4.7124,1.2217,-0.0076) (4.7997,1.2217,-0.0082) (4.8869,1.2217,-0.0100) (4.9742,1.2217,-0.0133) (5.0615,1.2217,-0.0185) (5.1487,1.2217,-0.0260) (5.2360,1.2217,-0.0365) (5.3233,1.2217,-0.0506) (5.4105,1.2217,-0.0687) (5.4978,1.2217,-0.0914) (5.5851,1.2217,-0.1190) (5.6723,1.2217,-0.1518) (5.7596,1.2217,-0.1897) (5.8469,1.2217,-0.2325) (5.9341,1.2217,-0.2799) (6.0214,1.2217,-0.3313) (6.1087,1.2217,-0.3857) (6.1959,1.2217,-0.4423) (6.2832,1.2217,-0.5000)
(0.0000,1.3090,-0.5000) (0.0873,1.3090,-0.5536) (0.1745,1.3090,-0.6060) (0.2618,1.3090,-0.6563) (0.3491,1.3090,-0.7034) (0.4363,1.3090,-0.7465) (0.5236,1.3090,-0.7851) (0.6109,1.3090,-0.8188) (0.6981,1.3090,-0.8474) (0.7854,1.3090,-0.8709) (0.8727,1.3090,-0.8896) (0.9599,1.3090,-0.9040) (1.0472,1.3090,-0.9146) (1.1345,1.3090,-0.9221) (1.2217,1.3090,-0.9270) (1.3090,1.3090,-0.9301) (1.3963,1.3090,-0.9319) (1.4835,1.3090,-0.9327) (1.5708,1.3090,-0.9330) (1.6581,1.3090,-0.9327) (1.7453,1.3090,-0.9319) (1.8326,1.3090,-0.9301) (1.9199,1.3090,-0.9270) (2.0071,1.3090,-0.9221) (2.0944,1.3090,-0.9146) (2.1817,1.3090,-0.9040) (2.2689,1.3090,-0.8896) (2.3562,1.3090,-0.8709) (2.4435,1.3090,-0.8474) (2.5307,1.3090,-0.8188) (2.6180,1.3090,-0.7851) (2.7053,1.3090,-0.7465) (2.7925,1.3090,-0.7034) (2.8798,1.3090,-0.6563) (2.9671,1.3090,-0.6060) (3.0543,1.3090,-0.5536) (3.1416,1.3090,-0.5000) (3.2289,1.3090,-0.4464) (3.3161,1.3090,-0.3940) (3.4034,1.3090,-0.3437) (3.4907,1.3090,-0.2966) (3.5779,1.3090,-0.2535) (3.6652,1.3090,-0.2149) (3.7525,1.3090,-0.1812) (3.8397,1.3090,-0.1526) (3.9270,1.3090,-0.1291) (4.0143,1.3090,-0.1104) (4.1015,1.3090,-0.0960) (4.1888,1.3090,-0.0854) (4.2761,1.3090,-0.0779) (4.3633,1.3090,-0.0730) (4.4506,1.3090,-0.0699) (4.5379,1.3090,-0.0681) (4.6251,1.3090,-0.0673) (4.7124,1.3090,-0.0670) (4.7997,1.3090,-0.0673) (4.8869,1.3090,-0.0681) (4.9742,1.3090,-0.0699) (5.0615,1.3090,-0.0730) (5.1487,1.3090,-0.0779) (5.2360,1.3090,-0.0854) (5.3233,1.3090,-0.0960) (5.4105,1.3090,-0.1104) (5.4978,1.3090,-0.1291) (5.5851,1.3090,-0.1526) (5.6723,1.3090,-0.1812) (5.7596,1.3090,-0.2149) (5.8469,1.3090,-0.2535) (5.9341,1.3090,-0.2966) (6.0214,1.3090,-0.3437) (6.1087,1.3090,-0.3940) (6.1959,1.3090,-0.4464) (6.2832,1.3090,-0.5000)
(0.0000,1.3963,-0.5000) (0.0873,1.3963,-0.5410) (0.1745,1.3963,-0.5811) (0.2618,1.3963,-0.6195) (0.3491,1.3963,-0.6553) (0.4363,1.3963,-0.6880) (0.5236,1.3963,-0.7171) (0.6109,1.3963,-0.7422) (0.6981,1.3963,-0.7633) (0.7854,1.3963,-0.7804) (0.8727,1.3963,-0.7938) (0.9599,1.3963,-0.8038) (1.0472,1.3963,-0.8109) (1.1345,1.3963,-0.8156) (1.2217,1.3963,-0.8185) (1.3090,1.3963,-0.8202) (1.3963,1.3963,-0.8210) (1.4835,1.3963,-0.8213) (1.5708,1.3963,-0.8214) (1.6581,1.3963,-0.8213) (1.7453,1.3963,-0.8210) (1.8326,1.3963,-0.8202) (1.9199,1.3963,-0.8185) (2.0071,1.3963,-0.8156) (2.0944,1.3963,-0.8109) (2.1817,1.3963,-0.8038) (2.2689,1.3963,-0.7938) (2.3562,1.3963,-0.7804) (2.4435,1.3963,-0.7633) (2.5307,1.3963,-0.7422) (2.6180,1.3963,-0.7171) (2.7053,1.3963,-0.6880) (2.7925,1.3963,-0.6553) (2.8798,1.3963,-0.6195) (2.9671,1.3963,-0.5811) (3.0543,1.3963,-0.5410) (3.1416,1.3963,-0.5000) (3.2289,1.3963,-0.4590) (3.3161,1.3963,-0.4189) (3.4034,1.3963,-0.3805) (3.4907,1.3963,-0.3447) (3.5779,1.3963,-0.3120) (3.6652,1.3963,-0.2829) (3.7525,1.3963,-0.2578) (3.8397,1.3963,-0.2367) (3.9270,1.3963,-0.2196) (4.0143,1.3963,-0.2062) (4.1015,1.3963,-0.1962) (4.1888,1.3963,-0.1891) (4.2761,1.3963,-0.1844) (4.3633,1.3963,-0.1815) (4.4506,1.3963,-0.1798) (4.5379,1.3963,-0.1790) (4.6251,1.3963,-0.1787) (4.7124,1.3963,-0.1786) (4.7997,1.3963,-0.1787) (4.8869,1.3963,-0.1790) (4.9742,1.3963,-0.1798) (5.0615,1.3963,-0.1815) (5.1487,1.3963,-0.1844) (5.2360,1.3963,-0.1891) (5.3233,1.3963,-0.1962) (5.4105,1.3963,-0.2062) (5.4978,1.3963,-0.2196) (5.5851,1.3963,-0.2367) (5.6723,1.3963,-0.2578) (5.7596,1.3963,-0.2829) (5.8469,1.3963,-0.3120) (5.9341,1.3963,-0.3447) (6.0214,1.3963,-0.3805) (6.1087,1.3963,-0.4189) (6.1959,1.3963,-0.4590) (6.2832,1.3963,-0.5000)
(0.0000,1.4835,-0.5000) (0.0873,1.4835,-0.5222) (0.1745,1.4835,-0.5439) (0.2618,1.4835,-0.5646) (0.3491,1.4835,-0.5839) (0.4363,1.4835,-0.6015) (0.5236,1.4835,-0.6171) (0.6109,1.4835,-0.6305) (0.6981,1.4835,-0.6417) (0.7854,1.4835,-0.6507) (0.8727,1.4835,-0.6576) (0.9599,1.4835,-0.6628) (1.0472,1.4835,-0.6663) (1.1345,1.4835,-0.6686) (1.2217,1.4835,-0.6700) (1.3090,1.4835,-0.6706) (1.3963,1.4835,-0.6709) (1.4835,1.4835,-0.6710) (1.5708,1.4835,-0.6710) (1.6581,1.4835,-0.6710) (1.7453,1.4835,-0.6709) (1.8326,1.4835,-0.6706) (1.9199,1.4835,-0.6700) (2.0071,1.4835,-0.6686) (2.0944,1.4835,-0.6663) (2.1817,1.4835,-0.6628) (2.2689,1.4835,-0.6576) (2.3562,1.4835,-0.6507) (2.4435,1.4835,-0.6417) (2.5307,1.4835,-0.6305) (2.6180,1.4835,-0.6171) (2.7053,1.4835,-0.6015) (2.7925,1.4835,-0.5839) (2.8798,1.4835,-0.5646) (2.9671,1.4835,-0.5439) (3.0543,1.4835,-0.5222) (3.1416,1.4835,-0.5000) (3.2289,1.4835,-0.4778) (3.3161,1.4835,-0.4561) (3.4034,1.4835,-0.4354) (3.4907,1.4835,-0.4161) (3.5779,1.4835,-0.3985) (3.6652,1.4835,-0.3829) (3.7525,1.4835,-0.3695) (3.8397,1.4835,-0.3583) (3.9270,1.4835,-0.3493) (4.0143,1.4835,-0.3424) (4.1015,1.4835,-0.3372) (4.1888,1.4835,-0.3337) (4.2761,1.4835,-0.3314) (4.3633,1.4835,-0.3300) (4.4506,1.4835,-0.3294) (4.5379,1.4835,-0.3291) (4.6251,1.4835,-0.3290) (4.7124,1.4835,-0.3290) (4.7997,1.4835,-0.3290) (4.8869,1.4835,-0.3291) (4.9742,1.4835,-0.3294) (5.0615,1.4835,-0.3300) (5.1487,1.4835,-0.3314) (5.2360,1.4835,-0.3337) (5.3233,1.4835,-0.3372) (5.4105,1.4835,-0.3424) (5.4978,1.4835,-0.3493) (5.5851,1.4835,-0.3583) (5.6723,1.4835,-0.3695) (5.7596,1.4835,-0.3829) (5.8469,1.4835,-0.3985) (5.9341,1.4835,-0.4161) (6.0214,1.4835,-0.4354) (6.1087,1.4835,-0.4561) (6.1959,1.4835,-0.4778) (6.2832,1.4835,-0.5000)
(0.0000,1.5708,-0.5000) (0.0873,1.5708,-0.5000) (0.1745,1.5708,-0.5000) (0.2618,1.5708,-0.5000) (0.3491,1.5708,-0.5000) (0.4363,1.5708,-0.5000) (0.5236,1.5708,-0.5000) (0.6109,1.5708,-0.5000) (0.6981,1.5708,-0.5000) (0.7854,1.5708,-0.5000) (0.8727,1.5708,-0.5000) (0.9599,1.5708,-0.5000) (1.0472,1.5708,-0.5000) (1.1345,1.5708,-0.5000) (1.2217,1.5708,-0.5000) (1.3090,1.5708,-0.5000) (1.3963,1.5708,-0.5000) (1.4835,1.5708,-0.5000) (1.5708,1.5708,-0.5000) (1.6581,1.5708,-0.5000) (1.7453,1.5708,-0.5000) (1.8326,1.5708,-0.5000) (1.9199,1.5708,-0.5000) (2.0071,1.5708,-0.5000) (2.0944,1.5708,-0.5000) (2.1817,1.5708,-0.5000) (2.2689,1.5708,-0.5000) (2.3562,1.5708,-0.5000) (2.4435,1.5708,-0.5000) (2.5307,1.5708,-0.5000) (2.6180,1.5708,-0.5000) (2.7053,1.5708,-0.5000) (2.7925,1.5708,-0.5000) (2.8798,1.5708,-0.5000) (2.9671,1.5708,-0.5000) (3.0543,1.5708,-0.5000) (3.1416,1.5708,-0.5000) (3.2289,1.5708,-0.5000) (3.3161,1.5708,-0.5000) (3.4034,1.5708,-0.5000) (3.4907,1.5708,-0.5000) (3.5779,1.5708,-0.5000) (3.6652,1.5708,-0.5000) (3.7525,1.5708,-0.5000) (3.8397,1.5708,-0.5000) (3.9270,1.5708,-0.5000) (4.0143,1.5708,-0.5000) (4.1015,1.5708,-0.5000) (4.1888,1.5708,-0.5000) (4.2761,1.5708,-0.5000) (4.3633,1.5708,-0.5000) (4.4506,1.5708,-0.5000) (4.5379,1.5708,-0.5000) (4.6251,1.5708,-0.5000) (4.7124,1.5708,-0.5000) (4.7997,1.5708,-0.5000) (4.8869,1.5708,-0.5000) (4.9742,1.5708,-0.5000) (5.0615,1.5708,-0.5000) (5.1487,1.5708,-0.5000) (5.2360,1.5708,-0.5000) (5.3233,1.5708,-0.5000) (5.4105,1.5708,-0.5000) (5.4978,1.5708,-0.5000) (5.5851,1.5708,-0.5000) (5.6723,1.5708,-0.5000) (5.7596,1.5708,-0.5000) (5.8469,1.5708,-0.5000) (5.9341,1.5708,-0.5000) (6.0214,1.5708,-0.5000) (6.1087,1.5708,-0.5000) (6.1959,1.5708,-0.5000) (6.2832,1.5708,-0.5000)
(0.0000,1.6581,-0.5000) (0.0873,1.6581,-0.4778) (0.1745,1.6581,-0.4561) (0.2618,1.6581,-0.4354) (0.3491,1.6581,-0.4161) (0.4363,1.6581,-0.3985) (0.5236,1.6581,-0.3829) (0.6109,1.6581,-0.3695) (0.6981,1.6581,-0.3583) (0.7854,1.6581,-0.3493) (0.8727,1.6581,-0.3424) (0.9599,1.6581,-0.3372) (1.0472,1.6581,-0.3337) (1.1345,1.6581,-0.3314) (1.2217,1.6581,-0.3300) (1.3090,1.6581,-0.3294) (1.3963,1.6581,-0.3291) (1.4835,1.6581,-0.3290) (1.5708,1.6581,-0.3290) (1.6581,1.6581,-0.3290) (1.7453,1.6581,-0.3291) (1.8326,1.6581,-0.3294) (1.9199,1.6581,-0.3300) (2.0071,1.6581,-0.3314) (2.0944,1.6581,-0.3337) (2.1817,1.6581,-0.3372) (2.2689,1.6581,-0.3424) (2.3562,1.6581,-0.3493) (2.4435,1.6581,-0.3583) (2.5307,1.6581,-0.3695) (2.6180,1.6581,-0.3829) (2.7053,1.6581,-0.3985) (2.7925,1.6581,-0.4161) (2.8798,1.6581,-0.4354) (2.9671,1.6581,-0.4561) (3.0543,1.6581,-0.4778) (3.1416,1.6581,-0.5000) (3.2289,1.6581,-0.5222) (3.3161,1.6581,-0.5439) (3.4034,1.6581,-0.5646) (3.4907,1.6581,-0.5839) (3.5779,1.6581,-0.6015) (3.6652,1.6581,-0.6171) (3.7525,1.6581,-0.6305) (3.8397,1.6581,-0.6417) (3.9270,1.6581,-0.6507) (4.0143,1.6581,-0.6576) (4.1015,1.6581,-0.6628) (4.1888,1.6581,-0.6663) (4.2761,1.6581,-0.6686) (4.3633,1.6581,-0.6700) (4.4506,1.6581,-0.6706) (4.5379,1.6581,-0.6709) (4.6251,1.6581,-0.6710) (4.7124,1.6581,-0.6710) (4.7997,1.6581,-0.6710) (4.8869,1.6581,-0.6709) (4.9742,1.6581,-0.6706) (5.0615,1.6581,-0.6700) (5.1487,1.6581,-0.6686) (5.2360,1.6581,-0.6663) (5.3233,1.6581,-0.6628) (5.4105,1.6581,-0.6576) (5.4978,1.6581,-0.6507) (5.5851,1.6581,-0.6417) (5.6723,1.6581,-0.6305) (5.7596,1.6581,-0.6171) (5.8469,1.6581,-0.6015) (5.9341,1.6581,-0.5839) (6.0214,1.6581,-0.5646) (6.1087,1.6581,-0.5439) (6.1959,1.6581,-0.5222) (6.2832,1.6581,-0.5000)
(0.0000,1.7453,-0.5000) (0.0873,1.7453,-0.4590) (0.1745,1.7453,-0.4189) (0.2618,1.7453,-0.3805) (0.3491,1.7453,-0.3447) (0.4363,1.7453,-0.3120) (0.5236,1.7453,-0.2829) (0.6109,1.7453,-0.2578) (0.6981,1.7453,-0.2367) (0.7854,1.7453,-0.2196) (0.8727,1.7453,-0.2062) (0.9599,1.7453,-0.1962) (1.0472,1.7453,-0.1891) (1.1345,1.7453,-0.1844) (1.2217,1.7453,-0.1815) (1.3090,1.7453,-0.1798) (1.3963,1.7453,-0.1790) (1.4835,1.7453,-0.1787) (1.5708,1.7453,-0.1786) (1.6581,1.7453,-0.1787) (1.7453,1.7453,-0.1790) (1.8326,1.7453,-0.1798) (1.9199,1.7453,-0.1815) (2.0071,1.7453,-0.1844) (2.0944,1.7453,-0.1891) (2.1817,1.7453,-0.1962) (2.2689,1.7453,-0.2062) (2.3562,1.7453,-0.2196) (2.4435,1.7453,-0.2367) (2.5307,1.7453,-0.2578) (2.6180,1.7453,-0.2829) (2.7053,1.7453,-0.3120) (2.7925,1.7453,-0.3447) (2.8798,1.7453,-0.3805) (2.9671,1.7453,-0.4189) (3.0543,1.7453,-0.4590) (3.1416,1.7453,-0.5000) (3.2289,1.7453,-0.5410) (3.3161,1.7453,-0.5811) (3.4034,1.7453,-0.6195) (3.4907,1.7453,-0.6553) (3.5779,1.7453,-0.6880) (3.6652,1.7453,-0.7171) (3.7525,1.7453,-0.7422) (3.8397,1.7453,-0.7633) (3.9270,1.7453,-0.7804) (4.0143,1.7453,-0.7938) (4.1015,1.7453,-0.8038) (4.1888,1.7453,-0.8109) (4.2761,1.7453,-0.8156) (4.3633,1.7453,-0.8185) (4.4506,1.7453,-0.8202) (4.5379,1.7453,-0.8210) (4.6251,1.7453,-0.8213) (4.7124,1.7453,-0.8214) (4.7997,1.7453,-0.8213) (4.8869,1.7453,-0.8210) (4.9742,1.7453,-0.8202) (5.0615,1.7453,-0.8185) (5.1487,1.7453,-0.8156) (5.2360,1.7453,-0.8109) (5.3233,1.7453,-0.8038) (5.4105,1.7453,-0.7938) (5.4978,1.7453,-0.7804) (5.5851,1.7453,-0.7633) (5.6723,1.7453,-0.7422) (5.7596,1.7453,-0.7171) (5.8469,1.7453,-0.6880) (5.9341,1.7453,-0.6553) (6.0214,1.7453,-0.6195) (6.1087,1.7453,-0.5811) (6.1959,1.7453,-0.5410) (6.2832,1.7453,-0.5000)
(0.0000,1.8326,-0.5000) (0.0873,1.8326,-0.4464) (0.1745,1.8326,-0.3940) (0.2618,1.8326,-0.3437) (0.3491,1.8326,-0.2966) (0.4363,1.8326,-0.2535) (0.5236,1.8326,-0.2149) (0.6109,1.8326,-0.1812) (0.6981,1.8326,-0.1526) (0.7854,1.8326,-0.1291) (0.8727,1.8326,-0.1104) (0.9599,1.8326,-0.0960) (1.0472,1.8326,-0.0854) (1.1345,1.8326,-0.0779) (1.2217,1.8326,-0.0730) (1.3090,1.8326,-0.0699) (1.3963,1.8326,-0.0681) (1.4835,1.8326,-0.0673) (1.5708,1.8326,-0.0670) (1.6581,1.8326,-0.0673) (1.7453,1.8326,-0.0681) (1.8326,1.8326,-0.0699) (1.9199,1.8326,-0.0730) (2.0071,1.8326,-0.0779) (2.0944,1.8326,-0.0854) (2.1817,1.8326,-0.0960) (2.2689,1.8326,-0.1104) (2.3562,1.8326,-0.1291) (2.4435,1.8326,-0.1526) (2.5307,1.8326,-0.1812) (2.6180,1.8326,-0.2149) (2.7053,1.8326,-0.2535) (2.7925,1.8326,-0.2966) (2.8798,1.8326,-0.3437) (2.9671,1.8326,-0.3940) (3.0543,1.8326,-0.4464) (3.1416,1.8326,-0.5000) (3.2289,1.8326,-0.5536) (3.3161,1.8326,-0.6060) (3.4034,1.8326,-0.6563) (3.4907,1.8326,-0.7034) (3.5779,1.8326,-0.7465) (3.6652,1.8326,-0.7851) (3.7525,1.8326,-0.8188) (3.8397,1.8326,-0.8474) (3.9270,1.8326,-0.8709) (4.0143,1.8326,-0.8896) (4.1015,1.8326,-0.9040) (4.1888,1.8326,-0.9146) (4.2761,1.8326,-0.9221) (4.3633,1.8326,-0.9270) (4.4506,1.8326,-0.9301) (4.5379,1.8326,-0.9319) (4.6251,1.8326,-0.9327) (4.7124,1.8326,-0.9330) (4.7997,1.8326,-0.9327) (4.8869,1.8326,-0.9319) (4.9742,1.8326,-0.9301) (5.0615,1.8326,-0.9270) (5.1487,1.8326,-0.9221) (5.2360,1.8326,-0.9146) (5.3233,1.8326,-0.9040) (5.4105,1.8326,-0.8896) (5.4978,1.8326,-0.8709) (5.5851,1.8326,-0.8474) (5.6723,1.8326,-0.8188) (5.7596,1.8326,-0.7851) (5.8469,1.8326,-0.7465) (5.9341,1.8326,-0.7034) (6.0214,1.8326,-0.6563) (6.1087,1.8326,-0.6060) (6.1959,1.8326,-0.5536) (6.2832,1.8326,-0.5000)
(0.0000,1.9199,-0.5000) (0.0873,1.9199,-0.4423) (0.1745,1.9199,-0.3857) (0.2618,1.9199,-0.3313) (0.3491,1.9199,-0.2799) (0.4363,1.9199,-0.2325) (0.5236,1.9199,-0.1897) (0.6109,1.9199,-0.1518) (0.6981,1.9199,-0.1190) (0.7854,1.9199,-0.0914) (0.8727,1.9199,-0.0687) (0.9599,1.9199,-0.0506) (1.0472,1.9199,-0.0365) (1.1345,1.9199,-0.0260) (1.2217,1.9199,-0.0185) (1.3090,1.9199,-0.0133) (1.3963,1.9199,-0.0100) (1.4835,1.9199,-0.0082) (1.5708,1.9199,-0.0076) (1.6581,1.9199,-0.0082) (1.7453,1.9199,-0.0100) (1.8326,1.9199,-0.0133) (1.9199,1.9199,-0.0185) (2.0071,1.9199,-0.0260) (2.0944,1.9199,-0.0365) (2.1817,1.9199,-0.0506) (2.2689,1.9199,-0.0687) (2.3562,1.9199,-0.0914) (2.4435,1.9199,-0.1190) (2.5307,1.9199,-0.1518) (2.6180,1.9199,-0.1897) (2.7053,1.9199,-0.2325) (2.7925,1.9199,-0.2799) (2.8798,1.9199,-0.3313) (2.9671,1.9199,-0.3857) (3.0543,1.9199,-0.4423) (3.1416,1.9199,-0.5000) (3.2289,1.9199,-0.5577) (3.3161,1.9199,-0.6143) (3.4034,1.9199,-0.6687) (3.4907,1.9199,-0.7201) (3.5779,1.9199,-0.7675) (3.6652,1.9199,-0.8103) (3.7525,1.9199,-0.8482) (3.8397,1.9199,-0.8810) (3.9270,1.9199,-0.9086) (4.0143,1.9199,-0.9313) (4.1015,1.9199,-0.9494) (4.1888,1.9199,-0.9635) (4.2761,1.9199,-0.9740) (4.3633,1.9199,-0.9815) (4.4506,1.9199,-0.9867) (4.5379,1.9199,-0.9900) (4.6251,1.9199,-0.9918) (4.7124,1.9199,-0.9924) (4.7997,1.9199,-0.9918) (4.8869,1.9199,-0.9900) (4.9742,1.9199,-0.9867) (5.0615,1.9199,-0.9815) (5.1487,1.9199,-0.9740) (5.2360,1.9199,-0.9635) (5.3233,1.9199,-0.9494) (5.4105,1.9199,-0.9313) (5.4978,1.9199,-0.9086) (5.5851,1.9199,-0.8810) (5.6723,1.9199,-0.8482) (5.7596,1.9199,-0.8103) (5.8469,1.9199,-0.7675) (5.9341,1.9199,-0.7201) (6.0214,1.9199,-0.6687) (6.1087,1.9199,-0.6143) (6.1959,1.9199,-0.5577) (6.2832,1.9199,-0.5000)
(0.0000,2.0071,-0.5000) (0.0873,2.0071,-0.4476) (0.1745,2.0071,-0.3961) (0.2618,2.0071,-0.3461) (0.3491,2.0071,-0.2986) (0.4363,2.0071,-0.2539) (0.5236,2.0071,-0.2128) (0.6109,2.0071,-0.1755) (0.6981,2.0071,-0.1422) (0.7854,2.0071,-0.1131) (0.8727,2.0071,-0.0882) (0.9599,2.0071,-0.0672) (1.0472,2.0071,-0.0499) (1.1345,2.0071,-0.0360) (1.2217,2.0071,-0.0253) (1.3090,2.0071,-0.0173) (1.3963,2.0071,-0.0118) (1.4835,2.0071,-0.0086) (1.5708,2.0071,-0.0076) (1.6581,2.0071,-0.0086) (1.7453,2.0071,-0.0118) (1.8326,2.0071,-0.0173) (1.9199,2.0071,-0.0253) (2.0071,2.0071,-0.0360) (2.0944,2.0071,-0.0499) (2.1817,2.0071,-0.0672) (2.2689,2.0071,-0.0882) (2.3562,2.0071,-0.1131) (2.4435,2.0071,-0.1422) (2.5307,2.0071,-0.1755) (2.6180,2.0071,-0.2128) (2.7053,2.0071,-0.2539) (2.7925,2.0071,-0.2986) (2.8798,2.0071,-0.3461) (2.9671,2.0071,-0.3961) (3.0543,2.0071,-0.4476) (3.1416,2.0071,-0.5000) (3.2289,2.0071,-0.5524) (3.3161,2.0071,-0.6039) (3.4034,2.0071,-0.6539) (3.4907,2.0071,-0.7014) (3.5779,2.0071,-0.7461) (3.6652,2.0071,-0.7872) (3.7525,2.0071,-0.8245) (3.8397,2.0071,-0.8578) (3.9270,2.0071,-0.8869) (4.0143,2.0071,-0.9118) (4.1015,2.0071,-0.9328) (4.1888,2.0071,-0.9501) (4.2761,2.0071,-0.9640) (4.3633,2.0071,-0.9747) (4.4506,2.0071,-0.9827) (4.5379,2.0071,-0.9882) (4.6251,2.0071,-0.9914) (4.7124,2.0071,-0.9924) (4.7997,2.0071,-0.9914) (4.8869,2.0071,-0.9882) (4.9742,2.0071,-0.9827) (5.0615,2.0071,-0.9747) (5.1487,2.0071,-0.9640) (5.2360,2.0071,-0.9501) (5.3233,2.0071,-0.9328) (5.4105,2.0071,-0.9118) (5.4978,2.0071,-0.8869) (5.5851,2.0071,-0.8578) (5.6723,2.0071,-0.8245) (5.7596,2.0071,-0.7872) (5.8469,2.0071,-0.7461) (5.9341,2.0071,-0.7014) (6.0214,2.0071,-0.6539) (6.1087,2.0071,-0.6039) (6.1959,2.0071,-0.5524) (6.2832,2.0071,-0.5000)
(0.0000,2.0944,-0.5000) (0.0873,2.0944,-0.4623) (0.1745,2.0944,-0.4248) (0.2618,2.0944,-0.3879) (0.3491,2.0944,-0.3519) (0.4363,2.0944,-0.3170) (0.5236,2.0944,-0.2835) (0.6109,2.0944,-0.2516) (0.6981,2.0944,-0.2217) (0.7854,2.0944,-0.1938) (0.8727,2.0944,-0.1683) (0.9599,2.0944,-0.1453) (1.0472,2.0944,-0.1250) (1.1345,2.0944,-0.1076) (1.2217,2.0944,-0.0931) (1.3090,2.0944,-0.0817) (1.3963,2.0944,-0.0736) (1.4835,2.0944,-0.0686) (1.5708,2.0944,-0.0670) (1.6581,2.0944,-0.0686) (1.7453,2.0944,-0.0736) (1.8326,2.0944,-0.0817) (1.9199,2.0944,-0.0931) (2.0071,2.0944,-0.1076) (2.0944,2.0944,-0.1250) (2.1817,2.0944,-0.1453) (2.2689,2.0944,-0.1683) (2.3562,2.0944,-0.1938) (2.4435,2.0944,-0.2217) (2.5307,2.0944,-0.2516) (2.6180,2.0944,-0.2835) (2.7053,2.0944,-0.3170) (2.7925,2.0944,-0.3519) (2.8798,2.0944,-0.3879) (2.9671,2.0944,-0.4248) (3.0543,2.0944,-0.4623) (3.1416,2.0944,-0.5000) (3.2289,2.0944,-0.5377) (3.3161,2.0944,-0.5752) (3.4034,2.0944,-0.6121) (3.4907,2.0944,-0.6481) (3.5779,2.0944,-0.6830) (3.6652,2.0944,-0.7165) (3.7525,2.0944,-0.7484) (3.8397,2.0944,-0.7783) (3.9270,2.0944,-0.8062) (4.0143,2.0944,-0.8317) (4.1015,2.0944,-0.8547) (4.1888,2.0944,-0.8750) (4.2761,2.0944,-0.8924) (4.3633,2.0944,-0.9069) (4.4506,2.0944,-0.9183) (4.5379,2.0944,-0.9264) (4.6251,2.0944,-0.9314) (4.7124,2.0944,-0.9330) (4.7997,2.0944,-0.9314) (4.8869,2.0944,-0.9264) (4.9742,2.0944,-0.9183) (5.0615,2.0944,-0.9069) (5.1487,2.0944,-0.8924) (5.2360,2.0944,-0.8750) (5.3233,2.0944,-0.8547) (5.4105,2.0944,-0.8317) (5.4978,2.0944,-0.8062) (5.5851,2.0944,-0.7783) (5.6723,2.0944,-0.7484) (5.7596,2.0944,-0.7165) (5.8469,2.0944,-0.6830) (5.9341,2.0944,-0.6481) (6.0214,2.0944,-0.6121) (6.1087,2.0944,-0.5752) (6.1959,2.0944,-0.5377) (6.2832,2.0944,-0.5000)
(0.0000,2.1817,-0.5000) (0.0873,2.1817,-0.4848) (0.1745,2.1817,-0.4692) (0.2618,2.1817,-0.4527) (0.3491,2.1817,-0.4349) (0.4363,2.1817,-0.4157) (0.5236,2.1817,-0.3950) (0.6109,2.1817,-0.3728) (0.6981,2.1817,-0.3494) (0.7854,2.1817,-0.3252) (0.8727,2.1817,-0.3008) (0.9599,2.1817,-0.2767) (1.0472,2.1817,-0.2538) (1.1345,2.1817,-0.2327) (1.2217,2.1817,-0.2143) (1.3090,2.1817,-0.1992) (1.3963,2.1817,-0.1879) (1.4835,2.1817,-0.1810) (1.5708,2.1817,-0.1786) (1.6581,2.1817,-0.1810) (1.7453,2.1817,-0.1879) (1.8326,2.1817,-0.1992) (1.9199,2.1817,-0.2143) (2.0071,2.1817,-0.2327) (2.0944,2.1817,-0.2538) (2.1817,2.1817,-0.2767) (2.2689,2.1817,-0.3008) (2.3562,2.1817,-0.3252) (2.4435,2.1817,-0.3494) (2.5307,2.1817,-0.3728) (2.6180,2.1817,-0.3950) (2.7053,2.1817,-0.4157) (2.7925,2.1817,-0.4349) (2.8798,2.1817,-0.4527) (2.9671,2.1817,-0.4692) (3.0543,2.1817,-0.4848) (3.1416,2.1817,-0.5000) (3.2289,2.1817,-0.5152) (3.3161,2.1817,-0.5308) (3.4034,2.1817,-0.5473) (3.4907,2.1817,-0.5651) (3.5779,2.1817,-0.5843) (3.6652,2.1817,-0.6050) (3.7525,2.1817,-0.6272) (3.8397,2.1817,-0.6506) (3.9270,2.1817,-0.6748) (4.0143,2.1817,-0.6992) (4.1015,2.1817,-0.7233) (4.1888,2.1817,-0.7462) (4.2761,2.1817,-0.7673) (4.3633,2.1817,-0.7857) (4.4506,2.1817,-0.8008) (4.5379,2.1817,-0.8121) (4.6251,2.1817,-0.8190) (4.7124,2.1817,-0.8214) (4.7997,2.1817,-0.8190) (4.8869,2.1817,-0.8121) (4.9742,2.1817,-0.8008) (5.0615,2.1817,-0.7857) (5.1487,2.1817,-0.7673) (5.2360,2.1817,-0.7462) (5.3233,2.1817,-0.7233) (5.4105,2.1817,-0.6992) (5.4978,2.1817,-0.6748) (5.5851,2.1817,-0.6506) (5.6723,2.1817,-0.6272) (5.7596,2.1817,-0.6050) (5.8469,2.1817,-0.5843) (5.9341,2.1817,-0.5651) (6.0214,2.1817,-0.5473) (6.1087,2.1817,-0.5308) (6.1959,2.1817,-0.5152) (6.2832,2.1817,-0.5000)
(0.0000,2.2689,-0.5000) (0.0873,2.2689,-0.5129) (0.1745,2.2689,-0.5244) (0.2618,2.2689,-0.5333) (0.3491,2.2689,-0.5386) (0.4363,2.2689,-0.5393) (0.5236,2.2689,-0.5350) (0.6109,2.2689,-0.5256) (0.6981,2.2689,-0.5113) (0.7854,2.2689,-0.4927) (0.8727,2.2689,-0.4707) (0.9599,2.2689,-0.4465) (1.0472,2.2689,-0.4215) (1.1345,2.2689,-0.3970) (1.2217,2.2689,-0.3746) (1.3090,2.2689,-0.3556) (1.3963,2.2689,-0.3411) (1.4835,2.2689,-0.3321) (1.5708,2.2689,-0.3290) (1.6581,2.2689,-0.3321) (1.7453,2.2689,-0.3411) (1.8326,2.2689,-0.3556) (1.9199,2.2689,-0.3746) (2.0071,2.2689,-0.3970) (2.0944,2.2689,-0.4215) (2.1817,2.2689,-0.4465) (2.2689,2.2689,-0.4707) (2.3562,2.2689,-0.4927) (2.4435,2.2689,-0.5113) (2.5307,2.2689,-0.5256) (2.6180,2.2689,-0.5350) (2.7053,2.2689,-0.5393) (2.7925,2.2689,-0.5386) (2.8798,2.2689,-0.5333) (2.9671,2.2689,-0.5244) (3.0543,2.2689,-0.5129) (3.1416,2.2689,-0.5000) (3.2289,2.2689,-0.4871) (3.3161,2.2689,-0.4756) (3.4034,2.2689,-0.4667) (3.4907,2.2689,-0.4614) (3.5779,2.2689,-0.4607) (3.6652,2.2689,-0.4650) (3.7525,2.2689,-0.4744) (3.8397,2.2689,-0.4887) (3.9270,2.2689,-0.5073) (4.0143,2.2689,-0.5293) (4.1015,2.2689,-0.5535) (4.1888,2.2689,-0.5785) (4.2761,2.2689,-0.6030) (4.3633,2.2689,-0.6254) (4.4506,2.2689,-0.6444) (4.5379,2.2689,-0.6589) (4.6251,2.2689,-0.6679) (4.7124,2.2689,-0.6710) (4.7997,2.2689,-0.6679) (4.8869,2.2689,-0.6589) (4.9742,2.2689,-0.6444) (5.0615,2.2689,-0.6254) (5.1487,2.2689,-0.6030) (5.2360,2.2689,-0.5785) (5.3233,2.2689,-0.5535) (5.4105,2.2689,-0.5293) (5.4978,2.2689,-0.5073) (5.5851,2.2689,-0.4887) (5.6723,2.2689,-0.4744) (5.7596,2.2689,-0.4650) (5.8469,2.2689,-0.4607) (5.9341,2.2689,-0.4614) (6.0214,2.2689,-0.4667) (6.1087,2.2689,-0.4756) (6.1959,2.2689,-0.4871) (6.2832,2.2689,-0.5000)
(0.0000,2.3562,-0.5000) (0.0873,2.3562,-0.5432) (0.1745,2.3562,-0.5842) (0.2618,2.3562,-0.6207) (0.3491,2.3562,-0.6510) (0.4363,2.3562,-0.6736) (0.5236,2.3562,-0.6875) (0.6109,2.3562,-0.6924) (0.6981,2.3562,-0.6886) (0.7854,2.3562,-0.6768) (0.8727,2.3562,-0.6583) (0.9599,2.3562,-0.6347) (1.0472,2.3562,-0.6083) (1.1345,2.3562,-0.5809) (1.2217,2.3562,-0.5550) (1.3090,2.3562,-0.5324) (1.3963,2.3562,-0.5148) (1.4835,2.3562,-0.5038) (1.5708,2.3562,-0.5000) (1.6581,2.3562,-0.5038) (1.7453,2.3562,-0.5148) (1.8326,2.3562,-0.5324) (1.9199,2.3562,-0.5550) (2.0071,2.3562,-0.5809) (2.0944,2.3562,-0.6083) (2.1817,2.3562,-0.6347) (2.2689,2.3562,-0.6583) (2.3562,2.3562,-0.6768) (2.4435,2.3562,-0.6886) (2.5307,2.3562,-0.6924) (2.6180,2.3562,-0.6875) (2.7053,2.3562,-0.6736) (2.7925,2.3562,-0.6510) (2.8798,2.3562,-0.6207) (2.9671,2.3562,-0.5842) (3.0543,2.3562,-0.5432) (3.1416,2.3562,-0.5000) (3.2289,2.3562,-0.4568) (3.3161,2.3562,-0.4158) (3.4034,2.3562,-0.3793) (3.4907,2.3562,-0.3490) (3.5779,2.3562,-0.3264) (3.6652,2.3562,-0.3125) (3.7525,2.3562,-0.3076) (3.8397,2.3562,-0.3114) (3.9270,2.3562,-0.3232) (4.0143,2.3562,-0.3417) (4.1015,2.3562,-0.3653) (4.1888,2.3562,-0.3917) (4.2761,2.3562,-0.4191) (4.3633,2.3562,-0.4450) (4.4506,2.3562,-0.4676) (4.5379,2.3562,-0.4852) (4.6251,2.3562,-0.4962) (4.7124,2.3562,-0.5000) (4.7997,2.3562,-0.4962) (4.8869,2.3562,-0.4852) (4.9742,2.3562,-0.4676) (5.0615,2.3562,-0.4450) (5.1487,2.3562,-0.4191) (5.2360,2.3562,-0.3917) (5.3233,2.3562,-0.3653) (5.4105,2.3562,-0.3417) (5.4978,2.3562,-0.3232) (5.5851,2.3562,-0.3114) (5.6723,2.3562,-0.3076) (5.7596,2.3562,-0.3125) (5.8469,2.3562,-0.3264) (5.9341,2.3562,-0.3490) (6.0214,2.3562,-0.3793) (6.1087,2.3562,-0.4158) (6.1959,2.3562,-0.4568) (6.2832,2.3562,-0.5000)
(0.0000,2.4435,-0.5000) (0.0873,2.4435,-0.5723) (0.1745,2.4435,-0.6414) (0.2618,2.4435,-0.7045) (0.3491,2.4435,-0.7588) (0.4363,2.4435,-0.8026) (0.5236,2.4435,-0.8343) (0.6109,2.4435,-0.8534) (0.6981,2.4435,-0.8602) (0.7854,2.4435,-0.8555) (0.8727,2.4435,-0.8410) (0.9599,2.4435,-0.8189) (1.0472,2.4435,-0.7917) (1.1345,2.4435,-0.7624) (1.2217,2.4435,-0.7336) (1.3090,2.4435,-0.7081) (1.3963,2.4435,-0.6881) (1.4835,2.4435,-0.6754) (1.5708,2.4435,-0.6710) (1.6581,2.4435,-0.6754) (1.7453,2.4435,-0.6881) (1.8326,2.4435,-0.7081) (1.9199,2.4435,-0.7336) (2.0071,2.4435,-0.7624) (2.0944,2.4435,-0.7917) (2.1817,2.4435,-0.8189) (2.2689,2.4435,-0.8410) (2.3562,2.4435,-0.8555) (2.4435,2.4435,-0.8602) (2.5307,2.4435,-0.8534) (2.6180,2.4435,-0.8343) (2.7053,2.4435,-0.8026) (2.7925,2.4435,-0.7588) (2.8798,2.4435,-0.7045) (2.9671,2.4435,-0.6414) (3.0543,2.4435,-0.5723) (3.1416,2.4435,-0.5000) (3.2289,2.4435,-0.4277) (3.3161,2.4435,-0.3586) (3.4034,2.4435,-0.2955) (3.4907,2.4435,-0.2412) (3.5779,2.4435,-0.1974) (3.6652,2.4435,-0.1657) (3.7525,2.4435,-0.1466) (3.8397,2.4435,-0.1398) (3.9270,2.4435,-0.1445) (4.0143,2.4435,-0.1590) (4.1015,2.4435,-0.1811) (4.1888,2.4435,-0.2083) (4.2761,2.4435,-0.2376) (4.3633,2.4435,-0.2664) (4.4506,2.4435,-0.2919) (4.5379,2.4435,-0.3119) (4.6251,2.4435,-0.3246) (4.7124,2.4435,-0.3290) (4.7997,2.4435,-0.3246) (4.8869,2.4435,-0.3119) (4.9742,2.4435,-0.2919) (5.0615,2.4435,-0.2664) (5.1487,2.4435,-0.2376) (5.2360,2.4435,-0.2083) (5.3233,2.4435,-0.1811) (5.4105,2.4435,-0.1590) (5.4978,2.4435,-0.1445) (5.5851,2.4435,-0.1398) (5.6723,2.4435,-0.1466) (5.7596,2.4435,-0.1657) (5.8469,2.4435,-0.1974) (5.9341,2.4435,-0.2412) (6.0214,2.4435,-0.2955) (6.1087,2.4435,-0.3586) (6.1959,2.4435,-0.4277) (6.2832,2.4435,-0.5000)
(0.0000,2.5307,-0.5000) (0.0873,2.5307,-0.5964) (0.1745,2.5307,-0.6891) (0.2618,2.5307,-0.7743) (0.3491,2.5307,-0.8489) (0.4363,2.5307,-0.9105) (0.5236,2.5307,-0.9574) (0.6109,2.5307,-0.9889) (0.6981,2.5307,-1.0050) (0.7854,2.5307,-1.0070) (0.8727,2.5307,-0.9966) (0.9599,2.5307,-0.9765) (1.0472,2.5307,-0.9496) (1.1345,2.5307,-0.9194) (1.2217,2.5307,-0.8890) (1.3090,2.5307,-0.8616) (1.3963,2.5307,-0.8400) (1.4835,2.5307,-0.8262) (1.5708,2.5307,-0.8214) (1.6581,2.5307,-0.8262) (1.7453,2.5307,-0.8400) (1.8326,2.5307,-0.8616) (1.9199,2.5307,-0.8890) (2.0071,2.5307,-0.9194) (2.0944,2.5307,-0.9496) (2.1817,2.5307,-0.9765) (2.2689,2.5307,-0.9966) (2.3562,2.5307,-1.0070) (2.4435,2.5307,-1.0050) (2.5307,2.5307,-0.9889) (2.6180,2.5307,-0.9574) (2.7053,2.5307,-0.9105) (2.7925,2.5307,-0.8489) (2.8798,2.5307,-0.7743) (2.9671,2.5307,-0.6891) (3.0543,2.5307,-0.5964) (3.1416,2.5307,-0.5000) (3.2289,2.5307,-0.4036) (3.3161,2.5307,-0.3109) (3.4034,2.5307,-0.2257) (3.4907,2.5307,-0.1511) (3.5779,2.5307,-0.0895) (3.6652,2.5307,-0.0426) (3.7525,2.5307,-0.0111) (3.8397,2.5307,0.0050) (3.9270,2.5307,0.0070) (4.0143,2.5307,-0.0034) (4.1015,2.5307,-0.0235) (4.1888,2.5307,-0.0504) (4.2761,2.5307,-0.0806) (4.3633,2.5307,-0.1110) (4.4506,2.5307,-0.1384) (4.5379,2.5307,-0.1600) (4.6251,2.5307,-0.1738) (4.7124,2.5307,-0.1786) (4.7997,2.5307,-0.1738) (4.8869,2.5307,-0.1600) (4.9742,2.5307,-0.1384) (5.0615,2.5307,-0.1110) (5.1487,2.5307,-0.0806) (5.2360,2.5307,-0.0504) (5.3233,2.5307,-0.0235) (5.4105,2.5307,-0.0034) (5.4978,2.5307,0.0070) (5.5851,2.5307,0.0050) (5.6723,2.5307,-0.0111) (5.7596,2.5307,-0.0426) (5.8469,2.5307,-0.0895) (5.9341,2.5307,-0.1511) (6.0214,2.5307,-0.2257) (6.1087,2.5307,-0.3109) (6.1959,2.5307,-0.4036) (6.2832,2.5307,-0.5000)
(0.0000,2.6180,-0.5000) (0.0873,2.6180,-0.6126) (0.1745,2.6180,-0.7210) (0.2618,2.6180,-0.8212) (0.3491,2.6180,-0.9096) (0.4363,2.6180,-0.9836) (0.5236,2.6180,-1.0413) (0.6109,2.6180,-1.0817) (0.6981,2.6180,-1.1050) (0.7854,2.6180,-1.1124) (0.8727,2.6180,-1.1058) (0.9599,2.6180,-1.0881) (1.0472,2.6180,-1.0625) (1.1345,2.6180,-1.0326) (1.2217,2.6180,-1.0021) (1.3090,2.6180,-0.9743) (1.3963,2.6180,-0.9522) (1.4835,2.6180,-0.9379) (1.5708,2.6180,-0.9330) (1.6581,2.6180,-0.9379) (1.7453,2.6180,-0.9522) (1.8326,2.6180,-0.9743) (1.9199,2.6180,-1.0021) (2.0071,2.6180,-1.0326) (2.0944,2.6180,-1.0625) (2.1817,2.6180,-1.0881) (2.2689,2.6180,-1.1058) (2.3562,2.6180,-1.1124) (2.4435,2.6180,-1.1050) (2.5307,2.6180,-1.0817) (2.6180,2.6180,-1.0413) (2.7053,2.6180,-0.9836) (2.7925,2.6180,-0.9096) (2.8798,2.6180,-0.8212) (2.9671,2.6180,-0.7210) (3.0543,2.6180,-0.6126) (3.1416,2.6180,-0.5000) (3.2289,2.6180,-0.3874) (3.3161,2.6180,-0.2790) (3.4034,2.6180,-0.1788) (3.4907,2.6180,-0.0904) (3.5779,2.6180,-0.0164) (3.6652,2.6180,0.0413) (3.7525,2.6180,0.0817) (3.8397,2.6180,0.1050) (3.9270,2.6180,0.1124) (4.0143,2.6180,0.1058) (4.1015,2.6180,0.0881) (4.1888,2.6180,0.0625) (4.2761,2.6180,0.0326) (4.3633,2.6180,0.0021) (4.4506,2.6180,-0.0257) (4.5379,2.6180,-0.0478) (4.6251,2.6180,-0.0621) (4.7124,2.6180,-0.0670) (4.7997,2.6180,-0.0621) (4.8869,2.6180,-0.0478) (4.9742,2.6180,-0.0257) (5.0615,2.6180,0.0021) (5.1487,2.6180,0.0326) (5.2360,2.6180,0.0625) (5.3233,2.6180,0.0881) (5.4105,2.6180,0.1058) (5.4978,2.6180,0.1124) (5.5851,2.6180,0.1050) (5.6723,2.6180,0.0817) (5.7596,2.6180,0.0413) (5.8469,2.6180,-0.0164) (5.9341,2.6180,-0.0904) (6.0214,2.6180,-0.1788) (6.1087,2.6180,-0.2790) (6.1959,2.6180,-0.3874) (6.2832,2.6180,-0.5000)
(0.0000,2.7053,-0.5000) (0.0873,2.7053,-0.6186) (0.1745,2.7053,-0.7329) (0.2618,2.7053,-0.8388) (0.3491,2.7053,-0.9328) (0.4363,2.7053,-1.0120) (0.5236,2.7053,-1.0745) (0.6109,2.7053,-1.1194) (0.6981,2.7053,-1.1467) (0.7854,2.7053,-1.1577) (0.8727,2.7053,-1.1543) (0.9599,2.7053,-1.1393) (1.0472,2.7053,-1.1160) (1.1345,2.7053,-1.0880) (1.2217,2.7053,-1.0589) (1.3090,2.7053,-1.0323) (1.3963,2.7053,-1.0109) (1.4835,2.7053,-0.9972) (1.5708,2.7053,-0.9924) (1.6581,2.7053,-0.9972) (1.7453,2.7053,-1.0109) (1.8326,2.7053,-1.0323) (1.9199,2.7053,-1.0589) (2.0071,2.7053,-1.0880) (2.0944,2.7053,-1.1160) (2.1817,2.7053,-1.1393) (2.2689,2.7053,-1.1543) (2.3562,2.7053,-1.1577) (2.4435,2.7053,-1.1467) (2.5307,2.7053,-1.1194) (2.6180,2.7053,-1.0745) (2.7053,2.7053,-1.0120) (2.7925,2.7053,-0.9328) (2.8798,2.7053,-0.8388) (2.9671,2.7053,-0.7329) (3.0543,2.7053,-0.6186) (3.1416,2.7053,-0.5000) (3.2289,2.7053,-0.3814) (3.3161,2.7053,-0.2671) (3.4034,2.7053,-0.1612) (3.4907,2.7053,-0.0672) (3.5779,2.7053,0.0120) (3.6652,2.7053,0.0745) (3.7525,2.7053,0.1194) (3.8397,2.7053,0.1467) (3.9270,2.7053,0.1577) (4.0143,2.7053,0.1543) (4.1015,2.7053,0.1393) (4.1888,2.7053,0.1160) (4.2761,2.7053,0.0880) (4.3633,2.7053,0.0589) (4.4506,2.7053,0.0323) (4.5379,2.7053,0.0109) (4.6251,2.7053,-0.0028) (4.7124,2.7053,-0.0076) (4.7997,2.7053,-0.0028) (4.8869,2.7053,0.0109) (4.9742,2.7053,0.0323) (5.0615,2.7053,0.0589) (5.1487,2.7053,0.0880) (5.2360,2.7053,0.1160) (5.3233,2.7053,0.1393) (5.4105,2.7053,0.1543) (5.4978,2.7053,0.1577) (5.5851,2.7053,0.1467) (5.6723,2.7053,0.1194) (5.7596,2.7053,0.0745) (5.8469,2.7053,0.0120) (5.9341,2.7053,-0.0672) (6.0214,2.7053,-0.1612) (6.1087,2.7053,-0.2671) (6.1959,2.7053,-0.3814) (6.2832,2.7053,-0.5000)
(0.0000,2.7925,-0.5000) (0.0873,2.7925,-0.6133) (0.1745,2.7925,-0.7226) (0.2618,2.7925,-0.8240) (0.3491,2.7925,-0.9142) (0.4363,2.7925,-0.9906) (0.5236,2.7925,-1.0514) (0.6109,2.7925,-1.0956) (0.6981,2.7925,-1.1235) (0.7854,2.7925,-1.1359) (0.8727,2.7925,-1.1348) (0.9599,2.7925,-1.1227) (1.0472,2.7925,-1.1026) (1.1345,2.7925,-1.0780) (1.2217,2.7925,-1.0522) (1.3090,2.7925,-1.0283) (1.3963,2.7925,-1.0091) (1.4835,2.7925,-0.9967) (1.5708,2.7925,-0.9924) (1.6581,2.7925,-0.9967) (1.7453,2.7925,-1.0091) (1.8326,2.7925,-1.0283) (1.9199,2.7925,-1.0522) (2.0071,2.7925,-1.0780) (2.0944,2.7925,-1.1026) (2.1817,2.7925,-1.1227) (2.2689,2.7925,-1.1348) (2.3562,2.7925,-1.1359) (2.4435,2.7925,-1.1235) (2.5307,2.7925,-1.0956) (2.6180,2.7925,-1.0514) (2.7053,2.7925,-0.9906) (2.7925,2.7925,-0.9142) (2.8798,2.7925,-0.8240) (2.9671,2.7925,-0.7226) (3.0543,2.7925,-0.6133) (3.1416,2.7925,-0.5000) (3.2289,2.7925,-0.3867) (3.3161,2.7925,-0.2774) (3.4034,2.7925,-0.1760) (3.4907,2.7925,-0.0858) (3.5779,2.7925,-0.0094) (3.6652,2.7925,0.0514) (3.7525,2.7925,0.0956) (3.8397,2.7925,0.1235) (3.9270,2.7925,0.1359) (4.0143,2.7925,0.1348) (4.1015,2.7925,0.1227) (4.1888,2.7925,0.1026) (4.2761,2.7925,0.0780) (4.3633,2.7925,0.0522) (4.4506,2.7925,0.0283) (4.5379,2.7925,0.0091) (4.6251,2.7925,-0.0033) (4.7124,2.7925,-0.0076) (4.7997,2.7925,-0.0033) (4.8869,2.7925,0.0091) (4.9742,2.7925,0.0283) (5.0615,2.7925,0.0522) (5.1487,2.7925,0.0780) (5.2360,2.7925,0.1026) (5.3233,2.7925,0.1227) (5.4105,2.7925,0.1348) (5.4978,2.7925,0.1359) (5.5851,2.7925,0.1235) (5.6723,2.7925,0.0956) (5.7596,2.7925,0.0514) (5.8469,2.7925,-0.0094) (5.9341,2.7925,-0.0858) (6.0214,2.7925,-0.1760) (6.1087,2.7925,-0.2774) (6.1959,2.7925,-0.3867) (6.2832,2.7925,-0.5000)
(0.0000,2.8798,-0.5000) (0.0873,2.8798,-0.5968) (0.1745,2.8798,-0.6902) (0.2618,2.8798,-0.7770) (0.3491,2.8798,-0.8544) (0.4363,2.8798,-0.9201) (0.5236,2.8798,-0.9726) (0.6109,2.8798,-1.0112) (0.6981,2.8798,-1.0360) (0.7854,2.8798,-1.0477) (0.8727,2.8798,-1.0479) (0.9599,2.8798,-1.0388) (1.0472,2.8798,-1.0229) (1.1345,2.8798,-1.0030) (1.2217,2.8798,-0.9820) (1.3090,2.8798,-0.9625) (1.3963,2.8798,-0.9467) (1.4835,2.8798,-0.9365) (1.5708,2.8798,-0.9330) (1.6581,2.8798,-0.9365) (1.7453,2.8798,-0.9467) (1.8326,2.8798,-0.9625) (1.9199,2.8798,-0.9820) (2.0071,2.8798,-1.0030) (2.0944,2.8798,-1.0229) (2.1817,2.8798,-1.0388) (2.2689,2.8798,-1.0479) (2.3562,2.8798,-1.0477) (2.4435,2.8798,-1.0360) (2.5307,2.8798,-1.0112) (2.6180,2.8798,-0.9726) (2.7053,2.8798,-0.9201) (2.7925,2.8798,-0.8544) (2.8798,2.8798,-0.7770) (2.9671,2.8798,-0.6902) (3.0543,2.8798,-0.5968) (3.1416,2.8798,-0.5000) (3.2289,2.8798,-0.4032) (3.3161,2.8798,-0.3098) (3.4034,2.8798,-0.2230) (3.4907,2.8798,-0.1456) (3.5779,2.8798,-0.0799) (3.6652,2.8798,-0.0274) (3.7525,2.8798,0.0112) (3.8397,2.8798,0.0360) (3.9270,2.8798,0.0477) (4.0143,2.8798,0.0479) (4.1015,2.8798,0.0388) (4.1888,2.8798,0.0229) (4.2761,2.8798,0.0030) (4.3633,2.8798,-0.0180) (4.4506,2.8798,-0.0375) (4.5379,2.8798,-0.0533) (4.6251,2.8798,-0.0635) (4.7124,2.8798,-0.0670) (4.7997,2.8798,-0.0635) (4.8869,2.8798,-0.0533) (4.9742,2.8798,-0.0375) (5.0615,2.8798,-0.0180) (5.1487,2.8798,0.0030) (5.2360,2.8798,0.0229) (5.3233,2.8798,0.0388) (5.4105,2.8798,0.0479) (5.4978,2.8798,0.0477) (5.5851,2.8798,0.0360) (5.6723,2.8798,0.0112) (5.7596,2.8798,-0.0274) (5.8469,2.8798,-0.0799) (5.9341,2.8798,-0.1456) (6.0214,2.8798,-0.2230) (6.1087,2.8798,-0.3098) (6.1959,2.8798,-0.4032) (6.2832,2.8798,-0.5000)
(0.0000,2.9671,-0.5000) (0.0873,2.9671,-0.5706) (0.1745,2.9671,-0.6387) (0.2618,2.9671,-0.7021) (0.3491,2.9671,-0.7586) (0.4363,2.9671,-0.8068) (0.5236,2.9671,-0.8453) (0.6109,2.9671,-0.8739) (0.6981,2.9671,-0.8923) (0.7854,2.9671,-0.9014) (0.8727,2.9671,-0.9021) (0.9599,2.9671,-0.8960) (1.0472,2.9671,-0.8849) (1.1345,2.9671,-0.8710) (1.2217,2.9671,-0.8561) (1.3090,2.9671,-0.8423) (1.3963,2.9671,-0.8311) (1.4835,2.9671,-0.8239) (1.5708,2.9671,-0.8214) (1.6581,2.9671,-0.8239) (1.7453,2.9671,-0.8311) (1.8326,2.9671,-0.8423) (1.9199,2.9671,-0.8561) (2.0071,2.9671,-0.8710) (2.0944,2.9671,-0.8849) (2.1817,2.9671,-0.8960) (2.2689,2.9671,-0.9021) (2.3562,2.9671,-0.9014) (2.4435,2.9671,-0.8923) (2.5307,2.9671,-0.8739) (2.6180,2.9671,-0.8453) (2.7053,2.9671,-0.8068) (2.7925,2.9671,-0.7586) (2.8798,2.9671,-0.7021) (2.9671,2.9671,-0.6387) (3.0543,2.9671,-0.5706) (3.1416,2.9671,-0.5000) (3.2289,2.9671,-0.4294) (3.3161,2.9671,-0.3613) (3.4034,2.9671,-0.2979) (3.4907,2.9671,-0.2414) (3.5779,2.9671,-0.1932) (3.6652,2.9671,-0.1547) (3.7525,2.9671,-0.1261) (3.8397,2.9671,-0.1077) (3.9270,2.9671,-0.0986) (4.0143,2.9671,-0.0979) (4.1015,2.9671,-0.1040) (4.1888,2.9671,-0.1151) (4.2761,2.9671,-0.1290) (4.3633,2.9671,-0.1439) (4.4506,2.9671,-0.1577) (4.5379,2.9671,-0.1689) (4.6251,2.9671,-0.1761) (4.7124,2.9671,-0.1786) (4.7997,2.9671,-0.1761) (4.8869,2.9671,-0.1689) (4.9742,2.9671,-0.1577) (5.0615,2.9671,-0.1439) (5.1487,2.9671,-0.1290) (5.2360,2.9671,-0.1151) (5.3233,2.9671,-0.1040) (5.4105,2.9671,-0.0979) (5.4978,2.9671,-0.0986) (5.5851,2.9671,-0.1077) (5.6723,2.9671,-0.1261) (5.7596,2.9671,-0.1547) (5.8469,2.9671,-0.1932) (5.9341,2.9671,-0.2414) (6.0214,2.9671,-0.2979) (6.1087,2.9671,-0.3613) (6.1959,2.9671,-0.4294) (6.2832,2.9671,-0.5000)
(0.0000,3.0543,-0.5000) (0.0873,3.0543,-0.5372) (0.1745,3.0543,-0.5731) (0.2618,3.0543,-0.6065) (0.3491,3.0543,-0.6364) (0.4363,3.0543,-0.6618) (0.5236,3.0543,-0.6822) (0.6109,3.0543,-0.6973) (0.6981,3.0543,-0.7072) (0.7854,3.0543,-0.7121) (0.8727,3.0543,-0.7126) (0.9599,3.0543,-0.7096) (1.0472,3.0543,-0.7039) (1.1345,3.0543,-0.6967) (1.2217,3.0543,-0.6890) (1.3090,3.0543,-0.6819) (1.3963,3.0543,-0.6761) (1.4835,3.0543,-0.6723) (1.5708,3.0543,-0.6710) (1.6581,3.0543,-0.6723) (1.7453,3.0543,-0.6761) (1.8326,3.0543,-0.6819) (1.9199,3.0543,-0.6890) (2.0071,3.0543,-0.6967) (2.0944,3.0543,-0.7039) (2.1817,3.0543,-0.7096) (2.2689,3.0543,-0.7126) (2.3562,3.0543,-0.7121) (2.4435,3.0543,-0.7072) (2.5307,3.0543,-0.6973) (2.6180,3.0543,-0.6822) (2.7053,3.0543,-0.6618) (2.7925,3.0543,-0.6364) (2.8798,3.0543,-0.6065) (2.9671,3.0543,-0.5731) (3.0543,3.0543,-0.5372) (3.1416,3.0543,-0.5000) (3.2289,3.0543,-0.4628) (3.3161,3.0543,-0.4269) (3.4034,3.0543,-0.3935) (3.4907,3.0543,-0.3636) (3.5779,3.0543,-0.3382) (3.6652,3.0543,-0.3178) (3.7525,3.0543,-0.3027) (3.8397,3.0543,-0.2928) (3.9270,3.0543,-0.2879) (4.0143,3.0543,-0.2874) (4.1015,3.0543,-0.2904) (4.1888,3.0543,-0.2961) (4.2761,3.0543,-0.3033) (4.3633,3.0543,-0.3110) (4.4506,3.0543,-0.3181) (4.5379,3.0543,-0.3239) (4.6251,3.0543,-0.3277) (4.7124,3.0543,-0.3290) (4.7997,3.0543,-0.3277) (4.8869,3.0543,-0.3239) (4.9742,3.0543,-0.3181) (5.0615,3.0543,-0.3110) (5.1487,3.0543,-0.3033) (5.2360,3.0543,-0.2961) (5.3233,3.0543,-0.2904) (5.4105,3.0543,-0.2874) (5.4978,3.0543,-0.2879) (5.5851,3.0543,-0.2928) (5.6723,3.0543,-0.3027) (5.7596,3.0543,-0.3178) (5.8469,3.0543,-0.3382) (5.9341,3.0543,-0.3636) (6.0214,3.0543,-0.3935) (6.1087,3.0543,-0.4269) (6.1959,3.0543,-0.4628) (6.2832,3.0543,-0.5000)
(0.0000,3.1416,-0.5000) (0.0873,3.1416,-0.5000) (0.1745,3.1416,-0.5000) (0.2618,3.1416,-0.5000) (0.3491,3.1416,-0.5000) (0.4363,3.1416,-0.5000) (0.5236,3.1416,-0.5000) (0.6109,3.1416,-0.5000) (0.6981,3.1416,-0.5000) (0.7854,3.1416,-0.5000) (0.8727,3.1416,-0.5000) (0.9599,3.1416,-0.5000) (1.0472,3.1416,-0.5000) (1.1345,3.1416,-0.5000) (1.2217,3.1416,-0.5000) (1.3090,3.1416,-0.5000) (1.3963,3.1416,-0.5000) (1.4835,3.1416,-0.5000) (1.5708,3.1416,-0.5000) (1.6581,3.1416,-0.5000) (1.7453,3.1416,-0.5000) (1.8326,3.1416,-0.5000) (1.9199,3.1416,-0.5000) (2.0071,3.1416,-0.5000) (2.0944,3.1416,-0.5000) (2.1817,3.1416,-0.5000) (2.2689,3.1416,-0.5000) (2.3562,3.1416,-0.5000) (2.4435,3.1416,-0.5000) (2.5307,3.1416,-0.5000) (2.6180,3.1416,-0.5000) (2.7053,3.1416,-0.5000) (2.7925,3.1416,-0.5000) (2.8798,3.1416,-0.5000) (2.9671,3.1416,-0.5000) (3.0543,3.1416,-0.5000) (3.1416,3.1416,-0.5000) (3.2289,3.1416,-0.5000) (3.3161,3.1416,-0.5000) (3.4034,3.1416,-0.5000) (3.4907,3.1416,-0.5000) (3.5779,3.1416,-0.5000) (3.6652,3.1416,-0.5000) (3.7525,3.1416,-0.5000) (3.8397,3.1416,-0.5000) (3.9270,3.1416,-0.5000) (4.0143,3.1416,-0.5000) (4.1015,3.1416,-0.5000) (4.1888,3.1416,-0.5000) (4.2761,3.1416,-0.5000) (4.3633,3.1416,-0.5000) (4.4506,3.1416,-0.5000) (4.5379,3.1416,-0.5000) (4.6251,3.1416,-0.5000) (4.7124,3.1416,-0.5000) (4.7997,3.1416,-0.5000) (4.8869,3.1416,-0.5000) (4.9742,3.1416,-0.5000) (5.0615,3.1416,-0.5000) (5.1487,3.1416,-0.5000) (5.2360,3.1416,-0.5000) (5.3233,3.1416,-0.5000) (5.4105,3.1416,-0.5000) (5.4978,3.1416,-0.5000) (5.5851,3.1416,-0.5000) (5.6723,3.1416,-0.5000) (5.7596,3.1416,-0.5000) (5.8469,3.1416,-0.5000) (5.9341,3.1416,-0.5000) (6.0214,3.1416,-0.5000) (6.1087,3.1416,-0.5000) (6.1959,3.1416,-0.5000) (6.2832,3.1416,-0.5000)
};
\addplot3[only marks,mark=x,mark size=4pt,corange,very thick] coordinates {(0.7854,2.7053,0)};
\end{axis}
\end{tikzpicture}

\caption{Paysage $F_1(\gamma,\beta)$ du QUBO jouet (calcul exact sur une grille $73\times37$). La croix marque le minimum, $F_1\approx-1{,}16$, à comparer à l'énergie moyenne $-0{,}5$ de l'état initial et au fondamental $-2$. Le paysage est périodique ($2\pi$ en $\gamma$ car les énergies sont entières, $\pi$ en $\beta$) et non convexe : l'optimiseur classique a besoin de plusieurs points de départ.}\label{fig:paysage-qaoa}
\end{figure}

\begin{figure}[htbp]
\centering
% figure générée par tools/make_figures.py : jouet_proba_p1
\begin{tikzpicture}
\begin{axis}[width=0.8\linewidth,height=5cm,ybar,bar width=12pt,axis lines=left,ymin=0,ymax=0.4,
  xtick={0,...,7},xticklabels={$|000\rangle$,$|100\rangle$,$|010\rangle$,$|110\rangle$,$|001\rangle$,$|101\rangle$,$|011\rangle$,$|111\rangle$},
  ylabel={probabilité},xticklabel style={font=\small},
  nodes near coords={},enlarge x limits=0.08]
\addplot[fill=cblue,draw=none] coordinates {(0,0.04482) (1,0.1365) (2,0.2777) (3,0.04093) (4,0.1365) (5,0.3104) (6,0.04093) (7,0.01218)};
\addplot[corange,thick,sharp plot,mark=none,update limits=false] coordinates {(-0.5,0.125) (7.5,0.125)};
\end{axis}
\end{tikzpicture}

\caption{Distribution de mesure après QAOA $p=1$ aux angles optimaux. La solution $\ket{101}$ sort avec probabilité $0{,}31$, contre $1/8=0{,}125$ pour la superposition uniforme (trait orange). À $p=2$ et $p=3$, elle monte à $0{,}54$ puis $0{,}67$.}\label{fig:jouet-proba}
\end{figure}

\section{Résultats sur notre QUBO aéro (9 qubits)}
Le QUBO de traînée du chapitre~\ref{ch:qubo} (512 chaînes) a été traité par QAOA simulé exactement, avec une optimisation COBYLA. Pour $p=1$, on part de 8 points de départ aléatoires ; pour $p\ge2$, on part de l'interpolation des angles optimaux de $p-1$ (stratégie INTERP).

\begin{center}
\begin{tabular}{cccccc}
\toprule
$p$ & CNOT & $P(\text{optimum})$ & $P(\text{5 meilleurs})$ & ratio $r$ & shots moyens pour voir l'optimum\\
\midrule
1 & 72 & 0,018 & 0,088 & 0,867 & 56\\
2 & 144 & 0,062 & 0,218 & 0,933 & 16\\
3 & 216 & 0,097 & 0,302 & 0,952 & 10\\
4 & 288 & 0,127 & 0,384 & 0,966 & 8\\
5 & 360 & 0,142 & 0,436 & 0,972 & 7\\
\bottomrule
\end{tabular}
\end{center}
Le ratio d'approximation est $r=1-\expval{\tilde E}$, où $\tilde E=(E-E_{\min})/(E_{\max}-E_{\min})\in[0,1]$ ; il vaut 1 si l'on mesure toujours l'optimum. Le hasard uniforme donne $P(\text{optimum})=1/512\approx0{,}002$.

\begin{figure}[htbp]
\centering
% figure générée par tools/make_figures.py : qaoa_proba_rang
\begin{tikzpicture}
\begin{axis}[width=0.85\linewidth,height=5.8cm,axis lines=left,xmin=0.5,xmax=25.5,ymin=0,ymax=0.155,
  xlabel={rang de l'énergie QUBO (1 = minimum)},ylabel={probabilité de mesure},
  legend style={draw=none,font=\small,at={(0.98,0.98)},anchor=north east}]
\addplot[cblue,thick,mark=*,mark size=1.6pt] coordinates {(1,0.01771) (2,0.02594) (3,0.01501) (4,0.01886) (5,0.01022) (6,0.011) (7,0.03356) (8,0.0103) (9,0.007724) (10,0.01746) (11,0.007663) (12,0.01215) (13,0.01348) (14,0.005126) (15,0.009068) (16,0.004641) (17,0.02513) (18,0.007766) (19,0.007132) (20,0.00968) (21,0.01745) (22,0.004169) (23,0.005312) (24,0.005923) (25,0.0109)};\addlegendentry{$p=1$}
\addplot[corange,thick,mark=*,mark size=1.6pt] coordinates {(1,0.09662) (2,0.06801) (3,0.05075) (4,0.02962) (5,0.05687) (6,0.1012) (7,0.03086) (8,0.02036) (9,0.02548) (10,0.04532) (11,0.04281) (12,0.03432) (13,0.01705) (14,0.01967) (15,0.02236) (16,0.007503) (17,0.03259) (18,0.01114) (19,0.01694) (20,0.007158) (21,0.008521) (22,0.007297) (23,0.01132) (24,0.02091) (25,0.01149)};\addlegendentry{$p=3$}
\addplot[cgreen,thick,mark=*,mark size=1.6pt] coordinates {(1,0.142) (2,0.1067) (3,0.06809) (4,0.04338) (5,0.07616) (6,0.128) (7,0.04926) (8,0.02372) (9,0.02495) (10,0.02896) (11,0.06437) (12,0.03374) (13,0.01058) (14,0.0238) (15,0.01275) (16,0.005626) (17,0.00962) (18,0.005358) (19,0.02331) (20,0.004033) (21,0.005568) (22,0.003729) (23,0.005821) (24,0.01116) (25,0.002006)};\addlegendentry{$p=5$}
\addplot[cmuted,dashed,thick] coordinates {(0.5,0.00195) (25.5,0.00195)};\addlegendentry{hasard : $1/512$}
\end{axis}
\end{tikzpicture}

\caption{Probabilité de mesurer chaque chaîne, classée par énergie QUBO croissante (25 premières). QAOA concentre la probabilité sur le bas du spectre, d'autant plus que $p$ est grand. Pour $p=1$, le rang 2 dépasse le rang 1 : les deux profils sont quasi dégénérés (§\ref{sec:notre-qubo}).}\label{fig:qaoa-rang}
\end{figure}

\begin{figure}[htbp]
\centering
\begin{minipage}{0.49\linewidth}\centering% figure générée par tools/make_figures.py : qaoa_convergence_p
\begin{tikzpicture}
\begin{axis}[width=0.7\linewidth,height=5.4cm,axis lines=left,xmin=0.7,xmax=5.3,ymin=0,ymax=1,xtick={1,...,5},
  xlabel={profondeur $p$},ylabel={probabilité / ratio},
  legend style={draw=none,font=\small,at={(0.02,0.62)},anchor=west}]
\addplot[cblue,thick,mark=*] coordinates {(1,0.01771) (2,0.06215) (3,0.09662) (4,0.1269) (5,0.142)};\addlegendentry{$P(\text{optimum})$}
\addplot[corange,thick,mark=square*] coordinates {(1,0.08774) (2,0.2185) (3,0.3019) (4,0.3836) (5,0.4363)};\addlegendentry{$P(\text{5 meilleurs})$}
\addplot[cgreen,thick,mark=triangle*] coordinates {(1,0.8673) (2,0.9325) (3,0.9515) (4,0.9658) (5,0.972)};\addlegendentry{ratio d'approximation $r$}
\end{axis}
\end{tikzpicture}
\end{minipage}\hfill
\begin{minipage}{0.49\linewidth}\centering% figure générée par tools/make_figures.py : qaoa_angles
\begin{tikzpicture}
\begin{axis}[width=0.7\linewidth,height=5.2cm,axis lines=left,xmin=0.7,xmax=5.3,xtick={1,...,5},
  xlabel={couche $k$},ylabel={angle (rad)},
  legend style={draw=none,font=\small,at={(0.02,0.98)},anchor=north west}]
\addplot[cblue,thick,mark=*] coordinates {(1,0.5558) (2,1.153) (3,1.706) (4,2.239) (5,2.563)};\addlegendentry{$\gamma_k$ (coût)}
\addplot[corange,thick,mark=square*] coordinates {(1,0.6973) (2,0.5187) (3,0.3131) (4,0.2204) (5,0.1343)};\addlegendentry{$|\beta_k|$ (mélangeur)}
\end{axis}
\end{tikzpicture}
\end{minipage}
\caption{À gauche : qualité en fonction de la profondeur. À droite : angles optimaux pour $p=5$. $\gamma_k$ croît et $|\beta_k|$ décroît le long du circuit : c'est le profil d'un \emph{recuit} ($s$ croissant), retrouvé par l'optimiseur sans qu'on le lui impose.}\label{fig:qaoa-p}
\end{figure}

\begin{piege}[ (échelle de l'énergie)]
$\gamma$ multiplie $E$ dans une phase. Avec $E$ en counts de traînée (de $57$ à $117$), une variation de $\gamma$ de $0{,}1$ fait tourner les phases de plusieurs radians : le paysage $F_1$ devient une « tôle ondulée » où l'optimiseur se perd. Sans normalisation, on obtient $P(\text{optimum})=0{,}0045$ à $p=1$ ; en centrant et réduisant $E$ (\texttt{normalize}), on obtient $0{,}018$. Il faut toujours normaliser le QUBO avant QAOA.
\end{piege}

\section{Limites à connaître}
\begin{itemize}
\item \textbf{Optimisation classique difficile} : le paysage est non convexe, avec de nombreux minima locaux. Il faut des stratégies d'initialisation (INTERP, transfert d'angles d'une instance à une autre).
\item \textbf{Estimation bruitée} : sur un vrai processeur, $F_p$ est estimé avec $N$ shots, donc avec un bruit statistique en $1/\sqrt N$, plus le bruit matériel.
\item \textbf{Profondeur} : $2|E|p$ CNOT. Avec un graphe complet, la profondeur explose avec $n$, ce qui est incompatible avec le matériel bruité actuel au-delà de quelques couches.
\item \textbf{Plateaux arides} : pour des circuits profonds et larges, le gradient de $F_p$ peut devenir exponentiellement petit.
\item \textbf{Pas d'avantage prouvé} : QAOA à $p$ fixe n'est pas connu pour battre les meilleurs algorithmes classiques sur des QUBO généraux.
\end{itemize}

\exercices
\begin{exercice}
Généraliser l'exemple à un qubit à $H_C=h\,Z$ avec $h>0$. Montrer que $F_1(\gamma,\beta)=h\,\sin(2h\gamma)\,\sin2\beta$. Où sont les minima ? Que se passe-t-il pour le paysage si $h$ est très grand ? (Lien avec le piège de l'échelle.)
\end{exercice}
\begin{exercice}
Pour notre QUBO aéro (9 qubits, 36 couplages), compter les portes d'un circuit QAOA à $p=3$ : $H$, $R_Z$, CNOT, $R_X$. Avec $P(\text{optimum})=0{,}097$, combien de shots faut-il pour observer l'optimum au moins une fois avec une probabilité de 99~\% ?
\end{exercice}

\chapter{Simuler un ordinateur quantique : méthodes et outils}\label{ch:outils}
\iterbanner{V0}

\section{Trois niveaux de simulation}
\begin{center}
\begin{tabular}{p{3.2cm}p{5.6cm}p{5.2cm}}
\toprule
Niveau & Ce qu'on calcule & Usage\\
\midrule
\textbf{Vecteur d'état exact} & les $2^n$ amplitudes ; probabilités et $\expval{H_C}$ exactes & mise au point, petits $n$ ($\lesssim30$) ; \emph{notre cas}\\
\textbf{Échantillonnage} (shots) & $N$ chaînes tirées selon $|\psi_x|^2$ ; estimateurs bruités & reproduit le comportement statistique d'une vraie machine\\
\textbf{Simulation bruitée} & matrice densité ou trajectoires, avec modèles d'erreur des portes & prévoir la dégradation sur matériel réel\\
\bottomrule
\end{tabular}
\end{center}
Le premier niveau est le plus simple et le plus instructif. Pour un QUBO, il est même trivial, car le coût est diagonal.

\section{Notre implémentation (NumPy, \texttt{quantique/qaoa.py})}
Toute la mécanique de QAOA tient en une vingtaine de lignes :
\begin{lstlisting}
def apply_mixer(psi, beta, n):
    """exp(-i beta sum_k X_k) |psi> : une rotation R_X(2 beta) par qubit."""
    psi = psi.reshape((2,) * n)                 # un axe par qubit
    c, s = np.cos(beta), -1j * np.sin(beta)
    for ax in range(n):
        psi = c * psi + s * np.flip(psi, axis=ax)   # X échange 0 et 1 sur cet axe
    return psi.reshape(-1)

def qaoa_state(gammas, betas, E, n):
    psi = np.full(2**n, 2 ** (-n / 2), dtype=complex)   # |+>^n
    for g, b in zip(gammas, betas):
        psi = psi * np.exp(-1j * g * E)          # coût : H_C diagonal -> phases
        psi = apply_mixer(psi, b, n)             # mélangeur
    return psi

def expectation(params, E, n):                    # F_p = <psi| H_C |psi>
    p = len(params) // 2
    psi = qaoa_state(params[:p], params[p:], E, n)
    return float(np.real(np.vdot(psi, E * psi)))
\end{lstlisting}
Ici \texttt{E} est le vecteur des $2^n$ énergies du QUBO, calculé une fois pour toutes (\texttt{brute\_force}). L'optimisation des angles est confiée à \texttt{scipy.optimize.minimize(method="COBYLA")}. Le script \texttt{quantique/experiences.py} produit tous les résultats chiffrés de ce cours.

\section{Les bibliothèques prévues par la note d'intention}
Les extraits suivants sont des \textbf{esquisses non testées} : ces bibliothèques ne sont pas installées sur la machine de rédaction, et leurs API évoluent vite. Il faut les valider sur notre QUBO en comparant au résultat NumPy, qui sert de référence.

\paragraph{D-Wave Ocean (recuit simulé, ou matériel réel).}
\begin{lstlisting}
import dimod
from dwave.samplers import SimulatedAnnealingSampler   # ancien paquet : neal
bqm = dimod.BinaryQuadraticModel.from_qubo(
    {(i, j): Q[i, j] for i in range(9) for j in range(i, 9) if Q[i, j] != 0}, offset=c)
res = SimulatedAnnealingSampler().sample(bqm, num_reads=200)
print(res.first.sample, res.first.energy)
# matériel : from dwave.system import DWaveSampler, EmbeddingComposite
#            EmbeddingComposite(DWaveSampler()).sample(bqm, num_reads=1000)  (plongement automatique)
\end{lstlisting}

\paragraph{Qiskit Optimization (QAOA).}
\begin{lstlisting}
from qiskit_optimization import QuadraticProgram
from qiskit_optimization.algorithms import MinimumEigenOptimizer
from qiskit_algorithms import QAOA
from qiskit_algorithms.optimizers import COBYLA
from qiskit.primitives import Sampler          # selon la version : StatevectorSampler (V2)
qp = QuadraticProgram()
for i in range(9):
    qp.binary_var(f"b{i}")
qp.minimize(constant=c, linear=[Q[i, i] for i in range(9)],
            quadratic={(f"b{i}", f"b{j}"): Q[i, j] for i in range(9) for j in range(i + 1, 9)})
H_C, offset = qp.to_ising()                    # hamiltonien d'Ising (SparsePauliOp)
qaoa = QAOA(sampler=Sampler(), optimizer=COBYLA(), reps=3)
print(MinimumEigenOptimizer(qaoa).solve(qp))
\end{lstlisting}

\paragraph{PennyLane.} Le module \texttt{qml.qaoa} fournit \texttt{cost\_layer} et \texttt{mixer\_layer}, avec une différentiation automatique des angles. Il est utile pour essayer des optimiseurs à gradient.

\section{Pièges pratiques}\label{sec:pieges-outils}
\begin{itemize}
\item \textbf{Ordre des bits} : Qiskit affiche $q_{n-1}\cdots q_0$. Une chaîne « \texttt{000010000} » ne se lit pas dans le même sens que notre \texttt{encode}. Il faut tester sur un profil connu.
\item \textbf{Conventions de signe et de facteur} : $e^{-i\beta X}$ ou $e^{-i\beta X/2}$, mélangeur $\sum X_k$ ou $-\sum X_k$. Les angles optimaux changent de signe ou de facteur 2 d'une bibliothèque à l'autre, mais la qualité de la solution ne change pas.
\item \textbf{Normalisation} de l'énergie avant QAOA (ch.~\ref{ch:qaoa}).
\item \textbf{Périodicité} : $\gamma$ n'est périodique que si les énergies sont commensurables. Il ne faut pas restreindre arbitrairement la recherche à $[0,\pi]$.
\item \textbf{Graines aléatoires} : l'optimiseur et l'échantillonnage sont stochastiques. On répète les expériences et on rapporte médiane et dispersion.
\end{itemize}

\section{Le matériel réel : bruit et ère NISQ}
Un qubit physique perd son information en un temps fini (\textbf{décohérence}), et chaque porte a une probabilité d'erreur non nulle. Ordres de grandeur actuels : quelques dixièmes de pourcent d'erreur par porte à deux qubits. Un circuit de quelques centaines de CNOT donne déjà un résultat fortement dégradé : nos 216 CNOT à $p=3$ sont à la limite. Des techniques d'\emph{atténuation d'erreurs} existent ; la \emph{correction} d'erreurs complète demande beaucoup plus de qubits qu'aujourd'hui. C'est pourquoi la note d'intention vise des \textbf{simulateurs} et une \textbf{analyse des ressources} plutôt qu'une exécution matérielle.

\exercices
\begin{exercice}
Quelle mémoire faut-il pour simuler exactement 30 qubits ? 40 qubits ? Un balayage de QAOA coûte $O(n\,2^n)$ opérations par couche : comparer le coût pour $n=9$ et $n=30$.
\end{exercice}
\begin{exercice}
On estime $F_p=\expval{H_C}$ avec $N$ shots. L'énergie normalisée a un écart-type $\sigma_E\approx1$ sur les shots. Combien de shots faut-il pour connaître $F_p$ à $\pm0{,}01$ (un écart-type) ? Pourquoi cela compte-t-il pour l'optimiseur classique ?
\end{exercice}

\chapter{De bout en bout : la formulation quantique du hackathon}\label{ch:projet}
\iterbanner{V0}
\begin{evolution}{V1}{la boucle devient bi-objectif}Ce chapitre décrit la boucle de la V0, sur la seule traînée. En V1, elle a été étendue aux deux objectifs (traînée et SER), avec les contraintes codées dans le QUBO (13 qubits : réduction de Rosenberg et variables d'écart), et comparée à NSGA-II : voir le cours \emph{Optimisation bi-objectif traînée / SER}, chapitres 6 et 7.\end{evolution}

\section{La chaîne complète}
\begin{figure}[htbp]
\centering
\begin{tikzpicture}[node distance=5mm and 6mm]
\node[flow] (g) {Grille $8^3$\\ 9 bits};
\node[flow,right=of g] (s1) {Phase 1\\ aéro : $C_D$};
\node[flow,below=of s1] (s2) {Phase 2\\ SER};
\node[flow,right=12mm of s1,yshift=-5mm] (m) {Modèles\\ quadratiques\\ $\widehat{C_D}$, $\widehat{\rm SER}$};
\node[flow,right=of m,fill=corange!10] (qubo) {QUBO\\ $w\widehat{C_D}+(1-w)\widehat{\rm SER}$\\ (ch.~\ref{ch:qubo})};
\node[flow,right=of qubo,fill=cblue!10] (solve) {QAOA /\\ recuit\\ (ch.~\ref{ch:recuit}, \ref{ch:qaoa})};
\node[flow,below=12mm of qubo] (front) {front de Pareto\\ des profils \emph{simulés}};
\draw[arr] (g) -- (s1); \draw[arr] (g) |- (s2);
\draw[arr] (s1) -- (m); \draw[arr] (s2) -- (m);
\draw[arr] (m) -- (qubo); \draw[arr] (qubo) -- (solve);
\draw[arr] (solve) |- node[pos=0.25,right,font=\small] {candidats} (front);
\draw[arr] (front) -| node[pos=0.75,left,font=\small,align=right] {simuler les\\ nouveaux} (s2.south);
\end{tikzpicture}
\caption{La formulation quantique de la note d'intention. Le front final ne contient que des profils réellement simulés : le modèle quadratique ne sert qu'à \emph{proposer} des candidats.}\label{fig:chaine}
\end{figure}

\section{La boucle itérative, et pourquoi elle est indispensable}
\begin{tcolorbox}[coursbox,colback=white,colframe=cink,title={Algorithme (une seule fonction objectif, ici $C_D$).}]
\begin{enumerate}
\item Simuler $n_0$ profils (plan d'expériences initial).
\item Ajuster le modèle quadratique (10 coefficients) sur tous les profils simulés.
\item Construire le QUBO et le résoudre : QAOA, recuit, ou énumération quand $n$ est petit.
\item Simuler les $k$ meilleurs candidats \emph{non encore simulés}.
\item Recommencer jusqu'à stabilisation, ou épuisement du budget de simulations.
\end{enumerate}
\end{tcolorbox}
On a vu au chapitre~\ref{ch:qubo} qu'un modèle ajusté \emph{une fois} sur toute la base désigne un profil seulement 38\textsuperscript{e} sur 512. La boucle corrige ce défaut, car elle raffine le modèle là où sont les bons candidats.

\begin{figure}[htbp]
\centering
% figure générée par tools/make_figures.py : boucle_convergence
\begin{tikzpicture}
\begin{axis}[width=0.85\linewidth,height=5.6cm,axis lines=left,xmin=10,xmax=60,ymin=49,ymax=62,
  xlabel={nombre de profils simulés},ylabel={meilleur $C_D$ trouvé (counts), médiane},
  legend style={draw=none,font=\small,at={(0.98,0.98)},anchor=north east}]
\addplot[cblue,very thick] coordinates {(10,61) (11,60.92) (12,60.85) (13,60.79) (14,58.9) (15,58.51) (16,57) (17,54.78) (18,54.78) (19,53.41) (20,52.79) (21,52.62) (22,52.2) (23,51.16) (24,50.85) (25,50.53) (26,50.53) (27,50.53) (28,50.53) (29,50.53) (30,50.53) (31,50.53) (32,50.53) (33,50.53) (34,50.53) (35,50.53) (36,50.53) (37,50.53) (38,50.53) (39,50.53) (40,50.53) (41,50.53) (42,50.53) (43,50.53) (44,50.53) (45,50.53) (46,50.53) (47,50.53) (48,50.53) (49,50.53) (50,50.53) (51,50.53) (52,50.53) (53,50.53) (54,50.53) (55,50.53)};\addlegendentry{boucle : modèle quadratique, puis QUBO}
\addplot[corange,thick] coordinates {(10,61.1) (11,60) (12,58.9) (13,58.8) (14,58.65) (15,58.65) (16,58.51) (17,57.76) (18,57) (19,56.99) (20,56.99) (21,56.72) (22,55.95) (23,55.95) (24,55.16) (25,55.1) (26,55.1) (27,53.81) (28,53.81) (29,53.81) (30,53.81) (31,53.7) (32,53.7) (33,53.7) (34,53.06) (35,53.06) (36,53.03) (37,53.03) (38,53.03) (39,53.03) (40,53.03) (41,53.03) (42,53.03) (43,52.98) (44,52.79) (45,52.79) (46,52.79) (47,52.79) (48,52.79) (49,52.62) (50,52.62) (51,52.62) (52,52.62) (53,52.62) (54,52.62) (55,52.62) (56,52.62) (57,52.41) (58,52.41) (59,52.2) (60,52.2)};\addlegendentry{tirage aléatoire}
\addplot[cmuted,dashed] coordinates {(10,50.533) (60,50.533)};\addlegendentry{vrai minimum (50.5)}
\end{axis}
\end{tikzpicture}

\caption{Meilleur $C_D$ trouvé en fonction du nombre de profils simulés (médiane sur 50 tirages du plan initial : $n_0=10$ profils aléatoires, puis $k=3$ candidats par itération). Le QUBO est résolu ici par énumération, qui donne la même réponse qu'un QAOA ou un recuit parfaitement réussi. La boucle trouve le vrai minimum en \textbf{25 simulations} (médiane ; quartiles 22 et 34), contre 256 en moyenne pour un tirage au hasard sur 512 profils.}\label{fig:boucle}
\end{figure}

\begin{projet}[ : ce qu'il reste à construire]
Ce résultat ne porte que sur $C_D$. Il faut encore :
\begin{itemize}
\item ajouter la SER (phase~2) ;
\item traiter le bi-objectif (section suivante) ;
\item remplacer l'énumération par QAOA et le recuit ;
\item compter les simulations avec le même compteur que NSGA-II (\texttt{AeroSimulator}).
\end{itemize}
Les briques sont dans \texttt{Code/quantique/} : \texttt{qubo.py}, \texttt{qaoa.py}, \texttt{recuit.py}, \texttt{experiences.py}.
\end{projet}

\section{Deux objectifs, un seul QUBO}
Un QUBO n'a qu'une énergie. Pour obtenir un front de Pareto $(C_D,\ {\rm SER})$ :
\begin{itemize}
\item \textbf{Somme pondérée} : minimiser $w\,\tilde C_D+(1-w)\,\widetilde{\rm SER}$ (objectifs normalisés) pour plusieurs $w\in[0,1]$. Elle reste quadratique, donc c'est un QUBO, mais elle ne trouve que les points de la \emph{partie convexe} du front.
\item \textbf{$\varepsilon$-contrainte} : minimiser $\tilde C_D$ sous $\widetilde{\rm SER}\le\varepsilon$, avec des bits d'écart et une pénalité (ch.~\ref{ch:qubo}). Elle trouve tout le front, au prix de qubits supplémentaires, et la pénalité doit rester de degré 2.
\end{itemize}
\begin{figure}[htbp]
\centering
\begin{tikzpicture}
\begin{axis}[width=0.7\linewidth,height=5.6cm,axis lines=left,xmin=0,xmax=5,ymin=0,ymax=5,xtick=\empty,ytick=\empty,
  xlabel={$\tilde C_D$},ylabel={$\widetilde{\rm SER}$}]
\addplot[only marks,mark=*,cmuted,mark size=1.5pt] coordinates {(1.6,4.6)(2.4,4.1)(3.2,3.5)(2.7,2.9)(3.8,2.6)(4.3,1.9)(3.4,4.4)(4.6,3.6)(2.2,3.6)(4.0,4.0)};
\addplot[cblue,thick,mark=*,mark size=2.2pt] coordinates {(0.6,4.4)(1.2,2.4)(2.9,1.1)(4.4,0.5)};
\addplot[corange,only marks,mark=square*,mark size=2.6pt] coordinates {(2.3,2.0)};
\draw[cblue,dashed] (axis cs:1.2,2.4) -- (axis cs:2.9,1.1);
\node[corange,font=\small,anchor=west] at (axis cs:2.4,2.1) {non supporté};
\node[cblue,font=\small,anchor=west] at (axis cs:0.65,4.6) {front de Pareto};
\draw[-{Stealth},cmuted] (axis cs:4.4,4.6) -- (axis cs:3.9,4.1) node[pos=0,anchor=south,font=\small] {$w$ fixé : lignes de niveau};
\end{axis}
\end{tikzpicture}
\caption{Le point orange est Pareto-optimal (aucun autre profil n'est meilleur sur les deux objectifs), mais il est situé \emph{au-dessus} du segment qui relie ses voisins : aucune pondération $w$ ne le rend minimal. La somme pondérée ne le trouvera jamais.}\label{fig:pareto}
\end{figure}

\section{La comparaison avec NSGA-II}
Pour que la comparaison soit honnête :
\begin{itemize}
\item même grille de profils et même point de vol ;
\item même \textbf{compteur de simulations uniques} ;
\item même critère d'arrêt, fixé à l'avance : le front exact retrouvé, ou un pourcentage de son hypervolume ;
\item plusieurs graines aléatoires, avec médiane et dispersion.
\end{itemize}
Le front exact est connu par énumération des 512 profils. C'est ce qui rend la démonstration rigoureuse.

\section{Ressources et passage à l'échelle}
Avec $d$ variables de forme codées sur $k$ bits chacune, en encodage binaire :
\begin{center}
\begin{tabular}{lccc}
\toprule
& notre cas ($d=3$, $k=3$) & $d=10$, $k=4$ & $d=50$, $k=4$\\
\midrule
qubits $n=dk$ & 9 & 40 & 200\\
couplages (graphe complet) $\binom n2$ & 36 & 780 & 19\,900\\
CNOT par couche QAOA ($2\binom n2$) & 72 & 1\,560 & 39\,800\\
coefficients du modèle ($\binom{d+2}2$) & 10 & 66 & 1\,326\\
simulation par vecteur d'état & triviale & $\approx16$~To & impossible\\
\bottomrule
\end{tabular}
\end{center}
Deux enseignements à présenter :
\begin{itemize}
\item le modèle quadratique en variables physiques demande peu de coefficients, donc peu de simulations, alors qu'un QUBO libre sur les bits en demanderait $\binom n2+n+1$ ;
\item le graphe complet est le vrai obstacle matériel : profondeur QAOA en $O(n^2)$ par couche, et plongement coûteux sur un recuit quantique.
\end{itemize}

\section{Messages pour le jury}
\begin{enumerate}
\item La chaîne complète, de la physique au QUBO puis à QAOA et au recuit, est construite et \emph{vérifiée} sur un cas résolu exactement.
\item La qualité se mesure en \textbf{nombre de simulations coûteuses} pour retrouver le front exact, et non en temps de calcul quantique.
\item Le goulot d'étranglement est la \emph{fidélité du modèle quadratique}, pas la résolution du QUBO : la boucle itérative est essentielle.
\item Les ressources quantiques croissent en $O(n^2)$ portes par couche pour un QUBO dense. C'est l'argument à confronter au matériel actuel (NISQ).
\end{enumerate}

\exercices
\begin{exercice}
Trois profils ont pour objectifs normalisés $A=(0,\,1)$, $B=(0{,}6,\,0{,}6)$, $C=(1,\,0)$. (a) Sont-ils tous Pareto-optimaux ? (b) Montrer qu'aucun $w\in[0,1]$ ne fait de $B$ le minimum de $w\,x+(1-w)\,y$. (c) Quelle méthode permet de trouver $B$ ?
\end{exercice}
\begin{exercice}
Un cas industriel a 20 variables de forme à 16 niveaux chacune. Nombre de qubits en binaire ? Nombre de CNOT par couche de QAOA si le QUBO est dense ? Combien de simulations au minimum pour ajuster le modèle quadratique en variables physiques ? Et un QUBO libre sur les bits ?
\end{exercice}

\chapter{V2 : que devient QAOA quand le problème grandit ?}\label{ch:montee}
\iterbanner{V2}

\begin{evolution}{V2}{pourquoi ce chapitre}La V1 a montré qu'à 13 qubits, le solveur du QUBO ne compte presque pas : le QUBO se résout exactement en une milliseconde, et QAOA fait aussi bien. La seule question quantique pertinente est donc celle du passage à l'échelle, que pose la note d'intention (« seuls les simulateurs et le nombre de variables changent »). Ce chapitre y répond par simulation exacte, jusqu'à 18 qubits.\end{evolution}

\section{Le protocole}
On garde \emph{le même} problème physique : le modèle quadratique de la somme pondérée traînée/SER ($w=0{,}5$), ajusté sur la base exhaustive V1. On affine seulement la grille : $b$ bits par variable, pour 3 variables, soit $n=3b$ qubits, de 6 à 18. Le QUBO grandit, mais sa structure reste la même : c'est exactement ce qui arriverait si l'on demandait une conception plus fine.

Pour chaque $n$ (\texttt{Code/quantique/montee\_en\_taille.py}) :
\begin{itemize}
\item les $2^n$ énergies sont calculées exactement, ce qui donne l'optimum de référence ;
\item QAOA ($p=1$ à 4) est simulé par vecteur d'état. Les angles de $p=1$ sont optimisés avec plusieurs départs ; ceux de $p\ge2$ partent du meilleur de deux initialisations : l'interpolation des angles optimaux de $p-1$ (méthode INTERP) et le transfert depuis la taille précédente, comme on le ferait sur une machine où chaque évaluation coûte cher ;
\item on mesure la probabilité de l'optimum, la probabilité de tomber dans le 1~\% meilleur des chaînes, et le rapport d'approximation $r=(E_{\max}-\expval{E})/(E_{\max}-E_{\min})$ ;
\item on compare au recuit simulé (part des lectures qui trouvent l'optimum, 100 balayages).
\end{itemize}

\section{Les résultats}
\begin{figure}[htbp]
\centering
% figure générée par tools/make_figures_v2.py
\begin{tikzpicture}
\begin{groupplot}[group style={group size=2 by 1,horizontal sep=1.6cm},width=0.48\linewidth,height=5.8cm,
  xlabel={nombre de qubits $n$},xtick={6,9,12,15,18},title style={font=\small},tick label style={font=\scriptsize},
  legend style={draw=none,font=\scriptsize}]
\nextgroupplot[ymode=log,title={probabilité de mesurer l'optimum},legend pos=south west]
\addplot[cblue,very thick,mark=*,mark size=1.3pt] coordinates {(6,0.079114) (9,0.016433) (12,0.0026563) (15,0.00038696) (18,5.1875e-05)};\addlegendentry{QAOA $p=1$}
\addplot[corange,very thick,mark=square*,mark size=1.3pt] coordinates {(6,0.13783) (9,0.043653) (12,0.0093765) (15,0.0015763) (18,0.00023126)};\addlegendentry{QAOA $p=2$}
\addplot[cviolet,very thick,mark=triangle*,mark size=1.6pt] coordinates {(6,0.17228) (9,0.077843) (12,0.020524) (15,0.0041741) (18,0.00068805)};\addlegendentry{QAOA $p=4$}
\addplot[cmuted,dashed,thick] coordinates {(6,0.015625) (9,0.0019531) (12,0.00024414) (15,3.0518e-05) (18,3.8147e-06)};\addlegendentry{hasard, $2^{-n}$}
\nextgroupplot[title={qualité de l'échantillon},ymin=0,ymax=1,legend pos=south west]
\addplot[corange,very thick,mark=square*,mark size=1.3pt] coordinates {(6,0.13783) (9,0.189) (12,0.26423) (15,0.28981) (18,0.29489)};\addlegendentry{$P$(1~\% meilleur), $p=2$}
\addplot[cviolet,very thick,mark=triangle*,mark size=1.6pt] coordinates {(6,0.17228) (9,0.32037) (12,0.50451) (15,0.57663) (18,0.59078)};\addlegendentry{$P$(1~\% meilleur), $p=4$}
\addplot[cgreen,very thick,mark=*,mark size=1.3pt] coordinates {(6,0.9572) (9,0.96404) (12,0.95918) (15,0.9523) (18,0.94742)};\addlegendentry{rapport d'approx., $p=4$}
\addplot[cmuted,very thick,dashed,mark=o,mark size=1.3pt] coordinates {(6,0.56) (9,0.58) (12,0.2) (15,0.08) (18,0)};\addlegendentry{recuit : lectures optimales}
\end{groupplot}
\end{tikzpicture}

\caption{Montée en taille de 6 à 18 qubits, sur le même problème physique. À gauche, la probabilité de mesurer l'optimum, en échelle logarithmique. À droite, des indicateurs plus robustes : probabilité d'une chaîne du 1~\% meilleur, rapport d'approximation, et réussite du recuit simulé.}\label{fig:montee}
\end{figure}

\begin{center}
\begin{tabular}{rrrrrr}
\toprule
$n$ & CNOT par couche & $P(\text{opt})$, $p=2$ & $P(\text{opt})$, $p=4$ & 1~\% meilleur, $p=4$ & hasard \\
\midrule
6 & 30 & 0,14 & 0,17 & 0,17 & 0,016 \\
9 & 72 & 0,044 & 0,078 & 0,32 & 0,002 \\
12 & 132 & 0,0094 & 0,021 & 0,50 & $2{,}4\cdot10^{-4}$ \\
15 & 210 & 0,0016 & 0,0042 & 0,58 & $3{,}1\cdot10^{-5}$ \\
18 & 306 & 0,0002 & 0,0007 & 0,59 & $3{,}8\cdot10^{-6}$ \\
\bottomrule
\end{tabular}
\end{center}

\paragraph{1. Trouver l'optimum devient exponentiellement difficile.} $P(\text{opt})$ est divisée par 3 à 7 tous les 3 qubits, en $2^{-0{,}7\,n}$ environ entre 6 et 18 qubits. QAOA garde un avantage sur le hasard (facteur 9 à 180 selon $p$ et $n$), mais cet avantage ne croît pas assez vite : le nombre de tirs nécessaires pour voir l'optimum croît, lui aussi, exponentiellement (environ 4\,000 tirs à 18 qubits avec $p=4$).

\paragraph{2. Mais l'échantillon reste bon.} La probabilité de tomber dans le 1~\% meilleur se \emph{stabilise} : environ 0,3 avec $p=2$, environ 0,59 avec $p=4$, et le rapport d'approximation reste au-dessus de 0,92. Quand le problème grandit, l'optimum exact devient une aiguille parmi de nombreuses chaînes presque aussi bonnes (la grille plus fine crée beaucoup de voisins quasi équivalents). Pour notre boucle, qui cherche des \emph{candidats} à simuler et non l'optimum exact du QUBO, c'est la bonne mesure.

\paragraph{3. La profondeur aide, à condition de bien initialiser les angles.} Avec l'initialisation INTERP, chaque couche supplémentaire améliore l'échantillon : à 18 qubits, la part de mesures dans le 1~\% meilleur passe de 0,29 ($p=2$) à 0,59 ($p=4$).

\begin{evolution}{V2}{correction en cours d'itération}Une première version de cette étude initialisait les angles de $p\ge3$ au hasard, ou par transfert seul : $p=3$ ne faisait alors pas mieux que $p=2$, et l'on avait conclu que la profondeur n'aidait pas. C'était un défaut de l'optimiseur classique, pas de QAOA : le paysage des angles devient plat et rugueux quand $p$ croît, et un départ aléatoire s'y perd. L'interpolation des angles optimaux de $p-1$ (Zhou et al., 2020) corrige le problème. Les résultats de la première version sont conservés dans \texttt{Code/resultats/montee\_en\_taille\_init\_aleatoire.json}.\end{evolution}

\paragraph{4. Le recuit simulé, réglé tel quel, s'effondre aussi.} Avec 100 balayages, il trouve l'optimum dans 58~\% des lectures à 9 qubits, mais plus jamais à 18. Comme pour QAOA, cela traduit surtout la multiplication des quasi-optima ; un recuit plus long, ou un critère « 1~\% meilleur », le remettrait à niveau.

\begin{resultat}
\begin{itemize}
\item Sur notre problème, QAOA reste meilleur que le hasard, mais trouver l'\emph{optimum exact} coûte exponentiellement plus de tirs quand $n$ croît.
\item En revanche, QAOA produit de bons \emph{candidats} à toute taille testée : environ un tiers des mesures tombent dans le 1~\% meilleur. C'est ce que la boucle exploite.
\item Le coût en portes croît comme $n^2$ (graphe complet) : 306 CNOT par couche à 18 qubits, avant le routage sur une vraie puce.
\end{itemize}
\end{resultat}

\begin{projet}
Le message à porter devant le jury change de nature. On ne prétend pas que QAOA résout mieux le QUBO : à ces tailles, l'énumération et le recuit classiques suffisent. On montre comment évoluent qualité et coût quand le problème grandit, et que le rôle utile du solveur, dans la boucle, est de fournir des candidats de bonne qualité, ce que QAOA fait de façon stable.
\end{projet}

\exercices
\begin{exercice}
Le nombre de tirs nécessaires pour voir au moins une fois l'optimum avec 95~\% de chances vaut $N=\ln0{,}05/\ln(1-P_{\text{opt}})$. Le calculer à 15 et 18 qubits pour $p=4$ (tableau), puis l'extrapoler à $n=24$ en supposant que la tendance continue. Comparer au nombre de chaînes $2^{24}$.
\end{exercice}

\begin{pistes}
\begin{itemize}
\item \iterbadge{V2}~\textbf{Fait} : initialisation INTERP, qui rend les couches supplémentaires utiles. Reste à aller au-delà de $p=4$ et à comparer aux angles fixes de la littérature.
\item Remplacer « probabilité de l'optimum » par « probabilité d'un candidat utile » comme indicateur principal.
\item Ajouter le bruit (dépolarisation par CNOT) et le routage sur une topologie réelle pour estimer le coût sur machine.
\end{itemize}
\end{pistes}


\appendix
\chapter{Corrigés des exercices}\label{annexe:corriges}
\iterbanner{V0}

\section*{Chapitre~\ref{ch:intro} — Pourquoi le calcul quantique}
\textbf{Ex.~1.1.} $H\ket1=\tfrac1{\sqrt2}(\ket0-\ket1)=\ket-$ : probabilités $\tfrac12,\tfrac12$. Puis $H\ket-=\ket1$ : on lit 1 avec certitude. Les deux chemins vers $\ket0$ ont des amplitudes $+\tfrac12$ et $-\tfrac12$, qui s'annulent. Deux pièces donnent 50/50 quel que soit le bit de départ : les probabilités ne s'annulent jamais.

\textbf{Ex.~1.2.} L'état final contient bien 512 amplitudes, mais une mesure ne renvoie \emph{qu'une} chaîne de 9 bits, tirée au hasard selon les $|\psi_x|^2$. On n'a pas accès aux 512 valeurs de coût. Un passage de QAOA produit \emph{un} profil candidat ; l'algorithme est réussi si les bons profils ont une grande probabilité, et l'on fait de nombreux shots pour les observer.

\section*{Chapitre~\ref{ch:bases} — Les bases}
\textbf{Ex.~2.1.} $|z_1|^2=|z_2|^2=\tfrac12$, donc $|z_1|^2+|z_2|^2=1$. Mais $z_2=\tfrac1{\sqrt2}e^{i\pi/4}e^{i\pi}=-z_1$, donc $z_1+z_2=0$ et $|z_1+z_2|^2=0$. $|z_1+z_2|^2$ est le module carré d'une somme d'\emph{amplitudes} : l'interférence est totalement destructive. $|z_1|^2+|z_2|^2$ est une somme de \emph{probabilités}.

\textbf{Ex.~2.2.} $\cos^230^\circ\times\cos^2(90^\circ-30^\circ)=\tfrac34\times\tfrac14=\tfrac3{16}\approx0{,}19$. Sans le polariseur intermédiaire : $\cos^290^\circ=0$.

\section*{Chapitre~\ref{ch:qubit} — Le qubit}
\textbf{Ex.~3.1.} (a) $|\tfrac1{\sqrt2}|^2=|\tfrac i{\sqrt2}|^2=\tfrac12$. (b) $\cos\tfrac\theta2=\tfrac1{\sqrt2}\Rightarrow\theta=\tfrac\pi2$ ; $e^{i\varphi}=i\Rightarrow\varphi=\tfrac\pi2$ : c'est $\ket{+i}$, sur l'axe $+y$. (c) $\braket+\psi=\tfrac12(1+i)$, de module carré $\tfrac12$ ; de même pour $\ket-$ : 50/50 aussi.

\textbf{Ex.~3.2.} (a) $-\ket-=\tfrac1{\sqrt2}(-\ket0+\ket1)$, soit le vecteur $\tfrac1{\sqrt2}(-1,1)$. Comme $-1=e^{i\pi}$, c'est $\ket-$ multiplié par une phase globale. Pour tout état $\ket u$, $|\braket{u}{-\psi}|^2=|-1|^2|\braket u\psi|^2=|\braket u\psi|^2$ : aucune mesure ne les distingue. (b) $\ket+$ et $\ket-$ ne diffèrent que par le signe de la composante sur $\ket1$, c'est-à-dire par une phase \emph{relative} $\varphi=\pi$. Leur produit scalaire vaut $\braket+-=\tfrac12(1-1)=0$ : ils sont orthogonaux, donc parfaitement distinguables. Procédure : appliquer $H$ puis mesurer en base de calcul. $H\ket+=\ket0$ donne toujours 0, et $H\ket-=\ket1$ donne toujours 1. (c) $\ket+$ est une lumière polarisée linéairement à $45^\circ$, $\ket-$ à $135^\circ$. Un polariseur à $45^\circ$ transmet toujours la première et bloque toujours la seconde, alors qu'un polariseur horizontal transmet chacune une fois sur deux.

\section*{Chapitre~\ref{ch:algebre} — Observables et évolution}
\textbf{Ex.~4.1.} (a) $X\ket+=\tfrac1{\sqrt2}(\ket1+\ket0)=\ket+$ et $X\ket-=\tfrac1{\sqrt2}(\ket1-\ket0)=-\ket-$. (b) $\bra+Z\ket+=\tfrac12(1-1)=0$ et $\bra+X\ket+=1$ : le point $(x,y,z)=(1,0,0)$, sur l'équateur.

\textbf{Ex.~4.2.} (a) $H^\dagger=H$ (réelle symétrique) et $H^2=\tfrac12\begin{pmatrix}2&0\\0&2\end{pmatrix}=\Id$, donc $H^\dagger H=\Id$. (b) Calcul direct : $HZH=\tfrac12\begin{pmatrix}1&1\\1&-1\end{pmatrix}\begin{pmatrix}1&1\\-1&1\end{pmatrix}=\begin{pmatrix}0&1\\1&0\end{pmatrix}=X$, et $HXH=H(HZH)H=Z$. Comme $H$ conjugue $Z$ en $X$, elle envoie les vecteurs propres de $Z$ sur ceux de $X$. (c) $X^2=\Id$, donc $e^{-i\theta X}=\sum_k\frac{(-i\theta)^k}{k!}X^k=\big(\sum_{k\text{ pair}}\cdots\big)\Id+\big(\sum_{k\text{ impair}}\cdots\big)X=\cos\theta\,\Id-i\sin\theta\,X$.

\section*{Chapitre~\ref{ch:multi} — Plusieurs qubits}
\textbf{Ex.~5.1.} (a) $\tfrac12(\ket0+\ket1)\otimes(\ket0-\ket1)=\ket+\ket-$ : produit (test : $\psi_{00}\psi_{11}-\psi_{01}\psi_{10}=-\tfrac14+\tfrac14=0$). (b) $\psi_{00}\psi_{11}-\psi_{01}\psi_{10}=0-\tfrac1{\sqrt2}\cdot(-\tfrac1{\sqrt2})=\tfrac12\neq0$ : intriqué (état singulet de Bell).

\textbf{Ex.~5.2.} Dans l'ordre $x=0,1,2,3$, c'est-à-dire $(b_0,b_1)=(0,0),(1,0),(0,1),(1,1)$ :
\[Z_0=\mathrm{diag}(1,-1,1,-1),\quad Z_1=\mathrm{diag}(1,1,-1,-1),\quad Z_0Z_1=\mathrm{diag}(1,-1,-1,1),\]
\[\tfrac14(\Id-Z_0)(\Id-Z_1)=\tfrac14\mathrm{diag}(0\cdot0,\ 2\cdot0,\ 0\cdot2,\ 2\cdot2)=\mathrm{diag}(0,0,0,1).\]

\section*{Chapitre~\ref{ch:circuits} — Portes et circuits}
\textbf{Ex.~6.1.} $(H\otimes\Id)\ket{01}=H\ket0\otimes\ket1=\tfrac1{\sqrt2}(\ket{01}+\ket{11})$. La CNOT (contrôle $q_0$) laisse $\ket{01}$ inchangé (contrôle à 0) et envoie $\ket{11}$ sur $\ket{10}$ (contrôle à 1, la cible passe de 1 à 0). On obtient $\tfrac1{\sqrt2}(\ket{01}+\ket{10})$. Probabilités : $\tfrac12$ pour 01, $\tfrac12$ pour 10, 0 pour 00 et 11 : les deux bits sont toujours \emph{différents}. Intriqué : $\psi_{00}\psi_{11}-\psi_{01}\psi_{10}=0-\tfrac12\neq0$. C'est un autre état de Bell, noté $\ket{\Psi^+}$.

\textbf{Ex.~6.2.} CNOT : $\ket{b_i,b_j}\mapsto\ket{b_i,b_i\oplus b_j}$. $R_Z(\theta)$ sur le second qubit multiplie par $e^{-i\theta/2}$ si $b_i\oplus b_j=0$, et par $e^{+i\theta/2}$ sinon. La seconde CNOT rétablit $\ket{b_i,b_j}$. Or $b_i\oplus b_j=0\Leftrightarrow s_is_j=+1$ : la phase est $e^{-i\theta s_is_j/2}$, qui est bien la valeur propre de $e^{-i\theta Z\otimes Z/2}$ sur $\ket{b_ib_j}$.

\section*{Chapitre~\ref{ch:hamiltoniens} — Hamiltoniens}
\textbf{Ex.~7.1.} $e^{-i\gamma Z}=\mathrm{diag}(e^{-i\gamma},e^{i\gamma})$, donc $e^{-i\gamma Z}\ket+=\tfrac1{\sqrt2}(e^{-i\gamma}\ket0+e^{i\gamma}\ket1)=e^{-i\gamma}\tfrac1{\sqrt2}(\ket0+e^{2i\gamma}\ket1)$. C'est une rotation d'angle $2\gamma$ autour de $z$ (longitude $\varphi=2\gamma$). Les probabilités restent $\tfrac12,\tfrac12$.

\textbf{Ex.~7.2.} $-X_k$ a pour fondamental $\ket+$ (valeur propre $-1$). Les $X_k$ agissent sur des qubits différents et commutent, donc le fondamental de la somme est le produit des fondamentaux, $\ket+^{\otimes n}$, d'énergie $-n$. En développant, $\ket+^{\otimes n}=2^{-n/2}\sum_x\ket x$ : toutes les chaînes ont la même amplitude.

\section*{Chapitre~\ref{ch:qubo} — QUBO et Ising}
\textbf{Ex.~8.1.} $h_0=-\tfrac{-1}2-\tfrac14(2)=0$, $h_1=\tfrac12-\tfrac14(2+2)=-\tfrac12$, $h_2=0$ ; $J_{01}=J_{12}=\tfrac24=\tfrac12$ ; constante $=0+\tfrac12(-3)+\tfrac14(4)=-\tfrac12$. Pour $\mathbf s=(-1,1,-1)$ : $E=-\tfrac12\cdot1+\tfrac12(-1)(1)+\tfrac12(1)(-1)-\tfrac12=-2$. On retrouve bien $E(101)$.

\textbf{Ex.~8.2.} $3\times8=24$ qubits au lieu de 9. Pénalité pour $m$ : $P\big(\sum_{j=0}^7b_{mj}-1\big)^2=P\big(1-\sum_jb_{mj}+2\sum_{j<l}b_{mj}b_{ml}\big)$. Elle ajoute $\binom82=28$ couplages par variable, soit 84 au total, en plus de ceux du coût. Le binaire est nettement plus économe pour notre projet. Le one-hot ne se justifierait que si le paysage binaire posait problème au recuit.

\section*{Chapitre~\ref{ch:recuit} — Recuit}
\textbf{Ex.~9.1.} $e^{-1/2}\approx0{,}61$ ; $e^{-2}\approx0{,}14$ ; $e^{-10}\approx4{,}5\cdot10^{-5}$. À haute température, on explore presque librement, ce qui permet de franchir les barrières. À basse température, on ne fait plus que descendre. Refroidir progressivement passe de l'exploration à l'exploitation.

\textbf{Ex.~9.2.} $T\gg1/g_{\min}^2\sim2^n$ : le temps adiabatique croît comme l'énumération exhaustive. L'adiabatique n'apporte alors aucun avantage asymptotique ; tout dépend de la façon dont le gap se comporte sur la famille de problèmes considérée.

\section*{Chapitre~\ref{ch:qaoa} — QAOA}
\textbf{Ex.~10.1.} $e^{-i\gamma hZ}=R_Z(2h\gamma)$ fait tourner $\ket+$ de $2h\gamma$ autour de $z$ : $(\cos2h\gamma,\sin2h\gamma,0)$. $R_X(2\beta)$ donne ensuite $z=\sin2h\gamma\,\sin2\beta$, d'où $F_1=h\,\sin(2h\gamma)\sin2\beta$. Les minima, $F_1=-h$, sont atteints pour $\sin(2h\gamma)=\pm1$ et $\sin2\beta=\mp1$, par exemple $\gamma=\tfrac{\pi}{4h}$ et $\beta=-\tfrac\pi4$. Si $h$ est grand, la période en $\gamma$ ($\pi/h$) devient minuscule : le paysage oscille très vite et l'optimiseur se perd. C'est le piège de l'échelle, d'où la normalisation.

\textbf{Ex.~10.2.} $H$ : 9 ; par couche, 9 $R_Z$, 36 $R_{ZZ}$ (72 CNOT et 36 $R_Z$) et 9 $R_X$. À $p=3$ : 216 CNOT, $3\times(9+36)=135$ $R_Z$, 27 $R_X$, plus 9 mesures. Shots : $1-(1-0{,}097)^N\ge0{,}99\Leftrightarrow N\ge\ln0{,}01/\ln0{,}903\approx45{,}2$, donc 46 shots.

\section*{Chapitre~\ref{ch:outils} — Outils}
\textbf{Ex.~11.1.} $16\cdot2^{30}\approx17$~Go ; $16\cdot2^{40}\approx17{,}6$~To. Coût par couche : $9\cdot2^9\approx4{,}6\cdot10^3$ opérations contre $30\cdot2^{30}\approx3{,}2\cdot10^{10}$, soit environ 7 millions de fois plus.

\textbf{Ex.~11.2.} $\sigma_E/\sqrt N=0{,}01\Rightarrow N=10^4$ shots \emph{par évaluation} de $F_p$. Un optimiseur qui fait des centaines d'évaluations demande donc des millions de shots. Sur matériel réel, c'est le coût dominant ; et le bruit statistique peut tromper un optimiseur qui compare des valeurs très proches.

\section*{Chapitre~\ref{ch:projet} — Le projet}
\textbf{Ex.~12.1.} $N\approx3/P_{\text{opt}}$ pour $P$ petit : environ 710 tirs à 15 qubits ($P=0{,}0042$) et 4\,300 à 18 qubits ($P=0{,}0007$), soit un facteur 6 tous les 3 qubits. En prolongeant : environ $4\,300\times36\approx1{,}5\cdot10^5$ tirs à 24 qubits, encore 100 fois moins que les $2^{24}\approx1{,}7\cdot10^7$ chaînes. Mais l'écart diminue avec $n$, et il faut ajouter les milliers d'évaluations nécessaires pour régler les angles, chacune demandant elle-même des centaines de tirs.

\chapter{Aide-mémoire}\label{annexe:memo}
\iterbanner{V0}

\begin{center}
\renewcommand{\arraystretch}{1.3}
\begin{tabularx}{\linewidth}{>{\raggedright\arraybackslash}p{4.3cm}X}
\toprule
\textbf{État} & $\ket\psi=\sum_x\psi_x\ket x$, $\sum|\psi_x|^2=1$ ; phase globale sans effet\\
\textbf{Mesure} & $P(x)=|\psi_x|^2$ ; $\expval A=\bra\psi A\ket\psi$ ; $N$ shots : erreur en $1/\sqrt N$\\
\textbf{Évolution} & $U$ unitaire ; $e^{-iHt}$ ; $e^{-i\theta P}=\cos\theta\,\Id-i\sin\theta\,P$ pour $P$ de Pauli\\
\textbf{Pauli} & $X=\begin{psmallmatrix}0&1\\1&0\end{psmallmatrix}$, $Z=\begin{psmallmatrix}1&0\\0&-1\end{psmallmatrix}$ ; $Z\ket b=(1-2b)\ket b$\\
\textbf{Bits $\leftrightarrow$ spins} & $b=(1-s)/2$ ; opérateur $\hat b=(\Id-Z)/2$\\
\textbf{Portes} & $H$ ; $R_P(\theta)=e^{-i\theta P/2}$ ; $R_{ZZ}(\theta)=\mathrm{CNOT}\,(\Id\otimes R_Z(\theta))\,\mathrm{CNOT}$\\
\textbf{QUBO} & $\min_{\mathbf b}\mathbf b^\top Q\mathbf b+c$, $Q$ triangulaire supérieure, termes linéaires sur la diagonale\\
\textbf{Ising} & $J_{ij}=Q_{ij}/4$ ; $h_i=-Q_{ii}/2-\tfrac14\sum_{j\neq i}Q_{ij}$ ; cte $=c+\tfrac12\sum Q_{ii}+\tfrac14\sum_{i<j}Q_{ij}$\\
\textbf{Hamiltonien de coût} & $H_C=\sum h_iZ_i+\sum J_{ij}Z_iZ_j$ (diagonal) ; fondamental = solution\\
\textbf{Pénalité} & égalité : $P\,g^2$ ; inégalité : bits d'écart ; $g$ doit être \emph{linéaire}\\
\textbf{Recuit simulé} & accepter avec $\min(1,e^{-\Delta E/T})$ ; $T$ décroissante\\
\textbf{Adiabatique} & $H(s)=(1-s)(-\sum X_k)+sH_C$ ; durée $\gg1/g_{\min}^2$\\
\textbf{QAOA} & $\prod_k e^{-i\beta_kH_M}e^{-i\gamma_kH_C}\ket+^{\otimes n}$ ; minimiser $F_p=\expval{H_C}$ ; $2p$ angles\\
\textbf{Ressources QAOA} & $n$ qubits ; par couche $n$ $R_Z$, $|E|$ $R_{ZZ}$ ($2|E|$ CNOT), $n$ $R_X$\\
\textbf{Notre cas} & 9 qubits, 36 couplages, 72 CNOT/couche ; $P(\text{opt})$ : 0,018 ($p=1$) $\to$ 0,142 ($p=5$)\\
\textbf{Boucle itérative} & vrai minimum de $C_D$ en 25 simulations (médiane) contre 256 au hasard\\
\textbf{Pièges} & normaliser $E$ ; ordre des bits (Qiskit) ; dégénérescence $m=0$ ; pénalité de degré 4 ; somme pondérée = front convexe seulement\\
\bottomrule
\end{tabularx}
\end{center}

\paragraph{Glossaire.}
\begin{description}[style=nextline,leftmargin=2.5cm]
\item[Amplitude] coefficient complexe d'une chaîne dans l'état ; son module au carré est une probabilité.
\item[Shot] une exécution du circuit suivie d'une mesure : une chaîne de bits.
\item[Fondamental] état d'énergie minimale d'un hamiltonien.
\item[Gap] écart entre les deux niveaux d'énergie les plus bas.
\item[Mélangeur] hamiltonien $\sum X_k$ qui fait passer l'amplitude d'une chaîne à ses voisines.
\item[Plongement] représentation d'un graphe de couplages sur la connectivité limitée du matériel.
\item[NISQ] \emph{Noisy Intermediate-Scale Quantum} : matériel actuel, bruité, de quelques dizaines à milliers de qubits.
\end{description}

\paragraph{Références.}
\begin{itemize}
\item M. A. Nielsen, I. L. Chuang, \emph{Quantum Computation and Quantum Information}, Cambridge UP, 2010, ch.~1, 2 et 4.
\item E. Farhi, J. Goldstone, S. Gutmann, « A Quantum Approximate Optimization Algorithm », arXiv:1411.4028, 2014.
\item A. Lucas, « Ising formulations of many NP problems », \emph{Frontiers in Physics} 2:5, 2014.
\item T. Albash, D. A. Lidar, « Adiabatic quantum computation », \emph{Reviews of Modern Physics} 90, 015002, 2018.
\item L. Zhou et al., « Quantum Approximate Optimization Algorithm: Performance, Mechanism, and Implementation on Near-Term Devices », \emph{Physical Review X} 10, 021067, 2020 (stratégie INTERP).
\item H. Ni et al., « Airfoil shape optimization via coherent Ising machine », arXiv:2603.10518, 2026 (référence [3] de la note d'intention).
\end{itemize}

\chapter{Pistes d'amélioration, chapitre par chapitre}\label{annexe:pistes}
\iterbanner{V1}

Cette annexe rassemble, pour chaque partie du cours, ce qu'il serait utile d'approfondir ou d'ajouter au code. Les pistes sont classées de la plus utile à la moins utile pour le hackathon. Le cours compagnon \emph{Optimisation bi-objectif traînée / SER} met en œuvre une partie d'entre elles (QUBO bi-objectif avec contraintes, chaîne complète).

\section*{Bases, qubit, algèbre, plusieurs qubits (ch.~\ref{ch:intro} à \ref{ch:multi})}
\begin{itemize}
\item Ajouter un petit carnet NumPy qui reproduit chaque exemple chiffré (sphère de Bloch, mesures dans trois bases, état de Bell) : vérifier soi-même les calculs est le meilleur exercice.
\item Introduire la matrice densité, pour décrire le bruit et les mélanges statistiques : indispensable pour comprendre ce que donne une vraie machine.
\end{itemize}

\section*{Portes, circuits, hamiltoniens (ch.~\ref{ch:circuits} et \ref{ch:hamiltoniens})}
\begin{itemize}
\item Compiler le circuit QAOA de notre QUBO sur une topologie réelle (grille ou « heavy-hex » d'IBM) et compter les SWAP : c'est le vrai coût en portes à deux qubits.
\item Mesurer l'erreur de Trotter sur l'exemple jouet, en comparant l'évolution exacte et l'évolution découpée.
\end{itemize}

\section*{QUBO et Ising (ch.~\ref{ch:qubo})}
\begin{itemize}
\item \iterbadge{V1}~\textbf{Fait dans la chaîne complète} : contraintes par réduction de Rosenberg (qubit auxiliaire) et par variables d'écart. Il reste à étudier l'effet du poids des pénalités sur la qualité des solutions de QAOA.
\item Comparer les encodages binaire, one-hot et « unaire » (domain wall) sur la même grille, en qubits, en couplages et en qualité de solution.
\end{itemize}

\section*{Recuit (ch.~\ref{ch:recuit})}
\begin{itemize}
\item Valider notre recuit simulé contre celui d'Ocean : le paquet \texttt{dwave-neal} s'installe dans l'environnement du projet, \texttt{Code/.venv}.
\item Étudier le plongement de notre graphe complet à 13 qubits sur la topologie Pegasus de D-Wave (\texttt{minorminer}) : longueur des chaînes, nombre de qubits physiques.
\end{itemize}

\section*{QAOA (ch.~\ref{ch:qaoa})}
\begin{itemize}
\item Réutiliser les angles optimaux d'une itération de la boucle à la suivante (le QUBO change peu) : on économiserait la plupart des évaluations de COBYLA.
\item Tester les variantes « warm-start » (état initial tiré d'une solution classique) et le mélangeur XY qui conserve le nombre de 1, adapté aux encodages one-hot.
\item Simuler l'effet du bruit (dépolarisation par porte à deux qubits) sur la probabilité de mesurer les bons profils.
\end{itemize}

\section*{Outils (ch.~\ref{ch:outils})}
\begin{itemize}
\item Installer Qiskit et Ocean dans \texttt{Code/.venv} et transformer les esquisses du chapitre en code testé, comparé à notre implémentation NumPy.
\end{itemize}

\section*{Chaîne du projet (ch.~\ref{ch:projet})}
\begin{itemize}
\item \iterbadge{V1}~\textbf{Fait} : la boucle bi-objectif traînée/SER, comparée à NSGA-II (cours d'optimisation). Résultat : avec le solveur exact, la boucle retrouve le front avec 1,6 à 3 fois moins de simulations que NSGA-II avec ses réglages par défaut (1,2 à 1,9 fois face à NSGA-II réglé).
\item Comparer la boucle à une optimisation bayésienne multi-objectif classique (même modèle de substitution, choix des candidats par amélioration attendue de l'hypervolume), pour isoler ce qu'apporte la résolution du QUBO.
\end{itemize}


\end{document}
